\chapter{C2: strong and single-angle sources}
\label{ch:cert-6}
Formula tags are local to their certificate section. All source and testing fields remain separately labelled.

\section{\texorpdfstring{A complete local strong-II Gram with independent testing rates}{A complete local strong-II Gram with independent testing rates}}
\label{component:pe16499ccc4c6}

\label{pe16499ccc4c6:sec:r7-local-strong-II}

\begin{theorem}[Strong-II local allocation]
Let the positive C2 source satisfy $z,y\in\Inf_{>0}$ and let
$B=H+L$ be a positive unit.  Keep the exact source coefficients of
(GE1), with $n\geq0$ and the remaining displayed weights positive.
Then the original C2 parent belongs to $\sum(K_c\otimes_{\mathbb R}
K_t)^2$ whenever both independent affine testing coordinates $u,v$
are infinitesimal, including either literal testing axis.
There is no comparison hypothesis between $u^2,z$, between $v^2,y$,
or between the two source rates.  All openings of $r$, $L/y$,
$1+zr-L/y$, and $1-H$ are retained, as is literal $n=0$.
\end{theorem}

This theorem is a full-parent local result.  The remote testing charts
and the original projective endpoints are separate obligations; they
must be joined before the strong II/HH checklist row is admitted.
All identities below are polynomial in testing coordinates.  Inverses
are taken only in the pure source field, and the final step uses the
bounded Gram strict-real-residue lemma already used in the completed
local UI certificate.  In particular no ordering or inverse of the
mixed tensor is used.

\paragraph{Source units and notation.}
Write $U=u^2,V=v^2,k=\sqrt y$ and
\[
 a_c=1+zr,\quad t=L/y,\quad j=a_c-t,\quad s=a_c/B,
 \quad c=CL,\quad d=Dyj.
\]
The source entry gives bounded $t,j,r$, positive units $B,c,d,m/z^2,A/y$,
bounded $nr$ and $\alpha(1-H)$, and strong $\alpha$.
These facts include the arbitrary openings $t\to0$, $j\to0$,
$r\to0$, and $1-H\to0$.  They follow directly from the positive
source formulas and were proved in (DI7).  Let $h=\operatorname{st}B$;
then $0<h\leq1$ and, with $r_0=\operatorname{st}r$,
\[
\begin{gathered}
 c_0=d_0^{\rm pair}=\frac{h}{2(1+h)},\quad
 \operatorname{st}(m/z^2)=\frac{1-r_0}{2h^2(1+h)},\quad
 \operatorname{st}(A/y)=\frac1{2h(1+h)},\\
 \operatorname{st}(nr)=\frac1{2(1+h)}-\frac{r_0}{2},\qquad
 \operatorname{st}\{\alpha(1-H)\}=\frac{h^2}{2(1+h)}.
\end{gathered}
\tag{LII1}
\]
These residues follow directly from (GE1): $L\to0$, $H\to h$,
$a_c\to1$, $G\to h$, $\delta\to h(1+h)$, and the source
numerator $M\to h^2$.  Substitution in the six rational coefficient
formulas gives (LII1), including the identity
$\operatorname{st}(nr)=1/[2(1+h)]-r_0/2$ without division by $n$.
Define the source quantities (with $q,g$ bounded)
\[
 \ell=y(C+D)=c/t+d/j,\qquad
 q=\frac{D}{C+D}=\frac{dt}{cj+dt},\qquad
 g=\frac{yCD}{C+D}=\frac{cd}{cj+dt}.
\]
Here $0<q<1$, $g$ is a positive unit with residue $h/[2(1+h)]$,
and $\ell$ has a positive real lower bound, although it may be
unbounded.  Indeed $cj+dt$ is a unit because $t+j=a_c$ is a unit.
Choose positive real $\ell_*\leq\ell$ and
$\alpha_*\leq\alpha$.

\paragraph{The state roots and the exact C,D completion.}
Use the source-fixed parity graph
\[
 X=-\frac{zr v}{B}W+\xi,\qquad
 Y=\frac{Hu}{B}W+\eta,\qquad
 Z=\frac{uv}{B}W+\zeta.                         \tag{LII2}
\]
It agrees with the actual kernel at $U=z,V=y$ and is a change of all
four homogeneous state coordinates with determinant one.  Set
\[
 Q_n=u\xi-\zeta,\quad K=ukW,\quad
 R_h=\frac{vQ_n-su(V-y)W}{k},\qquad
 N=\frac\eta k-qR_h.                              \tag{LII3}
\]
Thus $\eta=k(N+qR_h)$.  Every root in (LII3) vanishes at the actual
contact.  If $Q=uW-Y$ and $J=Q+v(uX-Z)$, then
\[
 Q/k=(t/B)K-N-qR_h,\qquad
 J/k=(1-q)R_h-N-(j/B)K.
\]
Consequently the exact weighted completion is
\[
\begin{split}
 4CQ^2+4DJ^2={}&4\ell N^2+4gR_h^2
 +\frac{8(d-c)}BKN-8gsKR_h
 +\frac{4(ct+dj)}{B^2}K^2,\\
 4CVQ^2={}&4CV\eta^2
 +\frac{4ct}{B^2}VK^2-\frac{8c}BVK(N+qR_h).
                                                        \end{split}\tag{LII4}
\]
There is no estimate in these identities.  Retain the following
four nonnegative, contact-centered packets outside the remainder:
\[
 4(\ell-\ell_*)N^2,\qquad 4CV\eta^2,\qquad
 4n(1+U)(1+V)\xi^2,\qquad
 4(\alpha-\alpha_*)(1+U)(v\eta-\zeta)^2.          \tag{LII5}
\]
Call the exact remainder $G$.  Its coefficients are bounded after
writing $\eta=k(N+qR_h)$ and using (LII4).  For completeness, its
entire formula is
\[
\begin{split}
G={}&4\ell_*N^2+4gR_h^2+\frac{8(d-c)}BKN-8gsKR_h\\
 &+\frac{4(ct+dj)}{B^2}K^2
 +\frac{4ct}{B^2}VK^2-\frac{8c}BVK(N+qR_h)\\
 &+4n(1+U)(1+V)
   \left\{\frac{z^2r^2V}{B^2}W^2-\frac{2zr v}B W\xi\right\}\\
 &+4\alpha_*(1+U)(v\eta-\zeta)^2
  +4\alpha(1+U)
   \left\{\frac{(1-H)^2UV}{B^2}W^2
        -\frac{2(1-H)uv}B W(v\eta-\zeta)\right\}\\
 &+4mM^2+4A(1+U)(1+V)Y^2\\
 &+2\{J^2+(1+U)(vY-Z)^2+(1+V)(X+uZ)^2
                      -(1+U)(1+V)Z^2\},        \end{split}\tag{LII6}
\]
where $M=vW-X+uvY-uZ$ and (LII2) is used in the last line.
Every potentially unbounded original coefficient occurs here only
through $nr$, $\alpha(1-H)$, or the bounded quantities in (LII1).
No division by a testing coordinate has been introduced.

\paragraph{One exact gauge removes the dependent-root first normal.}
In the formal bounded polynomial (LII6), the coefficient of $R_hW$
is $uk\rho(U,V)$ for a bounded source polynomial $\rho$.  Replace
this entire term using the exact identity
\[
 KR_h=uvWQ_n-sU(V-y)W^2.                         \tag{LII7}
\]
Call the resulting polynomial $\widetilde G$.  It is equal to $G$
on the roots (LII3), and has no $R_hW$ term.  This gauge is useful
because $R_h$ is a dependent root: it permits the state and angular
contact equations to be read as formal coefficient vanishings.

More precisely, by parity write uniquely
\[
 \widetilde G=Q_0(N,\xi,\zeta,R_h)
 +W\{vP_\xi(U,V)\xi+uvP_\zeta(U,V)\zeta
                         +ukP_N(U,V)N\}
 +\phi(U,V)W^2,                                  \tag{LII8}
\]
where all coefficients are bounded source polynomials and $Q_0$ is
quadratic in the four displayed roots.  Formula (LII6), followed by
the one replacement (LII7), is an exact coefficient-extraction
specification for every entry in (LII8).

The retained packets (LII5) vanish together with their first jets at
the actual contact.  The exact original source state equations thus
give
\[
 P_\xi(z,y)=P_\zeta(z,y)=P_N(z,y)=0.              \tag{LII9}
\]
Indeed the physical $\eta$ derivative first forces the $N$ coefficient
to vanish; the physical $\xi,\zeta$ derivatives then force their
formal coefficients to vanish because the $R_hW$ coefficient was
removed.  The exact original contact and the two angular derivatives
now give
\[
 \phi(z,y)=\phi_U(z,y)=\phi_V(z,y)=0.              \tag{LII10}
\]
This uses the six exact original source substitutions already checked
in R6.  It imposes no new relation between source and testing fields.

\paragraph{Bounded division and the complete angular allocation.}
There are two additional elementary support facts in (LII8):
\[
 P_\xi\in(U,z),\qquad \deg_U\phi\leq2,\qquad
 \phi(0,V)=V(d_s+e_sV),                           \tag{LII11}
\]
where
\[
 d_s=\frac4{B^2}\{m(B+zr)^2+(n+\tfrac12)z^2r^2\},\qquad
 e_s=\frac4{B^2}(n+\tfrac12)z^2r^2.
\]
Both $d_s/z^2$ and $e_s/z^2$ are bounded.  The first support assertion
is literal bounded-coefficient membership: setting $U=z=0$ in the
extracted polynomial gives zero.  It is also clear before expansion:
all $\xi W$ terms not supplied by the gauge contain $U$, $m$, $nzr$,
or $zr$; those supplied by the gauge contain $U$.

Monic division using (LII9) therefore gives bounded polynomials
\[
\begin{split}
 P_\xi&=(U-z)A_\xi+z(V-y)B_\xi,\\
 P_\zeta&=(U-z)A_\zeta+(V-y)B_\zeta,\\
 P_N&=(U-z)A_N+(V-y)B_N.                         \end{split}\tag{LII12}
\]
For the first line, evaluate $P_\xi$ at $U=z$: the support assertion
supplies the factor $z$ with bounded quotient, and its value at $V=y$
is zero.  All other divisions are by the monic polynomials $U-z$
or $V-y$.  In particular none creates an inverse of a small source
or a testing difference.

The scalar has the exact bounded decomposition
\[
 \phi=\frac{d_s+e_sV}{z^2}V(U-z)^2
       +U(U-z)(V-y)A_s(V)+U(V-y)^2C_s(V).        \tag{LII13}
\]
Here is a constructive proof, including boundedness.  Write
$\phi=a(V)U^2+b(V)U+V(d_s+e_sV)$ and subtract the first term of
(LII13).  The remaining $U^2,U$ coefficients $a_1,b_1$ vanish at
$V=y$ by the value and $U$ derivative in (LII10).  Divide them by
$V-y$ to obtain $A_s,B_s$.  The $V$ derivative gives
$zA_s(y)+B_s(y)=0$, so set $C_s=(B_s+zA_s)/(V-y)$.
All three divisions are monic divisions of bounded polynomials.
This proves (LII13) exactly, not asymptotically.

Introduce the two legal angular roots
\[
 R=v(U-z)W,\qquad T=uk(U-z)W.                    \tag{LII14}
\]
The following identities allocate every term in (LII12)--(LII13):
\[
\begin{split}
 zv(V-y)W&=\frac{uVQ_n-uvkR_h}{s}-(V-y)R,\\
 uv(V-y)W&=\frac{VQ_n-vkR_h}{s},\\
 uk(V-y)W&=\frac{vkQ_n-yR_h}{s},\\
 U(U-z)(V-y)W^2&=-\frac{R_hT-uRQ_n}{s},\\
 U(V-y)^2W^2&=\frac{VQ_n^2-2vkR_hQ_n+yR_h^2}{s^2}.
                                                        \end{split}\tag{LII15}
\]
Together with $v(U-z)W=R$, $uv(U-z)W=uR$, and
$uk(U-z)W=T$, these are an explicit complete allocation of the
full remainder as a bounded symmetric Gram in
\[
                    (N,\xi,\zeta,R_h,R,T).       \tag{LII16}
\]
The original first state normals are all present.  The scalar
restriction has not been used as a substitute for them.

\paragraph{The finite block and the missing pivot.}
The residue of $Q_0$ on the first four roots is
\[
             \operatorname{diag}(4\ell_*,2,4\alpha_*,4g_0).
                                                        \tag{LII17}
\]
For example the signed core at $W=0$ is exactly
$2(\xi+u\zeta-uv\eta)^2+2(1+U)V\xi^2+2(1+V)\eta^2$;
all its other entries vanish in the present residue.  The extra
terms in the first three lines of (LII15) have infinitesimal Gram
coefficients except for $A_\xi\xi R$ and the allowed $A_NNT$.
Likewise the last line has only infinitesimal coefficients, and
the fourth line has one potentially finite $R_hT$ entry.  Thus,
apart from entries in the last row and column $T$, the sole finite
off-diagonal entry is between $\xi$ and $R$.

The exact coefficient extraction yields
\[
 \operatorname{st}A_\xi=\frac4{h(1+h)},\qquad
 \operatorname{st}(d_s/z^2)=\frac2{h^2(1+h)}.     \tag{LII18}
\]
For a short independent check of the first value, the unscaled
original first-normal coefficient on the $V=y$ slice is
\[
 -4nzr(1+y)+4m(1-H)-(4D+2)yj+2(1+y)
 \longrightarrow 2-4d_0^{\rm pair}=\frac2{1+h}.
\]
Its full multiplier is $2/B$.  Passing from the original normal
to (LII8) changes this $U$-quotient only by a bounded multiple of
$U$ or $z$, so its residue is unchanged.  The second value follows
from (LII1):
\[
 \operatorname{st}\{m(B+zr)^2/z^2+(n+\tfrac12)r^2\}
 =\frac{1-r_0}{2(1+h)}+\frac{r_0}{2(1+h)}
 =\frac1{2(1+h)}.
\]
Hence the finite $(\xi,R)$ Gram block is exactly
\[
 \begin{pmatrix}
  2&\dfrac2{h(1+h)}\\[2mm]
  \dfrac2{h(1+h)}&\dfrac2{h^2(1+h)}
 \end{pmatrix},\qquad
 \text{Schur complement}=\frac2{h(1+h)^2}>0.      \tag{LII19}
\]
This calculation retains $r_0=0$, $n=0$, and $h=1$.

Finally the exact toric conductor is the zero identity
\[
 \boxed{\quad
 T^2-UR^2-\frac{U-z}{s}R_hT
             +\frac{u(U-z)}s RQ_n=0.\quad}       \tag{LII20}
\]
Add any positive real multiple $\Lambda$ of this zero identity to
the Gram.  Its only nonzero residue contribution is $\Lambda T^2$.
The five-by-five block excluding $T$ is strictly positive by
(LII17)--(LII19); all entries incident to $T$ are bounded.  Choose
a sufficiently large real $\Lambda$ by the ordinary real Schur
complement.  The entire six-by-six residue Gram is then strictly
positive.  The bounded Gram lemma gives a tensor SOS for $G$, and
adding the four packets (LII5) proves the theorem.

\paragraph{Quantifiers and endpoint discipline.}
No proof step divides by $u$, $v$, $U-z$, or $V-y$: parity quotients
and monic polynomial divisions are coefficient extractions whose
final identities are polynomial.  Thus the two literal affine axes
are included.  The source inverses $1/B$, $1/k$, and the source
packet square roots are legal even when their values are large;
all coefficients of the final Gram are bounded.  The source ratios
$t,j,r,1-H$ have not been bounded away from zero.  The new roots
raise testing degree, so the local certificate is not homogenized
to assert the original literal projective endpoints.  Those are
joined by the separate direct-parent remote certificate.

\paragraph{Source input.}
Source contact and stationarity follow from the native-entry identification in Chapter~\ref{ch:entry} and the exact source formulas in Section~\ref{component:p6e5d41ed80b3}. The scaled bounds used above are the II specialization of GE1--GE2 in Section~\ref{component:pf5726218aec8}, recorded there as DI7.


\section{\texorpdfstring{The entire remote complement of strong C2 II}{The entire remote complement of strong C2 II}}
\label{component:p590112f102bd}

\label{p590112f102bd:sec:r7-c2-remote}

This section is a full-parent tensor-SOS admission.  It applies to the
original C2 positive chart with $z,y$ infinitesimal and $B=H+L$ a
positive unit, and to every testing pair for which at least one testing
angle is not infinitesimal.  Literal testing axes and projective
infinities are included in the assertions below.  The local case in
which both testing angles are infinitesimal is not asserted here.

\paragraph{Source constants, including all openings.}
Put $t=L/y$, $j=1+zr-t$, and write $h=\operatorname{res}H$.
Then $0<h\leq1$, $t+j=1+zr\sim1$, and the exact source formulas give
\[
 Cyt\sim Dyj\sim\frac{h}{2(1+h)},\qquad
 \frac Ay\sim\frac1{2h(1+h)}.
\]
Consequently $C,D$ are unbounded, and the positive pure-source quantity
\[
 \gamma=\frac{yCD}{C+D}
 =\frac{(Cyt)(Dyj)}{(Cyt)j+(Dyj)t}
 \sim\frac{h}{2(1+h)}.                                      \tag{RR1}
\]
These statements are uniform when either $t$ or $j$ tends to zero.
If $h<1$, then
\[
 \alpha\sim\frac{h^2}{2(1-h^2)};
\]
if $h=1$, then $\alpha$ is unbounded.  Thus, for every real
$0<h_0<h$, the positive real number
\[
 \beta_0=\frac{h_0^2}{2(1-h_0^2)}<\alpha                     \tag{RR2}
\]
is an admissible cap.  No lower bound on $r,j,t,1-H$ is imposed,
and $n=0$ is retained.

\paragraph{Homogeneous testing coordinates.}
For finite $u,v$ define
\[
 a=\frac{u}{\sqrt{1+u^2}},\quad b=\frac1{\sqrt{1+u^2}},\qquad
 c=\frac{v}{\sqrt{1+v^2}},\quad d=\frac1{\sqrt{1+v^2}}.
\]
Use the same homogeneous formulas at infinity, with $(a,b)=(1,0)$
or $(c,d)=(1,0)$.  These are the original parent's homogeneous
endpoint values.  In particular no raised-degree graph is being
specialized at an endpoint.  Put
\[
 E=aW-bY,\quad K=aX-bZ,\quad L=bX+aZ,
 \quad J=dE+cK,\quad Q=cY-dZ.
\]
The original parent divided by $4(1+u^2)(1+v^2)$ is exactly
\[
 nX^2+m\{bcW-bdX+acY-adZ\}^2+CE^2+AY^2
 +(D+\tfrac12)J^2+(\alpha+\tfrac12)Q^2
 +\tfrac12(L^2-Z^2).                                    \tag{RR3}
\]
At projective endpoints this equation means the original homogeneous
form, not a division by an infinite number.  All radicals and
normalizations used here belong to a single testing field.

\begin{proposition}[All tests with $v$ not infinitesimal]
For every such $v$, including $v=\infty$, and every $u$, including
$u=0,\infty$, the original full parent is tensor-SOS.
\end{proposition}
\begin{proof}
Here $c$ is a unit.  Set $P=bL$, $R_1=\sqrt C E$, and
$R_2=\sqrt D J$.  Orthogonality gives
\[
 L^2-Z^2=P^2+2aKP-b^2K^2,
 \qquad
 K=\frac{R_2}{c\sqrt D}-\frac{dR_1}{c\sqrt C}.            \tag{RR4}
\]
Both coefficients on the right are infinitesimal bounded tensors:
$1/c$ is bounded in the testing field, while $1/\sqrt C$ and
$1/\sqrt D$ are infinitesimal in the source field.  Therefore
\[
 R_1^2+R_2^2+\tfrac12P^2+aKP-\tfrac12b^2K^2
\]
has a bounded Gram with residue $\operatorname{diag}(1,1,1/2)$.
It is tensor-SOS by the strict bounded-Gram lemma.  Every other term
of (RR3), namely the retained $n,m,A,\alpha$ packets and
$J^2/2+Q^2/2$, is already SOS.  This uses all four state coordinates.
\end{proof}

\paragraph{An exact harmonic identity for the remaining remote strip.}
Now let $v$ be infinitesimal, possibly zero, and let $u$ not be
infinitesimal, possibly infinity.  Thus $a$ is a unit.
Write $D_v=1+v^2$, and put
\[
 s=C+D,\quad\theta=D/s,\quad
 N=\sqrt{s}(E+\theta vK),\quad
 V=\frac{vK}{\sqrt y},\quad q=\sqrt yY.
\]
No comparison between $v$ and $\sqrt y$ is made.  The roots may be
unbounded; their Gram coefficients below are bounded.  The exact
pure-source harmonic identity is
\[
 CD_vE^2+D(E+vK)^2=N^2+\gamma V^2+Cv^2E^2.               \tag{RR5}
\]
Also set $x=bX$, $z_0=Z$, and $\ell=vY$.  The single toric relation
needed below is the polynomial identity
\[
 a x\ell-b^2z_0\ell=bVq.                                \tag{RR6}
\]
For a real cap $0<\beta_0<\alpha$, the original parent divided
by $4(1+u^2)$ is the sum of the explicitly positive packets
\[
 \begin{split}
 &nD_vX^2+m\{b(vW-X)+a(vY-Z)\}^2+N^2+Cv^2E^2\\
 &\hspace{15mm}+\tfrac12(E+vK)^2+(\alpha-\beta_0)(\ell-z_0)^2
 \end{split}                                             \tag{RR7}
\]
and the following exact bounded quadratic form:
\[
 G=\gamma V^2+\frac AyD_vq^2
 +(\beta_0+\tfrac12)(\ell-z_0)^2
 +\frac{D_v}{2}(x^2+2axz_0-b^2z_0^2).                   \tag{RR8}
\]
This decomposition is valid also at $u=\infty$, by the original
homogeneous convention.  It has used no state graph or state inverse.

\begin{proposition}[All tests with $u$ not infinitesimal and $v$ infinitesimal]
The full original parent is tensor-SOS for all these tests, including
$v=0$ and $u=\infty$.
\end{proposition}
\begin{proof}
For a real constant $k$ define
\[
 M_k=(\beta_0+\tfrac12)(\ell-z_0)^2
 +\tfrac12(x^2+2axz_0-b^2z_0^2)
 +k(ax\ell-b^2z_0\ell).
\]
By (RR6), the exact form (RR8) equals
\[
 \boxed{G=D_vM_k+
 \gamma V^2+\frac AyD_vq^2-kD_vbVq
 -v^2(\beta_0+\tfrac12)(\ell-z_0)^2.}                   \tag{RR9}
\]
This is a full-parent identity; every state-normal term is present.

First suppose $b$ has positive residue $b_0$.  Since $a$ is a unit,
$0<b_0<1$.  Choose $0<h_0<h$ sufficiently close to $h$ that
\[
 b_0(1+h-h_0)<1,
\]
and take the admissible cap (RR2).  There exists a positive real
$\kappa$ satisfying
\[
 \frac1{1+b_0h_0}<\kappa<
 \min\left\{\frac1{1-b_0h_0},\frac1{b_0(1+h)}\right\}.
                                                               \tag{RR10}
\]
Set $k=-\kappa$.  The Gram of $M_k$ on $(x,z_0,\ell)$ has
first pivot $1/2$, second Schur pivot $\beta_0$, and determinant
\[
 \det M_k=\frac{2\beta_0 b^2k^2-(2\beta_0+1)(k+1)^2}{8}.
                                                               \tag{RR11}
\]
Its residue is strictly positive by the first two inequalities in
(RR10), because $h_0^2=2\beta_0/(2\beta_0+1)$.  The Gram on $(V,q)$
in (RR9) has positive diagonal residues and determinant residue
\[
 \frac14\left\{\frac1{(1+h)^2}-k^2b_0^2\right\}>0.       \tag{RR12}
\]
Thus both blocks have strictly positive residue.  The final term in
(RR9) is a bounded infinitesimal perturbation on $(x,z_0,\ell)$,
so the combined bounded Gram remains strict and is tensor-SOS.

If $b$ is infinitesimal, including $b=0$, choose any positive real
cap $\beta_0<\alpha$ and take $k=-1$.  There is now an exact
three-square identity
\[
 M_{-1}=\tfrac12\{x+a(z_0-\ell)\}^2
       +\beta_0(z_0-\ell)^2+\tfrac12(b\ell)^2.           \tag{RR13}
\]
On these three roots and $(V,q)$, (RR9) has positive diagonal
residue; the only off-block term is $D_vbVq$, whose coefficient
is infinitesimal.  The negative perturbation changes only the
$(z_0-\ell)^2$ coefficient by an infinitesimal.  Its Gram is again
strict.  Formula (RR7) restores all retained positive packets.
The argument does not divide by $u,v,b$, or any state coordinate,
so the asserted literal boundaries are included.
\end{proof}

\paragraph{Scope and join.}
The two propositions cover precisely the complement of the region
where both testing angles are infinitesimal, for every strong-II
source with $B$ a unit.  They preserve independent source and testing
fields, arbitrary source openings, testing signs, and literal $n=0$.
Combined with a separately verified full-parent admission of that
local region, they close the strong II/HH orbit by the exact C2
reflection.  They do not close the $B$-infinitesimal weak face.


\section{\texorpdfstring{Strong HI with an arbitrary $n$ event}{Strong HI with an arbitrary n event}}
\label{component:p4363feef60dd}

\label{p4363feef60dd:sec:r7-hi-full-retained}
The preceding contact lift extends to every original $n\geq0$; no
additional comparison of source and testing scales is required. Retain
all notation in (N0M.2)--(N0M.5), except for the coefficient replacements
specified below. To distinguish the original source data from the four
toric roots, write $A_{\rm src}$ for the original coefficient $A$ and
$B_{\rm src}=H+yt$ for the original source $B$. Put
\[
 \nu=2nR=\frac{(B_{\rm src}+R)(1+R)S_c-R\Delta}{\Delta},
 \qquad h=H-\nu,\qquad \zeta=R+\nu.
\]
The source quantity $\nu$ is nonnegative and bounded, even when $R$ is
infinitesimal and $n$ is unbounded. Exact source identities are
\[
 \frac{2A_{\rm src}H}{qy}=\zeta,\qquad
 c-d=\frac y2\{\zeta(1+q)-t\},\qquad
 \mu-a_*=\frac q2\{E-\zeta(1+y)\}.             \tag{HI7.1}
\]

In addition to all four opening-excess squares $\mathcal E$ in
(N0M.3), retain the actual nonnegative square packet
\[
 \mathcal N=4nd_pd_vx^2,\qquad
 x=X+\frac RHpvY
   =-g s+(H+R)TQ+v\sqrt y\,N_c-p\sqrt q\,N_m.
                                                               \tag{HI7.2}
\]
This packet is zero at the original source contact. Its exact difference
from the original $4nd_pd_vX^2$ is
\begin{align*}
 4nd_pd_v(X^2-x^2)={}&4\nu d_pd_vg\,TsQ
 -2\nu(2H+R)d_pd_vT^2Q^2\\
 &-4\nu d_pd_vTQ(v\sqrt y\,N_c-p\sqrt q\,N_m).
\end{align*}
Every coefficient on the right is bounded. The unbounded coefficient
$n$ occurs only in the retained positive square.

The full exact identity (N0M.4) therefore holds with
$F-\mathcal E$ replaced by $F-\mathcal E-\mathcal N$, with the same
$h_c,h_m,f_A,f_B$, and with
\begin{align}
 C_s&=4d_v\left[y\{\zeta(1+q)-d_pg\}
                       +\frac{2dt}{S_c}\delta_v\right],\notag\\
 M_s&=4d_p\left[q\{d_vg-\zeta(1+y)\}
                       -\frac{2\mu E}{S_c}\delta_p\right],\notag\\
 \beta_c&=4h\sqrt y\,v d_pd_v,
 &\beta_m&=-4h\sqrt q\,p d_pd_v,\notag\\
 f_C&=2\{H^2-\nu(2H+R)\}d_pd_v,
 &f_D&=2H(R+\nu)d_pd_v,\notag\\
 e&=4h d_pd_vg.                                  \tag{HI7.3}
\end{align}
In this notation the complete first-normal couplings are
$N_c(C_sB+\beta_cC)+N_m(M_sA+\beta_mC)$, and the graph is
$\sum_i f_iA_i^2-eG^2$ on the positive sign chart. Every coefficient
is a polynomial in $p,v$ with bounded source coefficients. The normal
coupling coefficients are infinitesimal. Their exact contact sums vanish.

Let $r=\bar R\in[0,1]$. The original source formulas give
\[
 \bar\nu=\frac{1-r}{3+r},\qquad
 \bar c=\bar d=\bar\mu=\bar a_*
        =\frac{1+r}{2(3+r)}.
\]
Consequently all four graph coefficients have the same strictly positive
residue
\[
 \bar f_A=\bar f_B=\bar f_C=\bar f_D
       =\lambda(r)=\frac{2(1+r)^2}{3+r}\in[2/3,2],
 \qquad \bar e=4\lambda(r).                       \tag{HI7.4}
\]
The two state-normal diagonal coefficients have residue
$4/(3+r)\in[1,4/3]$.
These statements remain valid at $r=0$, where $n$ can be unbounded.

All exact contact jets (N0M.6) hold unchanged for the new coefficients;
also $C_s^*+\beta_c^*=M_s^*+\beta_m^*=0$. They can either be obtained
by subtracting the contact-vanishing packets from the original parent,
or checked directly using (HI7.1). These are rational identities for the complete source, without imposing
$n=0$ or replacing $R$ by $r$.

Choose $k=(f_C^*+f_D^*)/2$, now a positive bounded source unit with
residue $\lambda(r)$. The exact half-angle contact lift (N0M.8)--(N0M.9)
then applies verbatim. The coefficient errors are infinitesimal bounded
tensors, the finite circuit supplies strictly positive pivots on every
conductor root, and the state-normal block has a strictly positive
scalar residue. On
the opposite-sign or zero-state charts the direct strict raw Gram
argument applies as before. The strict bounded-Gram lemma supplies an
exact tensor SOS for $F-\mathcal E-\mathcal N$; adding the five retained
positive packets gives an exact SOS for the unchanged full original
parent throughout the critical HI testing corner.

No inversion or order on a mixed tensor is used. The only radicals of
state expressions are taken in the pure-testing real closed field after
a pure-testing sign choice. The result is a tensor SOS at each testing
state, not a claim of an SOS with the original polynomial degree in the
state amplitudes. Together with the already established remote tests
and literal axes, this supplies the remaining strong C2 HI admission.



\section{\texorpdfstring{Strong HI: the complete contact module}{Strong HI: the complete contact module}}
\label{component:p888d9f79ba3b}

\label{p888d9f79ba3b:sec:r7-hi-n0-retained}
All source and testing variables in this section retain their original
meanings. In particular $q,y,p,v>0$ are independently infinitesimal;
no comparison between a source scale and a testing scale is made.
Write
\[
 d_p=1+p^2,\quad d_v=1+v^2,\quad
 \delta_p=p^2-q,\quad\delta_v=v^2-y,\quad
 S_c=1+R+yt=j+t(1+y)=K+E(1+q).
\]
Use bounded source units
\[
 c=Cyt,\qquad d=Dyj,\qquad \mu=mqK,\qquad a_*=\alpha qE.
\]
At literal $n=0$, the exact source formulas give
\begin{align}
 2AH/(qy)&=R, & \bar R&=1,\\
 c-d&=\frac y2\{R(1+q)-t\},&
 \mu-a_*&=\frac q2\{E-R(1+y)\}.                 \tag{N0M.1}
\end{align}
The four starred weights have residue $1/4$, even when any positive
opening tends to zero. The last two identities follow from the rational
source formulas; before imposing $n=0$ their respective right sides are
$y\{(2n+1)R(1+q)-t\}/2$ and
$q\{E-(2n+1)R(1+y)\}/2$.

\paragraph{A common state with exact contact.}
Set
\[
 A_0=W-pY,\qquad B_0=vY-Z,\qquad
 S=X-vW-pZ+pvY,\qquad s=-S/S_c,
 \qquad Q=\frac{\sqrt{qy}}H Y.
\]
Then $J=d_vA_0+vS$ and $M=d_pB_0-pS$. Define
\[
 N_c=\frac{A_0-tv s}{\sqrt y},\qquad
 N_m=\frac{B_0+Ep s}{\sqrt q},\qquad
 P=\frac p{\sqrt q},\quad V=\frac v{\sqrt y},\quad T=PV,
\]
and the additional contact normals
\[
 N_D=d_vN_c+t\delta_vVs,\qquad
 N_M=d_pN_m-E\delta_pPs.                         \tag{N0M.2}
\]
At the actual source contact, $P=V=1$, $s=Q=\sqrt q$, and
$N_c=N_m=N_D=N_M=0$. All definitions are tensor identities; neither
$P$ nor $V$ is assumed bounded.

\paragraph{Retain all singular opening excesses.}
The following four positive squares are retained in the parent:
\begin{align}
 \mathcal E={}&\frac{4cj d_v}{tS_c}N_c^2
 +\frac{4dt(1+y)}{jS_c}N_D^2\notag\\
 &+\frac{4\mu E(1+q)}{KS_c}N_M^2
 +\frac{4a_*K d_p}{ES_c}N_m^2.                 \tag{N0M.3}
\end{align}
Their coefficients are products of positive pure-source and pure-testing
factors, so the displayed terms are actual tensor squares. They need
not be bounded. No excess is discarded before collecting the remainder.

Put $g=S_c-tv^2-Ep^2$. The unchanged full parent has the exact identity
\begin{align}
 F-\mathcal E={}&h_cN_c^2+h_mN_m^2
 -4\sqrt{qy}\,pv d_pd_vN_cN_m\notag\\
 &+N_c\{C_sVs+4H\sqrt y\,v d_pd_vTQ\}
 +N_m\{M_sPs-4H\sqrt q\,p d_pd_vTQ\}
 +\mathcal G,                                           \tag{N0M.4}
\end{align}
where
\begin{align*}
 h_c&=\frac{4d_v}{S_c}\{c(1+y)+dd_v\}
       +2yd_v(1+v^2+p^2v^2),\\
 h_m&=\frac{4d_p}{S_c}\{\mu d_p+a_*(1+q)\}
       +2qd_p(1+p^2+p^2v^2),\\
 C_s&=4d_v\left[y\{R(1+q)-d_pg\}
                    +\frac{2dt}{S_c}\delta_v\right],\\
 M_s&=4d_p\left[q\{d_vg-R(1+y)\}
                    -\frac{2\mu E}{S_c}\delta_p\right].
\end{align*}
The two diagonal coefficients have residue one. The normal--normal
off-diagonal coefficient and all four raw coefficients in the first
normal line are infinitesimal bounded tensors. In particular no inverse
opening occurs in (N0M.4).

The scalar graph is
\[
 \mathcal G=f_AP^2s^2+f_BV^2s^2
       +f_CT^2Q^2+f_DQ^2-eTsQ,                         \tag{N0M.5}
\]
with
\begin{align*}
 f_A={}&4\left[\mu K+a_*Ed_p-2\mu E\delta_p
                    +\frac{\mu E^2}{S_c}\delta_p^2\right]
 +2qd_v(E\delta_p-K+tv^2)^2-2E^2qv^2d_p,\\
 f_B={}&4\left[ctd_v+dj-2dt\delta_v
                    +\frac{dt^2}{S_c}\delta_v^2\right]
 +2y(t\delta_v-j)^2,\\
 f_C={}&2H^2d_pd_v,\qquad f_D=2HRd_pd_v,
 \qquad e=4Hd_pd_vg.
\end{align*}
All five are polynomials in $p,v$ with bounded source coefficients.
Their residues are, respectively, $2,2,2,2,8$. This assertion is uniform
over every opening permitted by the exact HI source entry.

\paragraph{The exact source contact jets.}
A star denotes $p=\sqrt q$, $v=\sqrt y$. Set
$\sigma=f_A+f_B+f_C+f_D-e$. The rational source formulas give
\begin{align}
 f_A^*+f_B^*&=f_C^*+f_D^*=e^*/2,\notag\\
 q^{1/4}(\partial_{\sqrt p}\sigma)^*
   &=-4(f_A^*-f_D^*),\notag\\
 y^{1/4}(\partial_{\sqrt v}\sigma)^*
   &=-4(f_B^*-f_D^*).                              \tag{N0M.6}
\end{align}
Also
\[
 C_s^*+4Hy(1+q)(1+y)=0,\qquad
 M_s^*-4Hq(1+q)(1+y)=0.                           \tag{N0M.7}
\]
These are actual-contact identities, not residue identities. The source-dependent identities hold modulo the exact numerator equation
for $n=0$, with $R$ retained rather than replaced by its residue.

\paragraph{Finite contact lift of the graph and both state normals.}
Here the SOS is a tensor SOS at each original pure-testing state. Its
roots need not be linear polynomials in $W,X,Y,Z$. If $sQ>0$, a common
sign change permits $s,Q>0$. This is a pure-testing sign choice, because
$sQ$ is a positive source factor times $-SY$. Put
\[
 A=Ps,\quad B=Vs,\quad C=TQ,\quad D=Q,
 \quad G=(ABCD)^{1/4}=\sqrt{TsQ},\quad E_i=A_i-G.
\]
Every $A_i$ and $G$ is a positive decomposable tensor. Thus every radical
below is a product of a pure-source radical and a pure-testing radical.
The basic circuit is
\[
 L=A^2+B^2+C^2+D^2-4G^2
   =\sum_iE_i^2+2G\sum_iE_i.                     \tag{N0M.8}
\]
For each pairing $ij|k\ell$,
\[
 G\sum_iE_i=G(\sqrt{A_i}-\sqrt{A_j})^2
 +G(\sqrt{A_k}-\sqrt{A_\ell})^2
 +2G\{(A_iA_j)^{1/4}-(A_kA_\ell)^{1/4}\}^2.
\]
Moreover,
\begin{align*}
 L&=(A-C)^2+(B-D)^2+2(\sqrt{AC}-\sqrt{BD})^2,\\
 (r-r')^2&=(r-G)^2+(r'-G)^2
             +2G(\sqrt r-\sqrt{r'})^2
 \quad(r=\sqrt{AC},\ r'=\sqrt{BD}),
\end{align*}
and the analogous identities hold for $BC|AD$. A finite average of
these decompositions gives strictly positive constant diagonal pivots
simultaneously on the $E_i$, the variance roots in (N0M.8), and
\[
 F_p=\sqrt{AC}-G,\qquad F_v=\sqrt{BC}-G.
\]
The exact half-angle conductors are
\[
 G(\sqrt p-q^{1/4})=q^{1/4}F_p,
 \qquad G(\sqrt v-y^{1/4})=y^{1/4}F_v.          \tag{N0M.9}
\]
Their coefficients are pure-source infinitesimals.

For completeness, the first-jet cancellation is explicit. Choose the
positive source unit $k=(f_C^*+f_D^*)/2\sim2$. Taylor-expand the bounded
coefficient polynomials in $z=\sqrt p,w=\sqrt v$ around
$z_0=q^{1/4},w_0=y^{1/4}$. All divided-difference coefficients remain
bounded, since these are monic polynomial divisions by pure-source
centers. Explicitly each monomial has the finite identity
\[
 z^iw^j=\sum_{a=0}^i\sum_{b=0}^j
 {i\choose a}{j\choose b}z_0^{i-a}w_0^{j-b}
 (z-z_0)^a(w-w_0)^b.
\]
Collecting total degrees zero, one, and at least two is the required
exact Taylor certificate. Its coefficients are bounded because all
exponents are nonnegative and all source coefficient denominators are
positive source units. No analytic remainder or mixed division enters.
After (N0M.9), every term in $\mathcal G-kL$ is an infinitesimal
bounded quadratic in the listed roots, except possibly
\[
 G\left[2\sum_i(f_i^*-k)E_i
       +z_0\sigma_z^*F_p+w_0\sigma_w^*F_v\right].
\]
By (N0M.6) this equals
\begin{align*}
 2(f_D^*-k)G\sum_iE_i
 &-\frac{z_0\sigma_z^*}{2}
       \{\sqrt G(\sqrt A-\sqrt C)\}^2\\
 &-\frac{w_0\sigma_w^*}{2}
       \{\sqrt G(\sqrt B-\sqrt C)\}^2.
\end{align*}
All three coefficients are source infinitesimals, and all their roots
already have strict pivots in the finite averaged circuit.

Finally the $N_c$ coupling in (N0M.4) is
$N_c(C_sB+\beta_cC)$, where
$\beta_c=4H\sqrt y\,v d_pd_v$ is infinitesimal.
The sum $C_s+\beta_c$ vanishes at the source contact by (N0M.7).
Writing $B=G+E_B$, $C=G+E_C$, monic-dividing that coefficient sum in
$z-z_0,w-w_0$, and applying (N0M.9) expresses this entire coupling by
infinitesimal bounded entries on $N_c,E_B,E_C,F_p,F_v$. The $N_m$
coupling is identical, using $A,C$ and
$\beta_m=-4H\sqrt q\,p d_pd_v$. Thus the two state normals have an
identity-residue block, and the graph has the strict circuit block.
The established strict bounded-Gram absorption lemma now gives an
actual tensor SOS of $F-\mathcal E$.

If $sQ<0$, the graph is instead a strict bounded diagonal form on the
four raw roots and $\sqrt{T|sQ|}$, with residue coefficients
$2,2,2,2,8$. The first-normal couplings already have infinitesimal raw
coefficients. If $sQ=0$, simply omit the last root. The same bounded-Gram
lemma applies directly. Consequently this argument covers literal
$n=0$ throughout the critical HI testing corner, at arbitrary independent
source/testing rates and through all positive source openings. It does
not assert the corresponding $n>0$ extension.



\section{\texorpdfstring{Exact source entry for the strong HI corner}{Exact source entry for the strong HI corner}}
\label{component:p54b8f20ad2dd}

Use the original positive rational C2 source with $d=1$, actual
$u_c$ unbounded and $v_c$ infinitesimal. Put
\[
 q=u_c^{-2},\quad y=v_c^2,\quad H=u_c^2h\in(0,1),
 \quad R=r/q,\quad B=b/q,\quad L=yt=B-H,\quad T=B-yR.
\]
The already proved source identities (WE1)--(WE3) show that $R$ is
bounded using only $n\geq0$, $q$ infinitesimal and $y$ bounded.
Here is their short contradiction argument, to make this use
explicit. The positive source gaps imply
\[
 0<T<H+y,\quad g=H+(1+y)r-1+qT>0,
\]
and $n=-qP/(2rQ)$ with $Q>0$ and
\[
 P\ge r^2(1+y)(H+yr)-Tqr(y+2)-q^2(T^2+T).
\]
If $q/r$ were infinitesimal, then $T/(H+yr)$ and
$(T^2+T)/(H+yr)$ would be bounded. For unit $r$ this follows
from $T<H+y\le(H+yr)/r$; for infinitesimal $r$, the inequality
$g>0$ forces $H$ to have residue one. Division by $r^2(H+yr)$
would give $P>0$, contrary to $n\geq0$. Thus $R$ is bounded.
(The exact positive remainder in this lower bound is printed in
(WE3), so this argument uses no asymptotic sign substitution.)

Now $r=qR$ is infinitesimal and $g>0$ gives $H_0=1$.
Define the positive opening
\[
 E=(1-H)/q.
\]
The same gap implies
$E<(1+y)R+T$, so $E$ is bounded. Since $j=1+R-t>0$,
$t$ is bounded; consequently
\[
 H=1-qE,\quad B=H+yt\sim1,\quad
 j=1+R-t,\quad K=B+R-E=g/q>0.
\]
All of $E,t,j,K$ are positive bounded, while
$t+j=1+R$ and $E+K=B+R$ are positive units.

Let $\Delta$ and $N$ denote the unchanged original source
polynomials, and set $\mathcal N=N/q$. Their exact expressions are
\begin{align*}
\Delta={}&(1+q)\{B^2+2BR+(1+y)B\}
 +(1-yq)R^2+2(1+y)R+1+y,\\
\mathcal N={}&(1+q)(B+R)(H+y)-(1+y)(1-qR)yt.
\end{align*}
Thus $\Delta$ is a positive bounded unit and
\[
 \bar\Delta=(1+\bar R)(3+\bar R),\qquad
 \overline{\mathcal N}=1+\bar R>0.
\]
Set $S=1+R+yt$, also a positive unit. Direct substitution into
the original coefficients gives the following exact equalities:
\[
\begin{aligned}
 mqK&=\frac{(1+y)(1-qR)(1+R)S}{2\Delta},&
 Cyt&=\frac{(1+R)\mathcal N}{2\Delta},\\
 \frac{A}{qy}&=\frac{(1+R)(B+R)S}{2H\Delta},&
 Dyj&=\frac{(1+q)(B+R)(B-yR)S}{2\Delta},\\
 \alpha qE&=\frac{(B+R)\mathcal N}{2\Delta},&
 n&=\frac{(B+R)(1+R)S-R\Delta}{2R\Delta}.
\end{aligned}                                                   \tag{HI-entry}
\]
Every one of the first five right-hand sides is a positive bounded
unit; in particular $B-yR$ has residue one. The last numerator has
residue $(1+\bar R)(1-\bar R)$, so $n\geq0$ implies
$0\le\bar R\le1$. This does not imply $R\le1$ as a source
inequality, and no such stronger claim is used.

Since $q,y$ are infinitesimal and $E,t,j,K$ bounded, the positive
numbers $qK,yt,yj,qE$ are infinitesimal. Dividing their unit
products in (HI-entry) proves
\[
 m,C,D,\alpha\quad\hbox{are all unbounded},\qquad
 A\asymp qy.
\]
This conclusion follows from the original source formulas and
$n\geq0$; no additional strong-weight assumption is needed for
an actual HI source.


\section{\texorpdfstring{A full-parent remote theorem for strong HI sources}{A full-parent remote theorem for strong HI sources}}
\label{component:pcee1372511ed}

\label{pcee1372511ed:sec:hi-remote}
\begin{theorem}[All tests outside the HI critical corner]
Let $K_c,K_t$ be independent real closed extensions of $\mathbb R$.
Consider the original C2 parent with $d=1$, $\lambda=\beta=0$,
$n,A\geq0$, and $m,C,D,\alpha$ positive unbounded source numbers.
Then at every projective testing pair except possibly
$1/u$ and $v$ both nonzero infinitesimal, and every real or complex testing
state, the
original parent is a sum of squares in $K_c\otimes_{\mathbb R}K_t$.
No comparison of source and testing rates is made.
\end{theorem}

\begin{proof}
For any positive real $\mu$, retain the four nonnegative event
packets with coefficients $m-\mu,C-\mu,D-\mu,\alpha-\mu$, and
retain the $n$ and $A$ packets in their entirety. It suffices to
prove the capped parent obtained by setting
$m=C=D=\alpha=\mu$ and $n=A=0$ is tensor SOS at the given test.
The real constant $\mu$ may be chosen after the real testing
residue is fixed, since each of the original four source
coefficients is larger than every real number.

Write $D_u=1+u^2$, $D_v=1+v^2$.
First suppose $v$ is infinitesimal, including zero. Use the roots
\[
 S=vW,\qquad E=Y-uW,\qquad X,\quad Z.
\]
All four capped event amplitudes are bounded linear forms in
these roots:
\[
 M=D_uS-X-uZ+uvE,\quad uW-Y=-E,\quad
 J=-E+uvX-vZ,\quad vY-Z=uS+vE-Z.
\]
The remaining core depends only on the same four roots, so the
exact capped parent has a bounded four-root Gram.
Let $a=\bar u$ and $d_a=1+a^2$. Its residue is
\[
 (8\mu+2)E^2+
 4\mu(d_aS-X-aZ)^2+(4\mu+2)d_a(aS-Z)^2
       +2(X+aZ)^2-2d_aZ^2.
\]
Put $T=X+aZ$ and $R=aS-Z$; this is an invertible real change of
coordinates. The last three-variable form is
\[
 4\mu(d_aS-T)^2+4\mu d_aR^2+2T^2
       -2d_a a^2S^2+4d_a aSR.
\]
The $T$ and $R$ pivots are $4\mu+2$ and $4\mu d_a$.
The remaining Schur pivot is exactly
\[
 \frac{d_a[4\mu^2-a^2(4\mu+1)]}{\mu(2\mu+1)}.
\]
Taking $\mu=1+a^2$ makes its numerator positive, because
$4\mu^2-a^2(4\mu+1)=3a^2+4$.
Thus this bounded Gram has strictly positive real residue and
is tensor SOS by the inherited strict-residue Gram lemma.

Next suppose $v$ is a bounded unit. At the real residues
$(a,b)=(\bar u,\bar v)$ the four event forms have common kernel
\[
 (W,X,Y,Z)=W(1,b,a,ab).
\]
The core restricted to this line is
$2(1+a^2)(1+b^2)b^2W^2$, which is strictly positive away from zero.
Consequently a sufficiently large positive real $\mu$ makes the
whole capped residue form positive definite. For completeness,
this elementary assertion has a direct quantitative proof.
Let $E_0\succeq0$ be the event Gram, split the state space into
$\ker E_0$ and its orthogonal complement, and let $\sigma>0$ be
its smallest eigenvalue on that complement. Let $\kappa>0$ be
the core's eigenvalue on the one-dimensional kernel and bound
the norms of the two remaining core blocks by $M$.
Then $4\mu\sigma>M+M^2/\kappa$ gives a positive Schur
complement. Every constant here belongs to $\mathbb R$.
The exact state Gram is bounded and has this positive residue.

Finally suppose $v$ is unbounded, including the projective
endpoint. Set $w=1/v$, homogenizing before setting $w=0$.
The four event forms become
\[
 W-wX+uY-uwZ,\quad uW-Y,\quad
 uwW+uX-wY-Z,\quad Y-wZ,
\]
with bounded positive factors $D_w$ or $D_u$ where appropriate.
At $w=0$ their common kernel is
$(W,X,Y,Z)=(0,X,0,aX)$, on which the homogenized core is
$2(1+a^2)X^2$. The same real Schur argument gives a positive
residue Gram for a sufficiently large $\mu$.
For finite unbounded $v$, undo homogenization by the pure
positive testing square $v^2$; for the projective endpoint the
homogeneous expression is already the required parent.
These three cases exhaust every projective Bob test when $u$ is bounded.

It remains to consider $p=1/u$ infinitesimal, after homogenizing
by $p^2$. If $v$ is a bounded unit, use the four roots
\[
 S=pY,\qquad T=vY-Z,\qquad W,\quad X.
\]
The exact event amplitudes are
\[
 W-S,\quad J=W+vX-D_vS+pvT,\quad
 M=pvW-pX+T,\quad T.
\]
The signed core also has bounded coefficients in these roots:
\[
 2J^2+2D_pT^2+2D_v
 [p^2X^2+2vSX-2pXT-v^2S^2+2pvST-p^2T^2].
\]
At $p=0$, set $b=\bar v\ne0$, $A_0=W-S$, and
$B_0=X-bS$. The full capped residue is
\[
 (8\mu+2)T^2+4\mu(1+b^2)A_0^2
 +(4\mu+2)(A_0+bB_0)^2
 +4(1+b^2)bSB_0+2(1+b^2)b^2S^2.
\]
As $\mu$ grows, the first three positive penalties constrain
$T=A_0=B_0=0$, and the residual kernel has positive coefficient
$2(1+b^2)b^2S^2$. The quantitative real Schur argument above
therefore gives a finite $\mu$ with a strictly positive residue.

If $u,v$ are both projectively unbounded, put $w=1/v$ and
homogenize by $p^2w^2$. Use $P=Y-wZ$, $S=pZ$, $W,X$.
Now the exact event amplitudes are
\[
 M=pW-pwX+P,\quad W-pP-wS,\quad
 J=wW+X-pwP-D_wS,\quad P.
\]
The exact core is
\[
 2J^2+2D_pP^2+2D_w[p^2X^2+2SX-S^2].
\]
At $p=w=0$ the capped residue is exactly
\[
 4\mu W^2+(8\mu+2)P^2+4\mu(X-S)^2+2X^2,
\]
which is strict for every $\mu>0$. The literal projective
endpoints cause no division. The literal critical axes also have direct exact certificates:
\begin{align*}
 F|_{p=0}={}&4nD_vX^2+(4m+4\alpha+2)(vY-Z)^2
 +4CD_vW^2+4AD_vY^2+(4D+2)(W+vX)^2,\\
 F|_{v=0}={}&4nD_pX^2+(4m+2)(pX+Z)^2
 +(4C+4D+2)(W-pY)^2+4AD_pY^2+4\alpha D_pZ^2.
\end{align*}
Thus only $p,v$ both nonzero infinitesimal remain untreated.
Applying the real-state result separately to real and imaginary
state components proves the complex-state claim.
\end{proof}

For a strong C2 HI source, the exact source-entry lemma below proves
that all four coefficients in the theorem are unbounded. Hence the
theorem applies without any additional source-shape assumption.
Its precise uncovered complement is the critical testing corner
$1/u$ and $v$ both nonzero infinitesimal. Exact sign gauges may
make both positive.
The original HI source leaf is therefore narrowed but not closed.


\section{\texorpdfstring{Weak-$C$ IH: the open testing axis}{Weak-C IH: the open testing axis}}
\label{component:pf95b1c5fc32f}

\label{pf95b1c5fc32f:sec:weak-ih-open-axis}
\paragraph{Domain and source input.}
The source and testing real closed fields are independent.  Let $B_c,B_t$
be their rings of bounded elements and use $B=B_c\otimes_{\mathbb R}B_t$
for Gram coefficients.  Source square roots and testing square roots are
always taken in their respective fields.  In the weak-$C$ IH source chart,
write $z=u_c^2$ infinitesimal, $y=v_c^2$ unbounded,
$H=zh\in(0,1)$, $B_*=zb$ unbounded and $\rho=B_*/y$ bounded.
The source entry calculation gives
\[
 0\leq\rho_0\leq1,\qquad
 (B_*C)_0=\left(\frac{AH}{B_*}\right)_0
 =\left(\frac{\alpha(1-H)}{B_*}\right)_0
 =c:=\frac1{2(1+\rho_0)}.
\tag{OA1}
\]
In particular $C$ is infinitesimal and $A,\alpha$ are unbounded.
For $0<\rho_0<1$,
\[
 D_0=\frac{\rho_0^2}{2(1-\rho_0^2)};
\tag{OA2}
\]
for $\rho_0=0$, $D$ is positive infinitesimal; for $\rho_0=1$, $D$ is
unbounded.  These statements follow from
$D\sim\rho^2/[2e(1+\rho)]$ and $e_0=1-\rho_0$.
No lower bound on $n\geq0$ is imposed.

\paragraph{Statement.}
For every testing $u$ with nonzero finite residue and every infinitesimal
$w=1/v$, including $w=0$, the complete reciprocal-$v$ canonical C2 parent
is a tensor sum of squares for every state.  The same assertion holds for
complex states by applying the real quadratic form to real and imaginary
parts separately.

\paragraph{Exact parent and retained squares.}
Put $D_u=1+u^2$, $D_w=1+w^2$ and
\[
 M=W-wX+uY-uwZ,\qquad J=uwW+uX-wY-Z.
\]
The parent is
\begin{align*}
\widehat F={}&4nD_uD_wX^2+4mM^2+4CD_w(uW-Y)^2
 +4AD_uD_wY^2+4DJ^2+4\alpha D_u(Y-wZ)^2\\
&+2\{J^2+D_u(Y-wZ)^2+D_w(X+uZ)^2-D_uD_wZ^2\}.
\end{align*}
Retain the entire $n,m$ packets. When $\rho_0=0$, retain the entire
$D$ packet as well and set $d=0$. When $0<\rho_0<1$, set $d=D$.
When $\rho_0=1$, choose a sufficiently large positive real $d$ below $D$
and retain $4(D-d)J^2$.  Denote the remainder by $F_*$.  All retained
terms are tensor squares or sums of tensor squares.

Use the graph $X=0$, $Y=HwZ$, and define
\[
 P=\frac{uW}{\sqrt{B_*}},\quad E=\sqrt{B_*}\,wZ,
 \quad V=uwW,\quad
 E_Y=\sqrt{\frac{B_*}{H(1-H)}}(Y-HwZ).
\tag{OA3}
\]
Crucially, the exact root relation is
\[
 PE=VZ.\tag{OA4}
\]
On the graph, the $C$, $A,\alpha$, $d$ packets are respectively
\begin{align*}
 &4B_*CD_w(P-HE/B_*)^2,\\
 &\frac{4D_u}{B_*}\{AH^2D_w+\alpha(1-H)^2\}E^2,\\
 &4d\{V-(1+Hw^2)Z\}^2.
\end{align*}
The remaining core is exactly
\[
 2V^2-4(1+Hw^2)PE
 +\frac{2}{B_*}\{H^2+u^2(1-H)^2+H^2w^2\}E^2.
\tag{OA5}
\]
Thus every graph coefficient is bounded.  The last $E^2$ coefficient is
infinitesimal, uniformly for all $0<H<1$.

\paragraph{Restoring the $Y$ normal.}
Let $L=B_*/[H(1-H)]$ and $Y_*=HwZ$.
The first normal, divided by four, is
\begin{align*}
 L_Y:=\frac14\partial_YF_*\big|_{X=0,Y=Y_*}
={}&2CD_w(Y_*-uW)+2D_u\{AH D_w-\alpha(1-H)\}wZ\\
&-2dw\{V-(1+Hw^2)Z\}
 -w\{V-(1+Hw^2)Z\}+D_u(H-1)wZ.
\tag{OA6}
\end{align*}
Every coefficient of $L_Y/\sqrt L$ in $(P,E,V,Z)$ is infinitesimal
when $d$ is bounded, including $d=0$.  Indeed the only potentially finite contribution is
\[
 2D_u\frac{AH-\alpha(1-H)}{B_*}\sqrt{H(1-H)}E,
\]
and it has zero residue by (OA1); the additional $w^2$ contribution is
infinitesimal.  The other terms contain $C$, $w/\sqrt L$ or $1/B_*$.
The quadratic $E_Y$ coefficient has residue $4cD_{u,0}>0$, and the
$X E_Y$ coefficient has zero residue.  These statements remain valid
when $H$ or $1-H$ is infinitesimal.  No inverse of $d$, $H$ or $1-H$ enters a Gram coefficient
unaccompanied by its compensating factor.

\paragraph{The strict residue when $0<\rho_0<1$.}
Write $s=\sqrt{1+u_0^2}>1$ and $r=\rho_0$.
Using (OA4), split the core cross term as
$-4PE=-4aPE-4(1-a)VZ$, where $a$ is any real number satisfying
\[
 \frac{s}{s+r}<a<\min\left\{\frac{s}{1+r},\frac{s}{s-r}\right\}.
\tag{OA7}
\]
The interval is nonempty: both upper endpoints exceed $s/(s+r)$.
The residue Gram splits into the positive $E_Y$ pivot, a $(P,E)$ block
\[
 G_{PE}=\begin{pmatrix}2/(1+r)&-2a\\-2a&2s^2/(1+r)\end{pmatrix},
\]
and a $(V,Z,X)$ block
\[
 G_{VZX}=\begin{pmatrix}
 2+4d_0&-4d_0-2(1-a)&2u_0(1+2d_0)\\
 -4d_0-2(1-a)&4d_0&-4d_0u_0\\
 2u_0(1+2d_0)&-4d_0u_0&2s^2+4d_0u_0^2
 \end{pmatrix}.
\tag{OA8}
\]
The $(V,X)$ principal block is positive definite.  Its Schur complement
is positive exactly when
\[
 \frac{s}{s+r}<a<\frac{s}{s-r},
\]
as direct substitution of $d_0=r^2/[2(1-r^2)]$ shows.  Hence (OA7)
makes both blocks strictly positive.  This checks the entire $X$ normal,
not merely the scalar graph.

\paragraph{The endpoint $r=1$.}
For finite $d>0$, positivity of the $(V,Z,X)$ block is equivalent to
$a_-<a<a_+$, where
\[
 a_\pm(d)=
 \frac{(2d+1)s^2\pm s\sqrt{2d(2d+1)}}{(2d+1)u_0^2+1}.
\]
As $d\to\infty$, the lower endpoint tends to $s/(s+1)$ and the upper
endpoint to $s/(s-1)$.  Choose first a real $a$ strictly between
$s/(s+1)$ and $\min\{s/2,s/(s-1)\}$, then a sufficiently large finite
$d$.  The interval is nonempty because $s>1$.  Since source $D$ is
unbounded this cap is admissible, and the full residue is strict.

\paragraph{The endpoint $r=0$.}
Take $a=1$, $d=0$, having retained the whole $D$ packet, and omit
$Z$ from the root list. The residue is the direct sum of
\[
 2(P-E)^2+2u_0^2E^2,
 \qquad 2(V+u_0X)^2+2X^2,
 \qquad 2D_{u,0}E_Y^2.
\tag{OA9}
\]
It is strictly positive. The exact rewriting has bounded coefficients:
(OA5) uses only $P,E,V$, and the first $X$ normal from the core is
$uV+u(1-H)w^2Z=uV+u(1-H)wE/\sqrt{B_*}$.
Thus no unscaled $Z$ or inverse of $D$ remains.

\paragraph{Conclusion.}
In all three cases the original remainder has a bounded Gram matrix with
strictly positive real residue.  The bounded-Gram lifting lemma therefore
makes it a tensor sum of squares.  Adding the retained source-positive
packets proves the assertion.  The argument is valid for arbitrary $n\geq0$,
including $n=0$, and literal projective $w=0$; no comparison between source
and testing infinitesimals was used.


\section{\texorpdfstring{Strong IH: the complete local corner}{Strong IH: the complete local corner}}
\label{component:pf3e0ca8bf054}

\label{pf3e0ca8bf054:sec:strong-ih-local}
Assume the strong IH source entry of \S\ref{pf5726218aec8:sec:strong-angle-entry}:
$z$ is infinitesimal, $y$ unbounded, $B=H+L$ a bounded positive unit,
$0<H<1$, $L>0$, and $C,\alpha$ strong. The source parameter
$\rho=B/y$ is infinitesimal, $s=r/\rho$ bounded, and $c_s=(1+L)/B$
is a positive unit. Testing $u,w=1/v$ are infinitesimal or zero.
We use the exact graph, five roots and both full normal identities
(IC3)--(IC12) of the preserved corner calculation; none of those
algebraic identities assumes $C$ weak.

The source quantities obey (GE3), and direct rational simplification
gives additionally
\[
 \frac{k}{B\rho}\longrightarrow
 \frac{2(B^2s+BHs+BH-Bs-H^2+H)}{B^2}.
\tag{SL1}
\]
No opening denominator occurs in this normalized expression.

Set $T_\rho=\rho T$ and
$\eta=\sqrt{H(1-H)/B}$, so $Y=HwZ+\eta E_Y$.
The full Gram's entire explicit dependence on $C$ is exactly
\[
 4C(U_1^2+U_2^2),\quad
 U_1=uW-Y-LwZ=\sqrt B R_0-\eta E_Y,
\]
\[
 U_2=w(uW-Y)-(L/y)Z
       =R_2-\sqrt{H/B}E-w\eta E_Y.
\tag{SL2}
\]
The scalar and normal identities verifying this assertion use
\[
 \sqrt H R_2-\sqrt{BH}\,wR_0-\sqrt B E=0.
\tag{SL3}
\]
Choose a positive real cap $c_*\le C$ and retain
$4(C-c_*)(U_1^2+U_2^2)$. In the remaining exact Gram replace explicit
$C$ by $c_*$, leaving $k$ and all other original source quantities
unchanged. The distinction matters: this is an exact square
subtraction from the complete Gram, not a perturbation of the source.
Add $\lambda$ times the exact zero (WIE2), where $\lambda>0$ is a
real number to be chosen. All resulting coefficients are bounded,
using (GE3), (SL1), and the opening-normalized roots.

We display the entire residue. The roots $R_1$ and $E_X$ have isolated
positive pivots $2c_s$ and $1$. The $(R_0,E_Y)$ block is
\[
 \begin{pmatrix}
 4Bc_*&-4c_*\sqrt{H(1-H)}\\
 -4c_*\sqrt{H(1-H)}&
 (4c_*+4A+4\alpha+2)H(1-H)/B
 \end{pmatrix}.
\tag{SL4}
\]
Here and below all expressions in a residue matrix mean their bounded
residues; products such as $AH$ are evaluated by (GE3), even when
$A$ itself is unbounded. Its determinant is exactly $8c_*(1+L)>0$.
Put $a_*=2(2c_*+1)$. The $(R_2,E)$ block is
\[
 \begin{pmatrix}
 a_*&-a_*\sqrt{H/B}\\
 -a_*\sqrt{H/B}&H(4A+a_*)/B
 \end{pmatrix},
\tag{SL5}
\]
with determinant $2a_*c_s>0$. The only remaining coupling is between
$R_2$ and $T_\rho$, whose Gram entry is
$c_s(H/B-s)$. The $T_\rho$ diagonal is $\lambda+2c_s s$.
Choose $\lambda$ larger than the finite Schur-complement threshold
of the positive block (SL5). This produces a strictly positive
residue matrix. Since this choice is made after taking a real
residue, it is a real scalar and introduces no mixed inverse.

We spell out why no normal term is hidden in this block description.
In (IC11), $f_X/\rho$ is bounded by (GE3) and every scalar-root
multiplier contains $u$ or $w$. In (IC12), the only nonzero limiting
term is the explicit $C$ contribution to $R_0E_Y$, already included
in (SL2)--(SL4). After the cap, all other normal multipliers contain
$z$, $u$, or $w$. The unbounded factors $A$ and $\alpha$ always occur
as $AH$ and $\alpha(1-H)$ after normalization, hence remain bounded
also at $H_0=0,1$. The normal $X$ coefficient divided by $N=4n+2$
has residue one, and its coupling to $E_Y$ vanishes. The remaining
$Y$ coefficient is exactly the lower-right entry of (SL4).
These facts follow directly from the displayed exact normal
factorizations, not from scalar-graph positivity.

The bounded positive-residue Gram lemma now gives tensor SOS for the
complete original parent, after restoring the retained positive
$C-c_*$ squares. This holds for all states and all independent local
tests, for both opening endpoints, for arbitrary $n\ge0$, and at
literal $u=0$ or $w=0$.


\section{\texorpdfstring{Strong IH: remote charts and the all-tests join}{Strong IH: remote charts and the all-tests join}}
\label{component:p7405dbf89341}

\label{p7405dbf89341:sec:strong-ih-remote}
Use the source entry of \S\ref{pf5726218aec8:sec:strong-angle-entry}. In particular
$B=H+L$ is a bounded positive unit, $m,D$ are infinitesimal,
$C,A,\alpha$ are strong, and the bounded products satisfy
\[
 CL\to\tfrac12,\qquad AH\to(1+L)/2,
 \qquad\alpha(1-H)\to B/2.
\tag{SR1}
\]
Whenever an event coefficient is unbounded, a positive finite real cap
may be used and its excess retained in its original positive square.
When a coefficient is bounded it may be used unchanged.

\paragraph{Finite-penalty principle.}
Suppose a real symmetric quadratic form is strictly positive on the
common kernel of finitely many real linear forms. Adding sufficiently
large positive multiples of their squares makes the complete form
strictly positive. To verify this without a limiting assertion, choose
an orthogonal splitting into their common kernel and its complement.
The upper block has a positive lower eigenvalue $c$, and the penalty
matrix on the complementary block has a positive lower eigenvalue
$e$. If the remaining lower block has norm at most $b$ and the cross
block norm is at most $a$, any penalty $M$ with $Me>b+a^2/c$ makes the
Schur complement strictly positive. All choices are real scalars.
This principle makes the opening cases below uniform; no unbounded
coefficient is inserted into a bounded Gram.

\paragraph{Testing $u$ infinitesimal, $v$ finite.}
Retain the $n,D$ packets and excess caps of $C,A,\alpha$.
Write their finite selected coefficients as $c,a,b>0$. Use roots
$U=uW,Y,Z,X,V=\sqrt m\,vW$. The $m$ packet is exactly
\[
 4(V-\sqrt m X+\sqrt m\,uvY-\sqrt m\,uZ)^2.
\]
Together with the selected packets and the original core this is an
exact bounded Gram expression. Its residue has isolated $X,V$ pivots
$2(1+v^2),4$ and the following $(U,Y,Z)$ block (residue notation suppressed):
\[
 \begin{pmatrix}
 2+4cD_v&-2-4cD_v&-2v\\
 -2-4cD_v&2D_v+4cD_v+4aD_v+4bv^2&-4bv\\
 -2v&-4bv&4b
 \end{pmatrix}.
\tag{SR2}
\]
In the interior, use $c=1/(2L)$, $a=(1+L)/(2H)$ and
$b=B/[2(1-H)]$. Its leading principal minors are
\[
 \frac{2(L+D_v)}L,\quad
 \frac{4(L+1)(D_v-H)(L+D_v)}{HL(1-H)},\quad
 \frac{8B(L+1)^2D_v}{HL(1-H)},
\tag{SR3}
\]
all positive. At an opening, apply the finite-penalty principle:
\begin{itemize}
\item $H_0=0$: the large $A$ packet forces $Y=0$; the restricted
      two-dimensional determinant is $4(L+1)>0$.
\item $H_0=1$: the large $\alpha$ packet forces $Z=vY$; the restricted
      determinant is $4(L+1)^2D_v/L>0$.
\item $L_0=0$: the large $C$ packet forces $U=Y$; the restricted
      determinant is $4/(1-H)>0$.
\item $L_0=0$, $H_0=1$: both constraints leave the strictly positive
      form $2Y^2$.
\end{itemize}
The case $L_0=H_0=0$ is excluded by $B_0>0$. Thus (SR2) can always be
made strictly positive with bounded selected coefficients.

\paragraph{Testing $u$ a unit, $w=1/v$ infinitesimal.}
Retain $n,m$ and the excess caps of $C,A,\alpha$; retain the full
$D$ packet using its natural root. Put
$U=uW$, $E=wZ$, $P_Z=\sqrt D Z$. In roots $(U,Y,E,X,P_Z)$ the exact
selected expression is
\begin{align*}
 &4cD_w(U-Y)^2+4aD_uD_wY^2+4bD_u(Y-E)^2\\
 &\quad+4\{\sqrt D(wU+uX-wY)-P_Z\}^2\\
 &\quad+2\{(wU+uX-wY)^2+D_u(Y-E)^2+D_wX^2-E^2
                   -2UE+2YE+2uwXE\}.
\tag{SR4}
\end{align*}
Its isolated $X,P_Z$ residue pivots are $2D_u,4$. The $(U,Y,E)$ block is
\[
 \begin{pmatrix}
 4c&-4c&-2\\
 -4c&4c+(4a+4b+2)D_u&-4bD_u-2u^2\\
 -2&-4bD_u-2u^2&4bD_u+2u^2
 \end{pmatrix}.
\tag{SR5}
\]
For interior openings the leading minors are
\[
 2/L,\qquad4(1+L)D_u/[HL(1-H)],\qquad
 8u^2(1+L)^2D_u/[HL(1-H)],
\tag{SR6}
\]
strict because $u_0\ne0$. The opening constraints and restricted
positive determinants are respectively
\[
 \begin{array}{c|c|c}
 H_0=0&Y=0&4u^2(1+1/L)\\
 H_0=1&Y=E&4u^2(1+L)/L\\
 L_0=0&U=Y&4u^2D_u/[H(1-H)]
 \end{array}
\]
and at $L_0=0,H_0=1$ the remaining form is $2u^2Y^2$.
Finite penalties therefore give a strict full Gram at every opening.

\paragraph{Both testing angles unbounded.}
Use (IR4)--(IR7), replacing $C$ by a finite positive cap $c_*$ and
setting $d=D$ (which is infinitesimal in the strong IH source).
Retain the excess $C$ packet and the entire $n,m$ packets.
The roots are
\[
 \begin{gathered}
 S=\sqrt{c_*}W,\quad E=\sqrt B\,wZ,\quad U=pwW,\quad P=\sqrt D\,pZ,\\
 E_X=X+wW,\qquad E_Y=\sqrt{B/[H(1-H)]}(Y-HwZ).
 \end{gathered}
\]
The exact graph identity (IR5) remains valid. The $E^2$ coefficient
has residue $2c_s$, $c_s=(1+L)/B$, because its core contribution
$2(1-H)^2/B$ must be retained when $B$ is a unit.
In (IR7), the potentially nonzero first-normal coefficient is now
\[
 2\{AH-\alpha(1-H)\}-(1-H),
\]
whose residue is zero by (SR1). The normal $E_Y$ pivot is also $2c_s$.
All scalar-normal cross coefficients are infinitesimal, since
$\sqrt D$ is infinitesimal. Thus the full diagonal residue is
\[
 \operatorname{diag}(4,2c_s,2,4,2,2c_s)>0.
\tag{SR7}
\]

\paragraph{The remaining remote charts.}
When $u$ is unbounded and $v$ is a unit, the homogeneous core has
residue $2\{(W+vX)^2+(vY-Z)^2\}$. Positive finite caps on the $C$
and $A$ packets supply $W^2$ and $Y^2$, respectively, so the residue
is strict on $(W,X,Y,Z)$. When $u$ is unbounded and $v$ infinitesimal,
put $p=1/u$ and use the five roots
$(P,Q,W,Y,Z)=(pX,W+vX,W,Y,Z)$. Positive caps $c,a,b$ give the residue
\[
 2Q^2+2(P+Z)^2+4cW^2+4aY^2+4bZ^2>0.
\tag{SR8}
\]
When $u$ is a unit and $v$ infinitesimal, the core has residue
$2\{(uW-Y)^2+(X+uZ)^2\}$; the positive $A$ and $\alpha$ caps give
$Y^2$ and $Z^2$, making the residue strict. In these expressions all
excess event coefficients are retained in their positive squares.
When both testing angles are units, the inherited noncommuting-test
theorem applies independently of source scales.

\paragraph{Exhaustive join.}
Classify each testing angle in the projective line as infinitesimal,
a bounded unit, or unbounded, with literal zero and infinity included
in the corresponding homogeneous chart. The infinitesimal/unbounded
cell is the local theorem of \S\ref{pf3e0ca8bf054:sec:strong-ih-local}.
The infinitesimal/finite cells follow from (SR2); the unit/unbounded
cell from (SR4); the unbounded/unbounded cell from (SR7); the
unbounded/finite and unit/infinitesimal cells from (SR8) and its
preceding argument; and the unit/unit cell from the inherited theorem.
These nine cells are exhaustive. Both exact sign gauges
$(u,W,X)\mapsto(-u,-W,-X)$ and
$(v,W,Y)\mapsto(-v,-W,-Y)$ preserve the full parent, so every
signed testing chart is included. Every certificate retains the full
state normals and the omitted original $\lambda,\beta$ remain zero.
Therefore the entire strong IH source leaf has all-tests tensor-SOS
admission, including $n=0$ and every projective endpoint.


\section{\texorpdfstring{Strong C2 source entry at an infinitesimal first angle}{Strong C2 source entry at an infinitesimal first angle}}
\label{component:pf5726218aec8}

\label{pf5726218aec8:sec:strong-angle-entry}
Let $z=u_c^2$ be infinitesimal, $y=v_c^2>0$ arbitrary, and use the exact
positive C2 source. Set
\[
 \begin{gathered}
 H=zh\in(0,1),\quad L=yt>0,\quad B=zb=H+L,\\
 A_c=1+zr,\quad S=A_c+L,\quad j=A_c-L/y>0.
 \end{gathered}
\]
Here $0<r<1$. Define
\begin{align*}
 \delta={}&z+(1+z)B+(1+z)B^2+2z(1+z)Br+z^3r^2+2z^2r\\
 &+y\{(1+z)B+z+z^2r(2-r)\},\\
 G={}&B+z(H+r-1),\\
 M={}&(1+z)(B+zr)(H+y)-z(1+y)(1-r)L\\
 ={}&(1+z)(B+zr)H+yG+yzj+zrL(1+y)>0.
\end{align*}
The unchanged original coefficient formulas are
\begin{align*}
 n&=\frac{(B+zr)A_cS-r\delta}{2r\delta},&
 m&=\frac{z^2(1+y)(1-r)A_cS}{2G\delta},\\
 C&=\frac{A_cM}{2L\delta},&
 A&=\frac{yA_c(B+zr)S}{2H\delta},\\
 D&=\frac{(1+z)(B+zr)(B-zyr)S}{2yj\delta},&
 \alpha&=\frac{(B+zr)M}{2(1-H)\delta}.
\tag{GE1}
\end{align*}
Assume $C$ and $\alpha$ are strong, meaning each has a positive real
lower bound in the coefficient field. Then $B$ is a bounded positive
unit. Indeed
\[
 \delta\ge(1+y)(B+z),\qquad
 M\le(1+z)(B+z)(1+y)
\]
give $CL\le A_c(1+z)/2$, so $L$ and $B$ are bounded.
If $B$ were infinitesimal then so would $H$, while
$\alpha(1-H)\le(B+zr)(1+z)/2$ would force $\alpha$ infinitesimal,
a contradiction.

Next, $n\ge0$ gives $r\delta\le(B+zr)A_cS$ and therefore
$yr\le A_cS$, since $\delta\ge y(B+zr)$.
Thus $G$, $B-zyr$ and $\delta/(1+y)$ are positive units.
Since $L<yA_c$, the ratio $(H+y)/(1+y)$ is a positive unit: if $y$
is infinitesimal then $H=B-L$ is a unit, and otherwise $y/(1+y)$
is a unit. The positive second expression for $M$ now shows that
$M/(1+y)$ is a positive unit.

Finally $1-r$ is a unit. To see this, let
$E(r)=(B+zr)A_cS-r\delta$, holding $B,H,y,z$ fixed.
Its coefficients divided by $1+y$ are bounded, while direct expansion
at $r=1$ gives
\[
 E(1)=-(H+y)(B+z)(1+z).
\]
The right side divided by $1+y$ is negative with nonzero residue.
If $1-r$ were infinitesimal, $E(r)/(1+y)$ would have the same negative
residue, contradicting $n\ge0$.

Consequently the following exact scaled quantities are positive units:
\[
 \frac m{z^2},\quad CL,\quad\alpha(1-H),\quad
 \frac{AH(1+y)}y,\quad Dyj(1+y).
\tag{GE2}
\]
For the IH source subcase $y$ is unbounded, hence $r=O(1/y)$ is
infinitesimal, $j$ is a unit, and $Dy^2$ is a positive unit.
Put $\rho=B/y$, $s=r/\rho$, and $c_s=(1+L)/B$.
Then $s$ is bounded and, with $N=4n+2$ and $\nu=Nr$,
\[
 \frac\nu\rho\longrightarrow2c_s,\quad
 \frac{nr}\rho\longrightarrow\frac{c_s-s}{2},\quad
 \frac D{\rho^2}\longrightarrow\frac{c_s}{2},\quad
 CL\longrightarrow\frac12,\quad
 \frac{Bm}{z^2}\longrightarrow\frac{c_s}{2},
\]
\[
 AH\longrightarrow\frac{1+L}{2},\qquad
 \alpha(1-H)\longrightarrow\frac B2.
\tag{GE3}
\]
Here arrows mean equality of residues of bounded expressions, with
$B,H,L,s$ replaced on the right by their residues. The expressions
remain valid at $H_0=0$, $H_0=1$, or $L_0=0$ whenever allowed by $B_0>0$.
The reason is that the scaled rational functions have denominators
bounded away from zero after the displayed opening factors have
been absorbed; no unscaled $1/H$, $1/(1-H)$ or $1/L$ is asserted bounded.

\label{component:pfb3a83f3849a}
\label{pfb3a83f3849a:sec:common-ii}
% Insertion after GE2 in the retained strong-angle source entry.
% Deliberately no component:pfb3a83f3849a label: the build already
% aliases that component to the retained GE source section.
\paragraph{The II source corollary.}
\label{sec:ii-source-bounds-refined}
For the original positive C2 source of GE1, suppose $z,y$ are
infinitesimal, without an initial strength assumption on $C$ or $\alpha$. Then $\alpha$ is strong if and only if $B=H+L$
is a positive unit.  On this face $C$ is unbounded, $H$ and
$1-r$ are positive units, and
\[
 \frac{m}{z^2},\quad\frac{A}{y},\quad CL,\quad Dyj,
 \quad\alpha(1-H)\asymp1,\qquad nr=O(1). \tag{DI7}
\]
Here every displayed comparison is in the coefficient field.
Indeed $0<L<yA_c$ makes $L$ infinitesimal and $B$ bounded.
If $B$ is infinitesimal, then $H$ is infinitesimal and GE1 gives
\[
 0<\alpha\le
 \frac{(B+zr)(1+z)(H+y)}{2(1-H)}\in\operatorname{Inf}_c,
\]
using $\delta\ge B+z$ and
$M\le(1+z)(B+zr)(H+y)$.  If $B$ is a unit, then
$H=B-L$ is a unit.  The positive expressions in GE1 show that
$G,\delta,M$ are units and hence that $CL$ and
$\alpha(1-H)$ are positive units.  Thus $C$ is unbounded
and $\alpha$ is strong, so the hypotheses used to prove GE2
hold.  That proof also gives $1-r$ a positive unit, from the
strictly negative value of $E(1)$ and $E(r)\ge0$.
Now GE2 yields the five bounds in DI7 because $H$ and $1+y$
are units.  Finally the exact GE1 identity
\[
 nr=\frac{(B+zr)A_cS-r\delta}{2\delta}
\]
is bounded.  No lower bound on $r$, $j$, or $1-H$ was imposed,
and the argument includes $n=0$.


\section{\texorpdfstring{Weak IH: all remote testing charts}{Weak IH: all remote testing charts}}
\label{component:p0fb32cd2f1bc}

\label{p0fb32cd2f1bc:sec:weak-ih-remote}
Assume the weak-$C$ IH source entry, with no restriction on
$\rho_0\in[0,1]$, on positive $r,H,1-H$, or on $n\ge0$.
The uniform source inputs are $B$ unbounded, $C,m$ infinitesimal,
$BC,AH/B,\alpha(1-H)/B$ positive units with common residue
$c=1/[2(1+\rho_0)]$, and $A,\alpha$ unbounded.
Every square retained below has nonnegative pure-source coefficient;
products $1+u^2$ and $1+v^2$ are themselves testing sums of squares.

\paragraph{Finite $v$, infinitesimal $u$.}
Write $Y_A=\sqrt A Y$ and $Z_\alpha=\sqrt\alpha(vY-Z)$, so
$Y=Y_A/\sqrt A$ and $Z=vY_A/\sqrt A-Z_\alpha/\sqrt\alpha$.
Retain the complete $n,C,D$ packets. Use the five roots
\[
 U=uW,\quad X,\quad V=\sqrt m\,vW,\quad Y_A,\quad Z_\alpha.
\]
The $m$ packet is exactly
$4(V-\sqrt m X+u\sqrt{m/\alpha}Z_\alpha)^2$.
All coefficients in the remaining core are bounded after substituting
$Y,Z$, and all coefficients involving $Y$ or $Z$ are infinitesimal.
The full residue Gram is diagonal:
\[
 2U^2+2(1+v_0^2)X^2+4V^2
       +4(1+v_0^2)Y_A^2+4Z_\alpha^2.
\tag{IR1}
\]
This includes literal $u=0$ or $v=0$; redundant zero roots cause no
problem for a positive Gram representation.

\paragraph{Finite $v$, unit $u$.}
Retain $n,m,C,D$. Use $(W,X,Y_A,Z_\alpha)$.
The residue of the core and selected $A,\alpha$ packets is
\[
 2\{u_0^2(W+v_0X)^2+(1+v_0^2)X^2\}
       +4(1+u_0^2)(1+v_0^2)Y_A^2
       +4(1+u_0^2)Z_\alpha^2,
\tag{IR2}
\]
which is strictly positive because $u_0\ne0$.

\paragraph{Finite $v$, unbounded $u$.}
Use the homogeneous coordinate $p=1/u$, infinitesimal or zero.
Retain $n,m,D$ and use
\[
 P=pX,\quad Q=W+vX,\quad S=\sqrt C W,\quad Y_A,\quad Z_\alpha.
\]
The $C$ packet becomes
$4(1+v^2)(S-p\sqrt{C/A}Y_A)^2$.
The entire selected expression has bounded coefficients, and residue
\[
 2Q^2+2(1+v_0^2)P^2+4(1+v_0^2)S^2
       +4(1+v_0^2)Y_A^2+4Z_\alpha^2.
\tag{IR3}
\]
Thus every finite-$v$ chart is complete, independently of $n$.

\paragraph{Both testing angles unbounded.}
Let $p=1/u$, $w=1/v$ be infinitesimal or zero and put
$D_p=1+p^2$, $D_w=1+w^2$.
Retain the entire $n,m$ packets. Set $d=\min(1,D)>0$ in the source
field and retain $4(D-d)J^2$. The remainder is exactly
\begin{align*}
 F_*={}&4CD_w(W-pY)^2+4AD_pD_wY^2+4dJ^2
       +4\alpha D_p(Y-wZ)^2\\
 &+2\{J^2+D_p(Y-wZ)^2+D_w(pX+Z)^2-D_pD_wZ^2\},\\
 J={}&wW+X-pwY-pZ.
\end{align*}
Use the graph $X=-wW$, $Y=HwZ$, the normals
$E_X=X+wW$, $E_Y=\sqrt L(Y-HwZ)$ with
$L=B/[H(1-H)]$, and scalar roots
\[
 S=\sqrt C W,\quad E=\sqrt B\,wZ,\quad
 U=pwW,\quad P=\sqrt d\,pZ.
\tag{IR4}
\]
On the graph the exact expression is
\begin{align*}
 F_*|_{\rm graph}={}&4D_w(S-pH\sqrt{C/B}E)^2
 +\frac{4D_p}{B}\{AH^2D_w+\alpha(1-H)^2\}E^2\\
 &+4(1+Hw^2)^2P^2+2D_wU^2
 -\frac{4pD_w}{\sqrt{BC}}SE
 +\frac2B\{(1-H)^2+H^2p^2D_w\}E^2.
\tag{IR5}
\end{align*}
All coefficients are bounded and its residue is
$4S^2+4cE^2+2U^2+4P^2$.

The first $X$ normal, divided by four, is exactly
\[
 L_X=-pD_wU+\frac{p(1-H)w}{\sqrt B}E
               -2\sqrt d(1+Hw^2)P.
\tag{IR6}
\]
The first $Y$ normal is
\begin{align*}
 L_Y={}&-2p\sqrt C D_w S+\frac{E}{\sqrt B}
 \{2Cp^2HD_w+2D_p(AHD_w-\alpha(1-H))\\
 &\hspace{38mm}+(2d+1)p^2(1+Hw^2)-D_p(1-H)\}.
\tag{IR7}
\end{align*}
After dividing by $\sqrt L$, all its coefficients are infinitesimal;
the only possible nonzero term is proportional to
$(AH-\alpha(1-H))/B$, whose residue is zero.
The quadratic normal coefficients have residues
$2+4d_0$ and $4c$, and their mixed coefficient has zero residue.
The only nonzero scalar-normal coupling is the $(P,E_X)$ block
\[
 \begin{pmatrix}4&-4\sqrt{d_0}\\-4\sqrt{d_0}&2+4d_0\end{pmatrix},
\tag{IR8}
\]
whose determinant is exactly $8$. Thus the full Gram is strict even
when $D$ is infinitesimal and even when $n=0$. The normalization uses
$\sqrt d$ but never $1/\sqrt d$ in a Gram coefficient.

Each displayed Gram is a representation of the complete selected
state quadratic form, not just its graph restriction. The bounded
positive-residue lemma and the retained squares prove tensor SOS.
These charts, the open-axis theorem, and the three local source-shape
proofs cover every independent testing pair.


\section{\texorpdfstring{The two missing weak-IH source endpoints}{The two missing weak-IH source endpoints}}
\label{component:p18dbedfff317}

\label{p18dbedfff317:sec:weak-ih-endpoints}
We use the exact parent, graph, roots, coefficients $K,k$, and both
normal factorizations (IC3)--(IC12) in the preserved interior-corner
proof.  Those identities are exact for every positive canonical C2
source; their validity does not require the interior-shape assumption.
Here $z$ and testing $u,w$ are infinitesimal (zero testing coordinates
are allowed), $B$ is unbounded, $H\in(0,1)$, $\rho=B/y$, and
$N=4n+2$, $\nu=Nr$. All source and testing rates are independent.

\paragraph{The endpoint $\rho_0=0$.}
Put $s=r/\rho$ and $\tau=1/B$. Nonnegativity of $n$ and the exact
source formulas imply that $s$ is bounded. The following normalized
source quantities have the indicated residues:
\[
 \nu/\rho\to2,\quad nr/\rho\to(1-s)/2,\quad
 D/\rho^2\to1/2,\quad BC,Bm/z^2,AH/B,\alpha(1-H)/B\to1/2,
\]
\[
 k/(B\rho)\to2s,\qquad 0\le s_0\le1.
\tag{WIE1}
\]
These are uniform opening identities: after substituting
$B=1/\tau$, $y=1/(\tau\rho)$, $r=s\rho$, their normalized rational
functions have a nonzero denominator at $z=\tau=\rho=0$, including
$H=0,1$. The fresh checker records these denominator residues.
Thus they concern simultaneous, not iterated, infinitesimal limits.

Replace $T$ by $T_\rho=\rho T$ in (IC7). Its coefficients are bounded;
for example $r^2n/\rho^2=s(nr/\rho)$ is bounded. Add the exact zero
\[
 T_\rho^2+(u^2-z)T_\rho R_2
 -\rho u(zR_0-uR_1)(uR_0-R_1)=0.
\tag{WIE2}
\]
This is (IC8a) multiplied by $\rho$, and uses no mixed inverse.
The residue Gram has $R_0,R_1,E$ diagonal entries $2,2,2$, and
\[
 (R_2,T_\rho):\quad
 \begin{pmatrix}2&-s_0\\-s_0&1+2s_0\end{pmatrix}.
\]
Its determinant $2+4s_0-s_0^2\ge2$ is strictly positive on $[0,1]$.

The entire normal restoration survives this rescaling. In (IC11),
replace $f_XT$ by $(f_X/\rho)T_\rho$, where
\[
 f_X/\rho=-2(nr/\rho)(1+w^2)-s(2D+1)
\]
is bounded. Every multiplier still contains testing $u$ or $w$, and
$1/\sqrt N$ is bounded. The remaining coefficients in (IC11) are
unchanged and infinitesimal. Formula (IC12) uses the root $R_1$ rather
than $T$ and remains bounded and infinitesimal unchanged. The normal
quadratic residue is $\operatorname{diag}(1,2)$. Hence the full
seven-root Gram is strict. This includes arbitrary positive $r,H,1-H$
and literal $n=0$.

\paragraph{The endpoint $\rho_0=1$.}
Set $e=1+zr-\rho+H/y=j>0$, an infinitesimal, and write
\[
 \rho=\frac{1+zr-e}{1-H\tau},\qquad \tau=1/B.
\]
The exact source formulas give the simultaneous normalized residues
\[
 De\to1/4,\quad\nu\to1,\quad nr\to(1-2r)/4,\quad
 BC\to1/4,\quad Bm/z^2\to(1-r)/4,
\]
\[
 AH/B,\alpha(1-H)/B\to1/4,\quad
 0\le r_0\le1/2.
\tag{WIE3}
\]
Put $\kappa=k-8DHr$; then $\kappa/B\to r$.
The scaled rational functions have regular nonzero denominators at
$z=\tau=e=0$ after the displayed powers of $e$ are absorbed.

The large coefficient $D$ must be retained in a square, not treated
as a bounded Gram coefficient. Introduce the full-state root
\[
 F_D=\sqrt D\,(J+eZ),\qquad a=\sqrt{H/B},\qquad
 Q=R_2-rT-aE.
\]
On the graph, $Q=J+eZ$. Subtracting its square from (IC7) gives the
exact identity
\begin{align*}
 G-4DQ^2={}&(4BC-\kappa z/B)R_0^2+4BmR_1^2/z^2
       +2(2C+1)R_2^2\\
 &+r^2(4nw^2+4n+2)T^2
 +\frac HB\{4A(1+u^2)+4C+2\}E^2\\
 &+\kappa uR_0R_1/B+(-4r+\kappa/B)R_2T
       -4a(2C+1)R_2E.
\tag{WIE4}
\end{align*}
Add the exact zero (IC8a). Restore the full normals
\[
 E_X=\sqrt N(X+ruZ),\qquad
 E_Y=\sqrt{B/[H(1-H)]}(Y-HwZ).
\]
Use the exact relation
\[
 R_2=F_D/\sqrt D+rT+aE-uE_X/\sqrt N
                         +wE_Y/\sqrt{B/[H(1-H)]}.
\tag{WIE5}
\]
The inverse $1/\sqrt D$ is purely in the source field and infinitesimal.
All coefficients after this substitution are bounded.

For completeness, the first-normal quotients of
$\widehat F-4F_D^2$ are precisely the D-removed versions of
(IC11)--(IC12). Explicitly,
\begin{align*}
 \bar f_X&=-2nr(1+V)-r,\\
 \bar g_X&=-2mB/z-2nr(1+z)+1-H-r-2m(H+r-1),\\
 \bar f&=2m-2C(1+V)-V,\\
 \bar f_Y&=2AH(1+V)-2\alpha(1-H)+(2m+1)(H+r-1),\\
 \bar g_Y&=-B(2C+1)+2AH(1+z)+2CH+H.
\end{align*}
The constant term in the $Y$ monic polynomial acquires $2De$, but
its monic quotient coefficients are unchanged; the $X$ polynomial
has the constant correction $-2De$. Exact source stationarity
still sets the constant remainders to zero because $J+eZ$ vanishes
at the contact. Thus the displayed quotient formulas reconstruct
both linear normals. Their normalized coefficients are infinitesimal
by (WIE3); in particular $De$ is bounded and no bare $D$ survives.
The normal quadratic block of the D-removed expression has residue
$\operatorname{diag}(1,1)$.

In the order $(R_0,R_1,T,E,F_D,E_X,E_Y)$, the complete residue Gram
is diagonal:
\[
 \operatorname{diag}\{1,1-r_0,1+r_0-r_0^2,1,4,1,1\}.
\tag{WIE6}
\]
For instance the $T^2$ entry comes from substituting $R_2=r_0T$
in (WIE4): it is $r_0\nu_0-r_0^2$, plus the unit contribution from
(IC8a). Every entry is positive for $0\le r_0\le1/2$.
This proves full-parent tensor SOS at the corner also when
$1-\rho$ is infinitesimal. Along with the preserved interior proof
and the $\rho_0=0$ argument, every allowed source shape is covered
at the independent testing corner $u,w$ infinitesimal.


\section{\texorpdfstring{An exact full-parent Gram at the interior IH double corner}{An exact full-parent Gram at the interior IH double corner}}
\label{component:pc00b84929b12}

\label{pc00b84929b12:sec:c2-weak-ih-interior-corner}

This is a local testing-chart theorem. It does not admit its source
subdomain at all independent tests. In particular, no conclusion for
remote commuting testing corners is included below.

\begin{theorem}[Full-parent IH corner, interior shape]
Use the exact positive C2 source with $d=1$, the omitted original
$\lambda,\beta$ coefficients zero, and $n\ge0$. Suppose its actual
angles have type IH, and $C$ is weak. Use the source notation of
\S\ref{p93b05ea6ad1e:sec:c2-weak-ih-frontier}; thus $z$ is infinitesimal and
$B$ unbounded. Write $\rho=B/y$ (a shape parameter, not an event
coefficient). Assume
\[
 \rho,\quad1-\rho
 \quad\hbox{are positive bounded units}.                 \tag{IC1}
\]
At every independent testing pair for which $u$ and the homogeneous
reciprocal coordinate $w=1/v$ are infinitesimal or zero, the entire
original parent has a separated-field tensor-SOS certificate, for
all real or complex state coordinates. The source openings $H,1-H$ and the positive radial parameter $r$
are arbitrary; $n$ may be unbounded. Literal
testing $u=0$ and $v=\infty$, either testing sign, and literal source
$n=0$ are included.
No state coordinate or certificate root is required to be bounded.
\end{theorem}

\begin{proof}
Put $N=4n+2$ and $\nu=Nr$. Source reduction (WI1)--(WI3)
shows that $\nu$ and $1-r$ are bounded positive units, even when
$N$ is unbounded. In residue,
\[
 \nu_0=\frac{2\rho_0}{1+\rho_0},\qquad 2r_0\le\nu_0,
 \qquad D_0=\frac{\rho_0^2}{2(1-\rho_0^2)}.             \tag{IC2}
\]
All assumptions and divisions in this sentence concern the coefficient
field alone.
The uniform coefficient estimates needed below are
\[
 BC\asymp1,\quad \frac{Bm}{z^2}\asymp1,\quad
 \frac{AH}B\asymp1,\quad\frac{\alpha(1-H)}B\asymp1,
 \quad D\asymp1,\quad nr\le\nu/4.
\]
In particular $C,m$ are infinitesimal, $nr,nr^2$ are bounded,
$mB/z\asymp z$, and $m\sqrt B/z\asymp z/\sqrt B$.
These estimates are uniform in $r,H,1-H$ and do not assert that
$A/B$ or $\alpha/B$ is bounded at an opening.

The homogeneous testing parent is exactly
\begin{align*}
 \widehat F={}&4nD_uD_wX^2+4mM^2+4CD_w(uW-Y)^2
     +4AD_uD_wY^2+4DJ^2+4\alpha D_u(Y-wZ)^2\\
 &+2\{J^2+D_u(Y-wZ)^2+D_w(X+uZ)^2-D_uD_wZ^2\},\\
 M={}&W-wX+uY-uwZ,\qquad
 J=uwW+uX-wY-Z,
 \qquad D_u=1+u^2,\quad D_w=1+w^2.
                                                               \tag{IC3}
\end{align*}
For $w\ne0$ it is $w^2F(u,1/w)$; (IC3) defines the literal
projective endpoint as well. Let $G$ be (IC3) restricted to the graph
\[
                         X=-ruZ,\qquad Y=HwZ.
\tag{IC4}
\]
Use the following five scalar roots, all defined using only pure
coefficient radicals and inverses:
\begin{align*}
 R_0&=(uW-BwZ)/\sqrt B,&
 R_1&=(zW-BuwZ)/\sqrt B,\\
 R_2&=uwW-\rho Z,&
 T&=(u^2-z)Z,&
 E&=\sqrt{BH}\,(w^2-1/y)Z.                               \tag{IC5}
\end{align*}
They vanish at every signed exact source contact. Define the pure
coefficient quantities
\begin{align*}
 K={}&4AH^2+8DHr+4\alpha(1-H)^2
      +4m(H+r-1)^2+2(H+r-1)^2+4nr^2,\\
 k={}&K-4B^2m/z^2+8znr^2+8AH^2/y.                      \tag{IC6}
\end{align*}
The exact source formulas and state/angular stationarity give the
following finite polynomial Gram identity:
\begin{align*}
G={}&\left(4BC-\frac{kz}{B}\right)R_0^2
     +\frac{4Bm}{z^2}R_1^2+2(2C+2D+1)R_2^2\\
 &+r^2\{4nw^2+2(2D+2n+1)\}T^2
  +\frac H B\{4Au^2+2(2A+2C+2D+1)\}E^2\\
 &+\frac{ku}{B}R_0R_1
  +\left\{-4r(2D+1)+\frac{k}{B}\right\}R_2T
  -\frac{4\sqrt H(2C+2D+1)}{\sqrt B}R_2E.                   \tag{IC7}
\end{align*}
For clarity, the only non-obvious algebraic step is the exact toric
conductor
\[
 B(u^2-z)(w^2-1/y)Z^2
 =(u^2-z)ZR_2-zR_0^2+uR_0R_1.                         \tag{IC8}
\]
Before (IC8), the last mixed radial term is
$k(u^2-z)(w^2-1/y)Z^2$. Expanding (IC7), its difference from
$G$ has only the four monomials $uwWZ,u^2Z^2,w^2Z^2,Z^2$;
the unchanged original source sets all four coefficients exactly to
zero. The supplied checker verifies these four rational identities,
not just a specialization or a leading-order equality.

All coefficients of (IC7) are bounded: $BC$, $Bm/z^2$, $AH/B$,
$\alpha(1-H)/B$, $nr,D$ and $k/B$ are bounded. No inverse of $r$
is used, and $1/\sqrt N$ is bounded because $N\ge2$. Moreover,
\[
 (k/B)_0=\frac{2r_0}{1+\rho_0}.
\]
Add to the right side of (IC7) the following exact zero, obtained
by one more toric conductor:
\[
 \rho T^2+(u^2-z)TR_2
       -u(zR_0-uR_1)(uR_0-R_1)=0.                    \tag{IC8a}
\]
The new $T^2$ term has positive unit residue $\rho_0$. Every other
coefficient introduced by (IC8a) is infinitesimal. This step is what
makes the proof uniform when $r$ is infinitesimal and $n$ unbounded.

Suppress residue bars temporarily. The $R_0,R_1,E$ diagonal
residues are respectively
\[
 \frac2{1+\rho},\qquad
 \frac{2(1-r)}{1+\rho},\qquad
 \frac2{1+\rho},
\]
which are positive. The only surviving off-diagonal block is
\[
 (R_2,T):\qquad
 \begin{pmatrix}
 \dfrac2{1-\rho^2}&-\dfrac r{1-\rho}\\[3pt]
 -\dfrac r{1-\rho}&\rho+r\nu+4Dr^2
 \end{pmatrix}.                                      \tag{IC9}
\]
Using (IC2), its determinant is exactly
\[
 \frac{2\rho}{1-\rho^2}
  +\frac{r\{4\rho-r(3\rho+1)\}}
       {(1-\rho)(1+\rho)^2}>0.
\]
Indeed $0\le r\le\rho/(1+\rho)$, and consequently
$4\rho-r(3\rho+1)\ge\rho(3+\rho)/(1+\rho)>0$.
The first determinant summand is already strictly positive, including
$r_0=0$. Thus (IC7) after the exact zero (IC8a) has strictly positive
definite residue, including literal $n=0$ and unbounded $n$.

It remains to restore the full state normals; the graph certificate
alone is not the claimed conclusion. Put
\[
 E_X=\sqrt N(X+ruZ),\qquad E_Y=\sqrt{B/[H(1-H)]}(Y-HwZ).
\]
We give exact factorizations of both linear normal terms. Let
$U=u^2,V=w^2,s=1/y$ and define $L_X,L_Y$ by
$4L_X=\partial_X\widehat F|_{\rm graph}$ and
$4L_Y=\partial_Y\widehat F|_{\rm graph}$. Direct expansion gives
\begin{align*}
 L_X={}&(2D+1)uR_2-\frac{2m\sqrt B}{z}wR_1+uZP_X(U,V),\\
 P_X={}&(2D+1)\rho-\frac{2mB}{z}V
    -2nr(1+U)(1+V)-r(2D+1)U-2D-r\\
 &+V\{1-H-r-2DH-2m(H+r-1)\}.                         \tag{IC10}
\end{align*}
Every coefficient of $P_X$ is bounded, and source stationarity gives
$P_X(z,s)=0$. Write its exact monic division as
\[
 P_X=(U-z)f_X(V)+(V-s)g_X.
\]
Both $f_X,g_X$ have bounded coefficients. Explicitly,
\begin{align*}
 f_X(V)&=-2nr(1+V)-r(2D+1),\\
 g_X&=-2mB/z-2nr(1+z)+1-H-r-2DH-2m(H+r-1).
\end{align*}
Thus no inverse of $r$ or of either opening is introduced. We obtain
\[
 L_X=\left(2D+1+\frac{g_X}B\right)uR_2
       -\frac{2m\sqrt B}{z}wR_1
       +u f_X(V)T
       -\frac{uwg_X}{\sqrt B}R_0.                       \tag{IC11}
\]
Here we used the exact identity
$B(V-s)Z=R_2-\sqrt B\,wR_0$; this avoids dividing by $\sqrt H$
when $H$ is infinitesimal. Every displayed multiplier is infinitesimal
in the bounded tensor ring; in particular $m\sqrt B/z\asymp z/\sqrt B$.

For the other normal, put
\begin{align*}
 f(V)={}&2m-2C(1+V)-(2D+1)V,\\
 P_Y(U,V)={}&Bf(V)+2AH(1+U)(1+V)+2CH(1+V)
       +2DHV+2DrU+2D\\
 &-2\alpha(1-H)(1+U)+2m(H+r-1)U
       +(H+r-1)U+H(V+1).
\end{align*}
Then
\[
 L_Y=\sqrt B f(V)R_0+wZP_Y(U,V),\qquad P_Y(z,s)=0.
\]
The coefficients of $P_Y/B$ are bounded. Write
$P_Y=(U-z)f_Y(V)+(V-s)g_Y$. The two quotient coefficients are explicitly
\begin{align*}
 f_Y(V)&=2AH(1+V)+2Dr-2\alpha(1-H)
                   +(2m+1)(H+r-1),\\
 g_Y&=-B(2C+2D+1)+2AH(1+z)+2CH+2DH+H.
\end{align*}
The displayed source estimates show directly that $f_Y/B$ and
$g_Y/B$ have bounded coefficients, including both opening limits.
The elementary exact identity
\[
          (U-z)\sqrt B\,wZ=zR_0-uR_1
\]
therefore gives
\[
 \frac{L_Y}{\sqrt{B/[H(1-H)]}}
 =\sqrt{H(1-H)}\left\{f(V)+\frac{z f_Y(V)}B\right\}R_0
       -\frac{\sqrt{H(1-H)}u f_Y(V)}B R_1
       +\frac{\sqrt{1-H}w g_Y}B E.          \tag{IC12}
\]
Again every displayed multiplier is infinitesimal. This argument
retains the full $X,Y$ linear terms; none is estimated on the graph
and then discarded.

Finally the normal quadratic block, in $(E_X,E_Y)$, is bounded.
Its residue is diagonal with entries
\[
             1,\qquad \frac2{1+\rho_0}>0.
\]
Indeed the $X^2$ coefficient before scaling is
$4nD_uD_w+4mw^2+(4D+2)u^2+2D_w$; the $Y^2$ coefficient is
$4CD_w+4AD_uD_w+(4D+2)w^2+(4\alpha+2)D_u+4mu^2$.
Their mixed coefficient is $-4uw(2m+2D+1)$.
For the $Y$ normalization, if $L=B/[H(1-H)]$, then $A/L$ and
$\alpha/L$ are bounded and
\[
 \frac{4(A+\alpha)}L
 =4\left\{\frac{AH}B(1-H)
               +\frac{\alpha(1-H)}B H\right\}
 \longrightarrow\frac2{1+\rho_0}.
\]
Thus the stated normal residue remains strict even when $H_0=0$
or $H_0=1$; no denominator involving a vanishing opening was dropped.
Equations (IC11)--(IC12) show that all cross terms from these normals
to the five scalar roots have zero residue. Thus the exact seven-root
Gram for the entire (IC3) is bounded and has positive definite
residue. The bounded positive-residue Gram lemma yields tensor SOS.
Its roots may be unbounded; it uses no mixed inverse or radical.
The homogeneous chart and real/imaginary splitting give all endpoints
and states stated in the theorem.
\end{proof}

\paragraph{Precise boundary of this result.}
The theorem covers only the independent testing corner $(u,v)=(0,\infty)$
in residue. It neither assumes nor concludes all-tests admission.
The proof does not cover source $\rho$ infinitesimal or
$1-\rho$ infinitesimal. The positive radial parameter $r$ and both
positive openings $H,1-H$ may be arbitrarily infinitesimal; $n$ may
be unbounded.
The local proof does include literal $n=0$, but a remote
infinity-infinity joining Gram has a pivot proportional to $n$;
that join has not been made uniform as $n\to0$. No source leaf is
removed from the all-tests frontier on the basis of this theorem.


\section{\texorpdfstring{The genuine directed IH weak-$C$ frontier}{The genuine directed IH weak-C frontier}}
\label{component:p93b05ea6ad1e}

\label{p93b05ea6ad1e:sec:c2-weak-ih-frontier}

\noindent\textbf{Scope of retained extract.} The current proof imports the complete exact source and opening/coquotient argument (WI1)--(WI3) below. Optional rational-curve examples and historical frontier prose are outside the mathematical scope used in this rederivation and are omitted from this presentation. The complete original file and the exact omitted spans remain in the source archive and presentation manifest.



Write $z=u_c^2$, $y=v_c^2$, and use the exact contact quantities
\[
 H=v_cY_c=zh,\quad L=v_c(u_cW_c-Y_c)=yt,\quad
 B=H+L=u_cv_cW_c=zb,\quad R=-u_cX_c=zr.
\]
Thus $0<H<1$, $0<R<z$, $L>0$, and
\[
 L<y(1+R),\quad B>yR,\quad
 G=B+R-z(1-H)>0.
\]
Set $S=1+R+L$ and
\begin{align*}
 \widetilde\Delta={}&(1+z)B^2+2(1+z)BR+(1+z)(1+y)B\\
 &+(z-y)R^2+2z(1+y)R+z(1+y),\\
 \widehat N={}&(1+z)(B+R)(H+y)-(1+y)(z-R)L.
\end{align*}
These are exactly $z\Delta,z^2N_{\rm src}$. In particular,
\[
 \widehat N=(1+z)(B+R)H+yG+z\{y(1+R)-L\}+RL(1+y)>0.
\]
The original coefficients, without an approximation, are
\begin{align*}
 C&=\frac{(1+R)\widehat N}{2L\widetilde\Delta},&
 \alpha&=\frac{(B+R)\widehat N}{2(1-H)\widetilde\Delta},\\
 A&=\frac{y(1+R)(B+R)S}{2H\widetilde\Delta},&
 D&=\frac{(1+z)(B+R)(B-yR)S}
 {2\{y(1+R)-L\}\widetilde\Delta},\\
 m&=\frac{z(1+y)(z-R)(1+R)S}{2G\widetilde\Delta},&
 n+\frac12&=\frac{z(B+R)(1+R)S}{2R\widetilde\Delta}.
                                                        \tag{WI1}
\end{align*}
The checker verifies all six against the unchanged original source.

\paragraph{Uniform source reduction for IH.}
Suppose $z$ is positive infinitesimal, $y$ is positive unbounded,
and $n\ge0$. Then $C$ is weak if and only if $B$ is unbounded.
To prove the nontrivial direction, suppose $B$ is bounded. Since
$R<z$ and $G>0$, one has $z\le4(B+R)$: use $B\ge1/2$
or, otherwise, $B+R>z(1-H)>z/2$. Direct division in (WI1)
gives bounded $\widetilde\Delta/[y(B+R)]$. Moreover,
\[
 \frac{\widehat N}{y(B+R)}
 =(1+z)(1+H/y)
 -(1+1/y)\frac{(z-R)L}{B+R}
\]
has residue $1$, since $L\le B$ and $z$ is infinitesimal.
Consequently $C\ge c/L$ for a positive real $c$, which is strong
because $L$ is bounded.

Now let $B$ be unbounded. Put
\[
 \rho=B/y,\qquad e=1+R-L/y>0.
\]
Here $0<\rho<1+R+H/y$, so $\rho$ is bounded and its residue
lies in $[0,1]$. The following exact normalized quantities have
residue $1$:
\[
 \frac{\widetilde\Delta}{B(B+y)},\qquad
 \frac{\widehat N}{yB},\qquad
 \frac{L}{B},\qquad \frac{S}{B},\qquad\frac{G}{B}.
                                                        \tag{WI2}
\]
They give
\[
 n+\tfrac12=\frac{\rho}{2r(1+\rho)}(1+o(1)).
\]
The inequality $n\ge0$ therefore gives $r/\rho$ bounded and
$r\le\rho/[1+\rho](1+o(1))\le1/2+o(1)$.
In particular $1-r$ is a positive unit, and
$R/\rho=z(r/\rho)$ is infinitesimal. The normalized leading
relations, uniform at arbitrary openings $H,1-H,e$, are
\begin{align*}
 C&=\frac{1+o(1)}{2B(1+\rho)},&
 A&=\frac{B(1+o(1))}{2H(1+\rho)},\\
 \alpha&=\frac{B(1+o(1))}{2(1-H)(1+\rho)},&
 m&=\frac{z^2(1-r)(1+o(1))}{2B(1+\rho)},\\
 D&=\frac{\rho^2(1+o(1))}{2e(1+\rho)}.
                                                        \tag{WI3}
\end{align*}
Here each $1+o(1)$ means a positive coefficient unit of residue one;
it never denotes a comparison with a testing infinitesimal. In
particular $A,\alpha$ are unbounded, $m,C$ are infinitesimal, and
$Cm$ is infinitesimal for every weak-$C$ IH source. Thus the four
reciprocal-product exit can never admit this entire directed leaf.
If $\rho$ is infinitesimal then $e$ is a unit and $D$ is weak;
if $\rho$ is a unit then $D$ is strong. Literal $n=0$ remains
allowed throughout the reduction. No lower bound on $r$ was used.



\section{\texorpdfstring{Completion of the strong C2 ordinary--small source chart}{Completion of the strong C2 ordinary--small source chart}}
\label{component:p6b22e944bd7c}

\label{p6b22e944bd7c:sec:ui-all-tests-recovered}

\begin{theorem}[Strong C2 UI, all independent tests]
Use the original C2 source (UR1)--(UR3), with $d=1$, actual source
angles $u_c=a>0$ a bounded unit and $v_c=k>0$ infinitesimal, and
$C,\alpha$ strong. Then the entire parent (UR6) is a sum of squares
in the tensor product of the independent source and testing fields,
for every testing angle, including projective endpoints, and every
homogeneous state. The statement retains the literal boundary $n=0$
and arbitrary positive rates of every source opening.
\end{theorem}

The local contact certificate in Section~\ref{p76545700d90b:sec:local-ui-certificate}
already covers $v$ infinitesimal and $u$ with the same residue as $a$.
We supply its complete complement. No comparison between $v$ and $k$
is made. The argument uses one elementary rank-one allocation, followed
by short positive frames at the testing endpoints.

\paragraph{A rank-one allocation lemma.}
Let $P$ be a positive definite real symmetric matrix, let $\ell,b$ be
real columns, and let $H,c_0>0$. Set
\[
 d=\ell^TP^{-1}\ell,\qquad e=\ell^TP^{-1}b,\qquad
 h=c_2-b^TP^{-1}b.
\]
Suppose
\[
 p(T)=c_0+(h+4Hdc_0)T+4H(dh+e^2)T^2>0\quad(T\ge0),
 \qquad dh+e^2>0.                                      \tag{UA1}
\]
There exists a real $\lambda$ such that
\[
 |\lambda|<\sqrt{4Hc_0},\qquad
 h+2e\lambda-d\lambda^2>0.                             \tag{UA2}
\]
Indeed put $s=\sqrt{4Hc_0}$. If $d>0$ and $|e|<ds$, take
$\lambda=e/d$. If $|e|\ge ds$, the value at the appropriate endpoint
is $h+2|e|s-ds^2$. Were this nonpositive, writing
$x=|e|-ds\ge0$ would give
$h+ds^2\le-2sx$ and $dh+e^2\le x^2$; the quadratic in (UA1)
would then have a nonpositive value at a positive argument. Thus a
nearby interior $\lambda$ works. This includes equality $|e|=ds$.
The scalar interval argument also holds in a limiting case $d=0$;
then $e\ne0$ follows from $dh+e^2>0$. (For a finite positive $P$
itself, $d=0$ would force $\ell=e=0$, so this case occurs only
when using the continuous scalar extension at a boundary.)

Consequently the two real matrices
\[
 \begin{pmatrix}P&b-\lambda\ell\\(b-\lambda\ell)^T&c_2\end{pmatrix},
 \qquad
 \begin{pmatrix}4H&\lambda\\\lambda&c_0\end{pmatrix}       \tag{UA3}
\]
are positive definite. This is an elementary real lemma. Its inverse
and square root are never taken in the tensor ring.

\paragraph{The finite ordinary angle away from contact.}
Assume $v$ infinitesimal and $u$ a bounded unit whose real residue has
square different from $a_0^2$. Keep the centered packets
$P_L,P_C,P_n,P_m,P_\alpha$ of (UR12), (UR14) outside the Gram, and
combine $P_H$ with the exact remainder. Use the six separated roots
\[
 r_0=uW-Y,\quad X,\quad Z,\quad V_0=vW,\quad K_0=kW,
 \quad R_0=\frac{v(uX-Z)}{k}.                           \tag{UA4}
\]
They satisfy the exact zero relation
\[
 R_0K_0-(uX-Z)V_0=0.                                  \tag{UA5}
\]
All divisions in (UA4) are pure source divisions.

For completeness the boundedness of the $C,D$ remainder in this frame
is explicit. Write $\ell_0=L\tau$, $c_R=C\rho k^2$, and
$Q=uX-Z$. Then (UR13) gives exactly
\begin{align*}
 G_{CD}+P_H={}&4\gamma R_0^2
 +8\ell_0kK_0r_0+8\ell_0qk^2V_0Q-4\ell_0\tau k^2K_0^2\\
 &+8c_RvV_0r_0-4c_R\rho k^2V_0^2.                  \tag{UA6}
\end{align*}
All displayed coefficients are bounded by (UR9)--(UR11); all terms
except the first have zero residue. The $A$ packet is exactly
$4(A/k^2)D_uD_v(uK_0-kr_0)^2$. Substituting
$Y=uW-r_0$ in (UR15) and the signed core gives a bounded Gram on
(UA4), with no discarded normal term. Its residue has an isolated
$2r_0^2$ term. First suppose the leading source weights are finite.
The remaining residue is
\[
 q_s^TPq_s+2V_0b_*^Tq_s+c_2V_0^2+4H R_0^2+c_0K_0^2,
 \quad q_s=(X,Z)^T,                                  \tag{UA7}
\]
where $D_u=1+u^2$, $H=H_0$, $A_w=A_0^{\rm weak}$ are (RC3), and
\[
\begin{gathered}
 P=\begin{pmatrix}
 4nD_u+4m+2&(4m+2)u\\
 (4m+2)u&(4m+2)u^2+4\alpha D_u
 \end{pmatrix},\quad \ell=(u,-1)^T,\\
 b_*=\binom{-4mD_u}{-D_uu(4m+4\alpha+2)},\quad
 c_0=4A_wD_uu^2,\quad
 c_2=4mD_u^2+(4\alpha+2)D_uu^2.                      
\end{gathered}
\tag{UA8}
\]
The matrix $P$ is strictly positive because $m,\alpha$ are strong.
The last sign-changing part of (UA7) is handled by adding
$2\lambda$ times (UA5). This gives precisely the two blocks (UA3).

Here is a direct proof of (UA1) for every ordinary angle away from
contact. Put $N_*=2n+1$, $T_*=2m+2\alpha+1$, $U=u^2$, and
$T=V^2$. The leading state polynomial (RC4), with its testing $a$
replaced by $u$ but its source unchanged, has determinant
\[
 \det f=8D_u^2u^2V^2\,\mathcal L(U,T),                \tag{UA9}
\]
where
\begin{align*}
\mathcal L={}&\frac{4\alpha mN_*}{U}
 +4A_wHT_*U+\frac{4A_w\alpha(2m+N_*)}{T}
 +\frac{4A_wnT_*U}{T}\\
 &+2HN_*T_*T+\frac{4HmN_*T}{U}
 +4A_wH(2m+N_*)-N_*T_*.
\end{align*}
After division by $N_*T_*$ all possible opening poles cancel. Define
\[
 i_N=N_*^{-1},\quad t_m=m/T_*,\quad t_{am}=\alpha m/T_*,
 \quad q_\alpha=\frac{\alpha(2m+N_*)}{N_*T_*},\quad i_\alpha=1/\alpha.
\]
Writing $z=a^2$, $A_c=1+zr$, and using the leading $\Delta$ of (RC3),
these quantities are exactly
\[
\begin{aligned}
 i_N&=\frac{r\Delta}{(b+r)A_c^2},&
 t_m&=\frac{(1-r)(1-zb)}{2b(b+r)(z+1)},\\
 t_{am}&=\frac{z^2(b+r)(1-r)}{4\Delta},&
 q_\alpha&=\frac{bz^2\{(z+1)b+(z+2)r-1\}}{2A_c^2},\\
 i_\alpha&=\frac{2(1-zb)\Delta}{z^2(z+1)b(b+r)^2}.
\end{aligned}\tag{UA10}
\]
All denominators here are units. The normalized Laurent polynomial is
\[
 \frac{\mathcal L}{N_*T_*}
 =\frac{4t_{am}}U+4A_wHi_NU+\frac{4A_wq_\alpha}T
  +2A_w(1-i_N)\frac UT+2HT+4Ht_m\frac TU
  +4A_wHq_\alpha i_\alpha-1.                         \tag{UA11}
\]
Its value and its two first derivatives vanish at $(U,T)=(z,1)$,
as direct substitution verifies. In logarithmic coordinates the
nonconstant terms have nonnegative coefficients. The coefficients of
$U^{-1}$ and $T$ are positive units. If $i_N>0$, the coefficients of
$U$ and $T^{-1}$ are positive; if $i_N=0$, that of $U/T$ is positive.
The function is therefore strictly convex in logarithmic coordinates
and has its unique minimum at $(z,1)$. This proves strict positivity
at fixed $U\ne z$ for every $T>0$. The constant and leading
coefficients of the equivalent quadratic (UA1) are strictly positive:
$c_0>0$, and the positive $T$ coefficient in (UA11) gives the second
assertion. Thus (UA2)--(UA3) prove the required strict residue Gram.

\paragraph{Opening boundaries and divergent source weights.}
The preceding proof extends to all allowed opening residues, rather
than assuming $n,m,\alpha$ bounded. The following formulas make the
needed uniformity checkable. Put
\[
 J_0=q_\alpha+\frac{u^2(1-i_N)}2.
\]
Then $J_0$ is a positive unit, and direct two-by-two inversion gives
\[
\begin{aligned}
 \frac{\det P}{N_*T_*}&=8D_uJ_0,\\
 d&=\frac{u^2i_N+q_\alpha i_\alpha}{4J_0},&
 e&=\frac{uD_u}{2J_0},&
 h&=\frac{D_u(4t_{am}-u^2)}{J_0},\\
 (P^{-1}b_*)_1&=\frac{i_N(u^2-4t_{am})}{2J_0},&
 (P^{-1}b_*)_2&=-\frac{u\{4t_{am}i_N+(1-i_N)u^2+1\}}{2J_0}.
\end{aligned}\tag{UA12}
\]
To check the unit assertion, if $i_N=0$ then $J_0=u^2/2>0$.
If $i_N>0$, the weight $N_*$ is finite and the strong lower bounds
on $m,\alpha$ give $q_\alpha>0$: this follows directly from
$\alpha(2m+N_*)/[N_*(2m+2\alpha+1)]$, whether both weights are
finite or one diverges. Thus (UA12) has only unit denominators on
the entire source residue domain.

The strictly positive leading coefficient is also explicit:
\[
 dh+e^2=\frac{D_u(u^2+2t_m)}{2J_0}>0.
\]
Equations (UA10)--(UA12) prove continuity of $d,e,h$ and of the
minimizing hard vector. The convexity proof for (UA11) remains strict
at every allowed boundary. Choose $\lambda$ from the strict interval
(UA2) at that boundary; the same strict inequalities hold nearby.
Moreover $P$ is bounded below by a positive real matrix, while
$P^{-1}b_*$ and $P^{-1}\ell$ are bounded. Its Schur factorization
therefore bounds the first block in (UA3) below by a positive real
multiple of the identity, uniformly in the nearby source. The second
block is a strictly positive real matrix as well.

The finite leading sources used in this limiting argument are valid
real interior sources. This point must be established separately:
the actual source has $b=h+(k^2/z)t$, and it is not legitimate to
substitute that actual $b$ into a leading formula with denominator
$1-zb$ when this denominator has an infinitesimal sign change.
For the limiting real parameters, put
\[
 f(b,r)=1+r(z-1)-brz(z+1)-r^2z(z+2).
\]
Their domain is $b>0$, $0\le r<1$, $g=(1+z)b-1+r\ge0$,
$1-zb\ge0$, and $f(b,r)\ge0$. For $0<s<1$, let
\[
 b_s=sb+(1-s)/(1+z),\qquad r_s=sr.
\]
Direct expansion gives
\[
 g_s=sg,\qquad 1-zb_s>0,\qquad
 f(b_s,r_s)=(1-s)(1+s-sr)+s^2f(b,r)>0.
\]
Increase $b_s$ by a sufficiently small $\epsilon>0$. This makes
$g_s$ strict, decreases the last two strict margins by respectively
$z\epsilon$ and $r_sz(z+1)\epsilon$, and therefore preserves them.
If $r=0$, then increase $r_s$ by a sufficiently small positive
amount; $g$ increases and $f$ remains positive by continuity.
Letting these perturbations tend to zero as $s\to1$ supplies valid
finite real leading sources approaching every allowed residue.
Their $t$ may be chosen as $(1+zr)/2$, since the leading harmonic
weight is independent of $t$. The normalized leading identity
$2rn\Delta=bf(b,r)$ verifies the nonnegative-$n$ requirement.
The same construction supplies the real-interior approximation
used in the finite-cap step of the local UI certificate.

At most two of $n,m,\alpha$ diverge: $m$ can diverge only at $g=0$,
and $\alpha$ only at $1-zb=0$, which cannot occur together. Their
centered excess rows have hard-state parts proportional to
\[
 X,\qquad X+uZ,\qquad Z.
\]
Any possible pair is independent because $u$ is a unit. For a vector
in their common limiting kernel, change only $(X,Z)$ by a vanishing
amount to annihilate the corresponding nearby exact rows. Their
large penalties then vanish exactly. Taking the limit in the uniform
strict lower bound proves that the bounded finite part is strictly
positive on that limiting kernel. The elementary finite-dimensional
penalty lemma used in the local theorem now supplies finite real caps
$n_*,m_*,\alpha_*$. Their centered excesses are precisely the retained
packets (UR14). The exact bounded remainder with those caps, after
adding (UA5), has a strictly positive residue Gram on (UA4).
This proves the finite ordinary-angle complement for all source
openings and literal $n=0$.

\paragraph{Endpoint frames.}
Choose positive real $\mu,m_*,\alpha_*$ below, respectively,
$k^2C,k^2D$, $m$, and $\alpha$. The products $k^2C,k^2D$ are strong
by (UR5), and $A_w=A/k^2$ is a positive unit. Retain $n$ and all
positive excesses in full. The following formulas concern the parent
with $C=D=\mu/k^2$, $m=m_*$, $\alpha=\alpha_*$ and $A=A_wk^2$.
They therefore suffice. The root coordinates need not themselves be
bounded; only their displayed Gram coefficients must be bounded.

At $u,v$ infinitesimal use
\[
 E=(uW-Y)/k,\quad R=v(uX-Z)/k,\quad V=vW,\quad K=kuW,
 \quad X,Z.
\]
The exact amplitudes are
\[
 J=k(E+R),\quad M=D_uV-X-uZ-uvkE,\quad
 vY-Z=uV-Z-vkE,\quad kY=K-k^2E.
\]
Their full residue form is
\[
 4\mu\{E^2+(E+R)^2\}+4m_*(V-X)^2+4\alpha_*Z^2
       +4\overline{A_w}K^2+2X^2,                    \tag{UA13}
\]
which is strict. Literal zero testing angles are included without
any division by a testing angle.

For bounded $u$ and finite unit $v$, put
$E=(uW-Y)/k$, $F=J/k$. Then
\[
 Y=uW-kE,\quad Z=uX+k(E-F)/v.
\]
The roots $(E,F,W,X)$ have strict residue
\[
 4\mu D_vE^2+4\mu F^2+
 \{4m_*D_u^2+(4\alpha_*+2)D_uu^2\}(vW-X)^2
 +2D_vD_uX^2.                                      \tag{UA14}
\]
For bounded $u$ and $v=\infty$ write $q=1/v$ and homogenize by $q^2$.
The same roots are defined using the homogeneous $J$, and
\[
 Y=uW-kE,\qquad Z=uX+kqE-kF.
\]
The strict residue is (UA14) with $D_v=1$ and $(vW-X)^2$ replaced
by $W^2$. This includes the literal endpoint $q=0$.

For $u=\infty$, put $p=1/u$, $D_p=1+p^2$, and homogenize by $p^2$.
Write $r_p=W-pY$, $Q_p=X-pZ$, and $P=vY-Z$. Then
\[
 J=r_p+vQ_p,\qquad M=D_pP-pQ_p+pv r_p.
\]
For $v$ infinitesimal use the six roots
\[
 E=r_p/k,\quad F=vQ_p/k,\quad P,\quad R=pQ_p,
 \quad K=kY,\quad S=pZ.
\]
The signed core is exactly
\[
 2k^2(E+F)^2+2D_pP^2+2D_vR^2-4D_pD_vRP
    +4D_pD_vpFK+2D_pD_vS^2.
\]
Its cross-term reduction uses $RvY=pFK$, an exact identity.
Adding every capped packet gives the strict residue
\[
 4\mu\{E^2+(E+F)^2\}+(4m_*+2)(P-R)^2
 +4\alpha_*P^2+4\overline{A_w}K^2+2S^2.             \tag{UA15}
\]
For finite unit $v$, instead use $E=r_p/k$, $F=J/k$, $P,K,S$.
Since $Q_p=k(F-E)/v$, the strict residue is
\[
 4\mu D_vE^2+4\mu F^2+(4m_*+4\alpha_*+2)P^2
 +4\overline{A_w}D_vK^2+2D_vS^2.                   \tag{UA16}
\]
Finally, at both infinite endpoints put $q=1/v$ and homogenize by
$p^2q^2$. Use
\[
 E=(W-pY)/k,\quad F=(qW+X-pqY-pZ)/k,\quad
 P=Y-qZ,\quad K=kY,\quad S=pZ.
\]
Here $M=D_pP+pkD_qE-pqkF$ and the entire signed core is
\[
 2k^2F^2+2D_pP^2+2D_qp^2k^2(F-qE)^2
 +4D_qD_pk(F-qE)S+2D_qD_pS^2.
\]
The strict residue is (UA16) with $D_v=1$.

\paragraph{Exhaustion and restoration.}
A testing angle is infinitesimal, a bounded unit, or projectively
infinite. Equations (UA13)--(UA16) cover every pair except ordinary
$u$ with infinitesimal $v$. The local theorem covers the equal-residue
contact, and (UA4)--(UA12) cover its complement. Exact sign gauges are
\[
 (u,W,X)\mapsto(-u,-W,-X),\qquad
 (v,W,Y)\mapsto(-v,-W,-Y);
\]
all unspecified variables are fixed. Both preserve the full parent,
so negative ordinary contact residues are included. All endpoint
frames have been written after homogenization, including literal
projective endpoints. Each proof has an exact bounded Gram with
strict real residue; the strict-residue Gram lemma gives tensor SOS.
Restoring all retained packets completes the original parent.
Real and imaginary parts treat complex states. This closes the
entire strong C2 UI admission and its exact symmetry image HU.



\section{\texorpdfstring{The exact strong-C2 ordinary--small source reduction}{The exact strong-C2 ordinary--small source reduction}}
\label{component:pc3daa35e9a9b}

\label{pc3daa35e9a9b:sec:ui-full-parent-reduction}

This section records the original source, every parent packet, and the
bounded products used in the following reciprocal and local certificates.
It is an exact algebraic reduction, not an all-tests admission. The
coefficient and testing fields remain independent real closed extensions
of $\mathbb R$.

\paragraph{Source definitions and physical domain.}
Take the actual source angles $u_c=a$ and $v_c=k$, where $a>0$ is a
bounded unit and $k>0$ is infinitesimal. Set $z=a^2$ and
$\varepsilon=k^2/z$. In the original canonical C2 chart put
\[
 h=b-\varepsilon t,\quad A_c=1+zr,\quad j=A_c-t,\quad
 g=(1+z)b-1+r-\varepsilon zt,\quad S=A_c+\varepsilon zt.
 \tag{UR1}\label{pc3daa35e9a9b:eq:ui-source-gaps}
\]
The native positive gaps are $r,1-r,t,h,g,j,1-zh$; assume all original
event signs, including $n\ge0$, and assume $C,\alpha$ are bounded below
by positive real numbers. No positive lower bound on a native opening
is imposed. Define the two complete source polynomials
\begin{align*}
 \Delta={}&1+(1+z)b+z(1+z)b^2+2z(1+z)br+z^2r^2+2zr\\
 &+\varepsilon z\{(1+z)b+1+zr(2-r)\},\\
 N={}&(1+z)(b+r)\{b+\varepsilon(1-t)\}
          -\varepsilon(1+\varepsilon z)(1-r)t.
 \tag{UR2}\label{pc3daa35e9a9b:eq:ui-source-polynomials}
\end{align*}
With $d=1$ and $\lambda=\beta=0$, the unchanged original coefficients
are
\[
 \begin{aligned}
 n&=\frac{(b+r)A_cS-r\Delta}{2r\Delta},&
 m&=\frac{(1+k^2)(1-r)A_cS}{2g\Delta},\\
 C&=\frac{A_cN}{2\varepsilon t\Delta},&
 A&=\frac{\varepsilon A_c(b+r)S}{2h\Delta},\\
 D&=\frac{(1+z)(b+r)(b-k^2r)S}{2\varepsilon j\Delta},&
 \alpha&=\frac{z^2(b+r)N}{2(1-zh)\Delta}.
 \end{aligned}                                                   \tag{UR3}
 \label{pc3daa35e9a9b:eq:ui-original-source}
\]
These are the original rational C2 source formulas, with its angle-square
parameter replaced by $z$ and its squared angle-ratio parameter by
$\varepsilon$; all quantities formerly supplied by a source-file
reference are displayed in (UR1)--(UR3).

The source bounds used below follow directly from these definitions.
First $t<A_c$, $h<1/z$, and $b=h+\varepsilon t$ give bounded $t,h,b$.
The inequality $(1+z)b+r>1$ gives a positive lower bound on $b+r$;
(UR2) makes $\Delta$ a positive bounded unit. The exact positive identity
\[
 N=(1+z)(b+r)h+
       \varepsilon\{g+j+rt(1+\varepsilon z)\}                 \tag{UR4}
\]
shows that infinitesimal $h$ would make $b,N$ infinitesimal and
$1-zh$ a unit, contradicting strong $\alpha$. Thus $h,b,N$ are units.
The numerator of $2rn\Delta$ evaluated at $r=1$ is exactly
\[
 -z(\varepsilon+h)(b+1)(z+1).
\]
It has negative unit residue, so $n\ge0$ forces $1-r$ to be a unit.
Consequently
\[
 mg,\quad A/\varepsilon,\quad \varepsilon Ct,\quad
 \varepsilon Dj,\quad \alpha(1-zh)
 \quad\hbox{are positive bounded units},\qquad nr\text{ is bounded}.
 \tag{UR5}\label{pc3daa35e9a9b:eq:ui-source-bounds}
\]
In particular $m$ is strong, $A$ is infinitesimal, and $C,D$ may have
arbitrarily large positive sizes. Literal $n=0$ is retained.

\paragraph{The entire parent and exact contact.}
For independent testing angles and a homogeneous state define
\[
 \begin{gathered}
 D_u=1+u^2,\quad D_v=1+v^2,\quad
 r_{\rm amp}=uW-Y,\quad Q=uX-Z,\\
 J=r_{\rm amp}+vQ=uW+uvX-Y-vZ,\\
 M=vW-X+uvY-uZ,\quad p=vY-Z,\quad B=X+uZ.
 \end{gathered}
\]
The signed core and the full original parent are
\[
 \mathcal K=2\{J^2+D_up^2+D_vB^2-D_uD_vZ^2\},
\]
\[
 F=4nD_uD_vX^2+4mM^2+4CD_vr_{\rm amp}^2
       +4AD_uD_vY^2+4DJ^2+4\alpha D_up^2+\mathcal K.
 \tag{UR6}\label{pc3daa35e9a9b:eq:ui-full-parent}
\]
Normalize the actual contact by $W_c=1$. Introduce the bounded pure-source
quantities
\[
 x=-r/b,\quad \rho=t/(ab),\quad \sigma=-j/(ab),\quad
 s=\sigma-\rho=-A_c/(ab),\quad
 \mu=g/b,\quad \pi=(zh-1)/(ab).                       \tag{UR7}
 \label{pc3daa35e9a9b:eq:ui-contact-amplitudes}
\]
Here $s$ is a nonzero bounded unit. The exact stationary contact is
\[
 (u_c,v_c;W_c,X_c,Y_c,Z_c)
 =\left(a,k;1,kx,a-\rho k^2,\frac{k}{ab}\right),       \tag{UR8}
\]
and its amplitudes are
\[
 r_{{\rm amp},c}=\rho k^2,\quad J_c=\sigma k^2,\quad
 M_c=\mu k,\quad p_c=\pi k,\quad Q_c=sk.
\]
Substitution of (UR3) gives $F=0$ and all four state and both angle
first derivatives zero at (UR8). These are the unchanged native
contact identities; no equality of independent testing rates is used.

\paragraph{Products with all possible opening poles cancelled.}
Define two positive bounded units
\[
 c_1=\frac{zA_cN}{2\Delta},\qquad
 c_2=\frac{z(1+z)(b+r)(b-k^2r)S}{2\Delta}.
\]
The exact products required by the local cap argument are
\[
 \begin{aligned}
 nx&=-\frac{(b+r)A_cS-r\Delta}{2b\Delta},&
 m\mu&=\frac{(1+k^2)(1-r)A_cS}{2b\Delta},\\
 \alpha\pi&=-\frac{z^2(b+r)N}{2ab\Delta},&
 C\rho k^2&=\frac{c_1}{ab},\qquad
 D\sigma k^2=-\frac{c_2}{ab}.
 \end{aligned}                                                   \tag{UR9}
 \label{pc3daa35e9a9b:eq:ui-cancelled-products}
\]
Their products with the bounded factors $x,\mu,\pi$ also show
$nx^2,m\mu^2,\alpha\pi^2$ bounded. Crucially, cancellation before
estimating gives the further exact identity
\[
 Ct-Dj=
 \frac{(1+z)(b+r)\{A_cj-zt(b-\varepsilon zr)\}
       -A_c(1+\varepsilon z)(1-r)t}{2\Delta}.          \tag{UR10}
 \label{pc3daa35e9a9b:eq:ui-cancelled-CD-difference}
\]
Thus $C\rho+D\sigma=(Ct-Dj)/(ab)$ is bounded for every positive rate
of $t,j$, including infinitesimal openings.

Set
\[
 L=C+D,\qquad q=D/L,\qquad H=CD/L,\qquad
 \tau=(C\rho+D\sigma)/L,\qquad \gamma=k^2H.
\]
Here $q,1/L,\tau,L\tau$ are bounded, and
\[
 \gamma=\frac{c_1c_2}{c_1j+c_2t}\asymp1,
 \qquad t+j=A_c\asymp1.                              \tag{UR11}
 \label{pc3daa35e9a9b:eq:ui-harmonic-unit}
\]
The denominator in (UR11) is a positive unit even if either opening
is infinitesimal. At finite source residues it gives precisely $H_0$
in (RC3). Equations (UR9)--(UR11) also show that the bounded finite
part extends continuously to the allowed boundary residues: no
inverse of $r,g,1-zh,t,j$ remains in these products.

\begin{proposition}[Exact bounded UI remainder]
\label{pc3daa35e9a9b:prop:ui-bounded-remainder}
For bounded source caps $0\le n_*\le n$ and positive bounded-unit
caps $m_*\le m$, $\alpha_*\le\alpha$, the parent (UR6) is the sum
of explicit positive centered packets and a bounded polynomial
remainder $G$. Every centered packet vanishes with its full first jet
at (UR8), so $G$ has exact zero value and all first derivatives there.
\end{proposition}
\begin{proof}
The exact original $C,D$ combination is
\[
 CD_vr_{\rm amp}^2+DJ^2
   =L(r_{\rm amp}+qvQ)^2+Hv^2Q^2+Cv^2r_{\rm amp}^2.
\]
Retain the three nonnegative packets
\[
 \begin{aligned}
 P_L&=4L(r_{\rm amp}+qvQ-\tau k^2W)^2,\\
 P_H&=4H(vQ-sk^2W)^2,\\
 P_C&=4Cv^2(r_{\rm amp}-\rho k^2W)^2.
 \end{aligned}                                                   \tag{UR12}
\]
Their subtraction leaves exactly
\begin{align*}
 G_{CD}={}&8L\tau k^2W(r_{\rm amp}+qvQ)-4L\tau^2k^4W^2\\
 &+8Hs k^2vWQ-4Hs^2k^4W^2\\
 &+8C\rho k^2v^2Wr_{\rm amp}-4C\rho^2k^4v^2W^2.
 \tag{UR13}
\end{align*}
All its coefficients are bounded by (UR9)--(UR11), including
$L\tau^2=(L\tau)^2/L$. Centering at $k^2W$ is essential to this
boundedness assertion.

Also retain
\[
 \begin{aligned}
 P_n&=4(n-n_*)D_uD_v(X-xvW)^2,\\
 P_m&=4(m-m_*)(M-\mu vW)^2,\\
 P_\alpha&=4(\alpha-\alpha_*)D_u(p-\pi vW)^2.
 \end{aligned}                                                   \tag{UR14}
\]
The exact corresponding remainders are
\[
 \begin{aligned}
 G_n={}&4n_*D_uD_vX^2+8(n-n_*)xvWD_uD_vX
                    -4(n-n_*)x^2v^2W^2D_uD_v,\\
 G_m={}&4m_*M^2+8(m-m_*)\mu vWM
                    -4(m-m_*)\mu^2v^2W^2,\\
 G_\alpha={}&4\alpha_*D_up^2+8(\alpha-\alpha_*)\pi vWD_up
                    -4(\alpha-\alpha_*)\pi^2v^2W^2D_u.
 \end{aligned}                                                   \tag{UR15}
\]
These coefficients are bounded by (UR9) and the bounded caps. With
\[
 G=G_{CD}+G_n+G_m+G_\alpha+4AD_uD_vY^2+\mathcal K,
\]
the entire exact identity is
\[
 F=P_L+P_H+P_C+P_n+P_m+P_\alpha+G.                 \tag{UR16}
 \label{pc3daa35e9a9b:eq:ui-entire-reduction}
\]
It includes the signed core and every original state-normal term.
All six retained roots vanish at (UR8), proving the contact-jet
assertion. This proposition does not claim positivity of $G$.
\end{proof}

For the local theorem retain $P_L,P_C,P_n,P_m,P_\alpha$ outside the
bounded Gram and combine $G$ with
\[
 P_H=4\gamma R_h^2,\qquad R_h=(vQ-sk^2W)/k.
 \tag{UR17}
\]
The factor $1/k$ is pure coefficient. The exact graph normals are
\[
 E_X=X-xvW,\quad E_Z=Z-(ux-s)vW,\quad
 E_r=r_{\rm amp}-\rho k^2W.
\]
Their inverse is polynomial:
\[
 X=E_X+xvW,\quad Z=E_Z+(ux-s)vW,\quad
 Y=uW-\rho k^2W-E_r,\quad Q=svW+uE_X-E_Z.
\]
Thus no state is discarded, including $W=0$. The following reciprocal
lemma identifies the leading relative-angle completion; the local
certificate supplies the complete normal lift, cap choice, and strict
Gram. The outer-angle and other commuting-axis join remains a separate
obligation, so this reduction adds no all-tests source admission.


\section{\texorpdfstring{A local strong-C2 certificate at the ordinary--small contact}{A local strong-C2 certificate at the ordinary--small contact}}
\label{component:p76545700d90b}

\label{p76545700d90b:sec:local-ui-certificate}

This section uses the exact full-parent reduction and source notation
in Section~\ref{pc3daa35e9a9b:sec:ui-full-parent-reduction}. It proves a local independent-test
statement. The other commuting axis and the outer projective join
are not included.

\begin{theorem}[Local ordinary--small C2 certificate]
Use the unchanged physical C2 source with $d=1$, actual contact
$u_c=a$ a positive bounded unit and $v_c=k$ a positive infinitesimal,
and $C,\alpha$ bounded below by positive real numbers. Then the full
parent is tensor SOS for every independent test with $v$ infinitesimal
and $u-a$ infinitesimal, and every projective state. The original
$n=0$ boundary and all positive source opening rates are included.
No upper bound on $n,m,C,D,\alpha$ is imposed.
\end{theorem}

Here $u-a$ is the difference of the corresponding tensor factors.
The residue condition means that their real residues agree; it does
not identify the source and testing fields or compare their rates.
The exact sign gauges give the corresponding negative-angle charts.

\paragraph{Bounded full remainder and its exact graph.}
The packet reduction leaves a bounded polynomial remainder $G$ and
the positive relative square
\[
 4\gamma R_h^2,\qquad
 \gamma=k^2CD/(C+D)\asymp1,\qquad
 R_h=\frac{v(uX-Z)-s k^2W}{k}.
\]
Every other retained packet is an exact positive square vanishing
at the full actual contact. For positive bounded caps $m_*,\alpha_*$
and $0\le n_*\le n$, set
\[
 E_X=X-xvW,\quad E_Z=Z-(ux-s)vW,\quad
 E_r=uW-Y-\rho k^2W,
 \qquad Q_N=uE_X-E_Z.
\]
The full inverse is polynomial. All terms of $G$, including its
complete first-normal rows and quadratic normal block, are retained.
There are no higher state-normal terms.

Put $w=u-a$, $\xi=v^2$, and use the additional separated roots
\[
 S=vwW,\qquad T=kwW,\qquad R=(v^2-k^2)W/k.
\]
The exact identities
\[
 kR=\frac{kR_h-vQ_N}{s},\qquad
 RT=\frac{R_hT-SQ_N}{s},\qquad
 S^2-T^2=\frac w s(R_hT-SQ_N)                      \tag{UI1}
\]
use only the bounded pure-source inverse $1/s$.

\paragraph{Complete scalar and normal lifting.}
On the graph $E_X=E_Z=E_r=0$, write $G=\phi(u,\xi)W^2$.
Its bounded polynomial coefficients and the original parity imply
$\phi\in(k^2,\xi)$. The exact contact identities give
$\phi=\phi_u=\phi_\xi=0$ at $(u,\xi)=(a,k^2)$.
Successive monic Taylor division gives
\[
 \phi=w^2(k^2A_0+\xi A_1)
             +w(\xi-k^2)B+(\xi-k^2)^2C_0,          \tag{UI2}
\]
with bounded coefficients. To see the asserted first factor directly,
divide first in $u-a$; the second $u$ coquotient stays in
$(k^2,\xi)$. The constant and first $u$ terms are then divided in
$\xi-k^2$. No source infinitesimal is inverted by these operations.
Equations (UI1) allocate the three terms of (UI2) on
$(S,T,R_h,E_X,E_Z)$.

The full first $E_X,E_Z$ coefficients are $vL_X(u,\xi)$ and
$vL_Z(u,\xi)$. Exact contact gives $L_i(a,k^2)=0$, so they are
\[
 v\{wL_{i0}+(\xi-k^2)L_{i1}\}.
\]
The first summand is allocated on $S$, and the second uses
$vkR=v(kR_h-vQ_N)/s$. The entire $E_r$ coefficient belongs to
$(k^2,\xi)$ and vanishes at contact. The same division writes it as
\[
 w(k^2L_{r0}+\xi L_{r1})+(\xi-k^2)L_{r2}.
\]
It consequently uses $kT,vS,kR$. All these $E_r$ couplings have
zero residue. Thus (UI1)--(UI2) produce an exact bounded Gram on
\[
             (E_r,E_X,E_Z,R_h,S,T).                 \tag{UI3}
\]
This is a full-parent lift, not a restriction to a state kernel.

\paragraph{The strict matched Hessian.}
First take finite positive residue source weights. The leading model
before setting the ordinary test angle equal to $a$ is (RC4), with
every testing occurrence of $a$ replaced by $u$; its source
coefficients remain those of (RC3). Let $\mathsf G$ denote one half
of its angular Hessian after eliminating its two state directions,
in variables $(V,u)$ at the exact contact. Direct differentiation and
two-by-two Schur elimination give
\[
 \mathsf G_{uu}=
 \frac{8(1-r)(1+zr)^2}{zb\Delta}>0,                 \tag{UI4}
\]
\[
 \det\mathsf G=
 \frac{16(b+r)^2(1+zr)^3
                \{3z+4-(z+2)^2r\}}{b^2\Delta^2}.  \tag{UI5}
\]
The other two entries, making boundedness explicit, are
\begin{align*}
 \mathsf G_{Vu}&=
 \frac{4(1+zr)^2\{zb+zr+2r-2\}}{ab\Delta},\\
 \mathsf G_{VV}&=
 \frac{8(1+zr)\{z^2b^2+z^2br+zb^2+2zbr-zb-r+1\}}
      {b\Delta}.
\end{align*}
For a short check of strictness, the exact leading $n$ numerator is
$bQ_n$, where
\[
 Q_n=1+r(z-1)-z(z+2)r^2-z(z+1)br
     =(1-(z+1)r)(1+zr)-zr g.
\]
Since $Q_n\ge0$ and $g\ge0$, one has $r\le1/(z+1)$.
Consequently
\[
 3z+4-(z+2)^2r\ \ge\ \frac{z(2z+3)}{z+1}>0.       \tag{UI6}
\]
All denominators in (UI4)--(UI6) are positive bounded units.
These formulas therefore give a strict, uniformly bounded-away-from-zero
angular Schur form on the source residue domain, including $n=0$.
They contain no inverse of $r$, $g$, or $1-zb$.

\paragraph{Large weights and finite real caps.}
We use the elementary finite-dimensional penalty lemma: if a real
symmetric form $B$ is strictly positive on the common kernel of the
rows of a matrix $V$, then $B+M V^TV$ is positive definite for some
finite real $M>0$. A proof is to split into $\ker V$ and its
orthogonal complement. The first diagonal block is strict; taking
$M$ sufficiently large makes its Schur complement strict too.

The bounded finite part obtained by the centered packets is
continuous at all source boundary residues. This assertion is
algebraic: its only large-weight products are
\[
 nx,\ nx^2,\ m\mu,\ m\mu^2,\ \alpha\pi,\ \alpha\pi^2,
 \quad C\rho k^2,\ D\sigma k^2,\ k^2H,
\]
and their bounded cap analogues. Substitution of the original source
formulas cancels the possible $r,g,1-zh,t,j$ poles in those products;
their denominators are the displayed positive units. For example,
\[
 m\mu=\frac{(1+k^2)(1-r)A_cS}{2b\Delta},\qquad
 \alpha\pi=-\frac{z^2(b+r)N}{2ab\Delta}.
\]
The other cancellations are
\[
 nx=-\frac{(b+r)A_cS-r\Delta}{2b\Delta},\quad
 nx^2=x(nx),\quad m\mu^2=\mu(m\mu),\quad
 \alpha\pi^2=\pi(\alpha\pi),
\]
\[
 C\rho k^2=\frac{c_1}{ab},\qquad
 D\sigma k^2=-\frac{c_2}{ab},\qquad
 c_1=\frac{zA_cN}{2\Delta},\quad
 c_2=\frac{z(1+z)(b+r)(b-k^2r)S}{2\Delta}.
\]
Finally $C\rho+D\sigma=(Ct-Dj)/(ab)$, using the exact canceled
formula in the full-parent reduction. Thus even this difference has
no hidden small-opening denominator.
For $k^2H$, write $k^2C=c_1/t$, $k^2D=c_2/j$, with
$c_1,c_2$ positive bounded units. Then
$k^2H=c_1c_2/(c_1j+c_2t)$, whose denominator is a unit because
$t+j=A_c$ is a unit. Its residue is exactly $H_0$ in (RC3).

At most two of $n,m,\alpha$ can diverge. Indeed $m$ can diverge
only at $g=0$, while $\alpha$ can diverge only at $1-zb=0$;
these two equalities would give $g=1/z+r>0$. The hard-state parts
of the three possible diverging centered rows are, up to nonzero
unit factors,
\[
                   E_X,\qquad E_X+aE_Z,\qquad E_Z.
\]
Every possible pair is independent with a unit determinant.

Here is why their constrained limiting finite form is strict.
Approach the boundary by the positive leading-source domain. The
state block is bounded below by a positive real form because $m$ and
$\alpha$ are strong. Its minimizing hard vector for a fixed angular
vector is bounded: solve the at most two diverging packet equations
in the hard variables, using the unit determinants just displayed;
this supplies a bounded competitor, and all remaining coefficients
stay bounded. If $h_0$ is that competitor and $h_{\min}$ the state
minimizer for angular vector $q$, the exact identity
\[
 (h_0-h_{\min})^TP(h_0-h_{\min})
       =F(h_0,q)-q^T\mathsf Gq
\]
and $P\succeq cI$ make boundedness of $h_{\min}$ explicit.
Together with (UI4)--(UI6), the state Schur factorization
then bounds the full matched Hessian below by a positive real multiple
of the Euclidean norm. For a vector in the common limiting packet
kernel, adjust only its hard coordinates to lie in the exact nearby
packet kernels. The adjustment tends to zero and annihilates the
large weights exactly. Taking the limit proves strict positivity of
the bounded finite part on that common kernel. The penalty lemma now
chooses finite real caps for every diverging weight. Bounded source
weights can be retained in full. All caps are below their original
weights, so the remaining centered excesses are positive squares.

This limiting argument uses only real finite-dimensional forms to
choose caps. No inverse of a tensor matrix and no mixed root are
used in the full-parent identity. The boundary residues are limits
of the positive leading-source domain: $g$ and $1-zb$ specify the
interval for $b$, and decreasing a positive $r$ if necessary makes
$Q_n\ge0$ strict. Thus the argument also covers a literal original
$n=0$ and all combinations of the allowed boundary residues.

\paragraph{Separate angular roots and full strictness.}
The residue of the hard block in (UI3) is
\[
 2E_r^2+4n_*(1+a^2)E_X^2
 +(4m_*+2)(E_X+aE_Z)^2+4\alpha_*(1+a^2)E_Z^2
 +4\gamma R_h^2,
\]
which is strict. In the residue Gram, $S$ couples only to the two
ordinary hard normals, and $T$ only to $R_h$; the $E_r$ couplings
vanish. Setting $S=T$ gives the matched leading Hessian under an
invertible bounded change of coordinates. Indeed at $V=1$ the
ordinary normals are its two state deviations, $R_h$ supplies the
relative-angle deviation with nonzero coefficient $2s$, and
$S=T$ supplies the ordinary-angle deviation. The chosen caps therefore
make the sum of the two angular Schur diagonals strictly positive.

Add an appropriate real multiple of the last zero identity in (UI1).
It redistributes mass between $S^2$ and $T^2$; all accompanying
coefficients contain $w$ and have zero residue. Since their Schur
diagonal sum is positive, choose this multiple so that both Schur
diagonals are positive. The entire six-root residue Gram is now
positive definite. The strict-residue Gram lemma proves the exact
bounded remainder plus relative packet is tensor SOS.

Finally restore every centered excess and both other C/D packets.
The construction was homogeneous in $W$ and its state inverse was
polynomial, so it includes $W=0$ and every projective state. Real and
imaginary parts give complex states. This completes only the stated
local independent-test domain; the outer-angle and other commuting-axis
certificates are still required before admitting the full UI source leaf.


\section{\texorpdfstring{A reciprocal contact completion for a singular angle}{A reciprocal contact completion for a singular angle}}
\label{component:pa142a027a282}

\label{pa142a027a282:sec:ui-reciprocal-contact}

The following exact algebraic lemma isolates a useful mechanism. It
does not assert a complete C2 all-tests theorem.

\begin{lemma}[Reciprocal contact completion]
Let $A_0,A_2$ be symmetric matrices over a real closed field, let
$q_0$ be a column vector, and consider
\[
 f(q,V,W)=q^T(A_0+V^2A_2)q+2VW b^Tq+(c_0+c_2V^2)W^2.
\]
Suppose $(q,V,W)=(q_0,1,1)$ is a zero with all first derivatives
zero. Then
\[
 b=-(A_0+A_2)q_0,\qquad
 c_0=q_0^TA_2q_0,\qquad c_2=q_0^TA_0q_0,
\]
and the entire polynomial has the exact decomposition
\[
 \boxed{f=(q-Vq_0W)^TA_0(q-Vq_0W)
             +(Vq-q_0W)^TA_2(Vq-q_0W).}             \tag{RC1}
\]
Consequently $A_0,A_2\succeq0$ imply a sum of squares. If
$A_0\succ0$ and $q_0^TA_2q_0>0$, its nonzero zero states have
$V=\pm1$ and $q=Vq_0W$.
\end{lemma}
\begin{proof}
The $q$ equation gives the displayed value of $b$. The $V$
equation then gives $c_2=q_0^TA_0q_0$, and the $W$ equation gives
$c_0=q_0^TA_2q_0$. Expanding the right side of (RC1) proves the
identity. The zero statement follows first from the strictly positive
$A_0$ term and then from $(V^2-1)^2W^2q_0^TA_2q_0$.
\end{proof}

\paragraph{Homogeneous form without a mixed inverse.}
For a positive pure-coefficient $k$, substitute $q=Q/k$ and
$V=v/k$, and multiply (RC1) by $k^2$. It becomes
\[
 (Q-vq_0W)^TA_0(Q-vq_0W)
 +\frac1{k^2}(vQ-k^2q_0W)^TA_2(vQ-k^2q_0W).       \tag{RC2}
\]
Thus its roots use only the pure-coefficient factor $1/k$.
There is no inverse or radical of a mixed tensor element and no
comparison between a testing rate and $k$.

\paragraph{The bounded-source C2 UI leading model.}
Here $u_c=a$ is a positive bounded unit and $v_c=k$ is
infinitesimal, with $d=1$. To state a residue model, first restrict
to the subdomain where the original $n,m,\alpha$ are bounded and
the source parameters $r,g,1-a^2b$ have nonzero residues. In the
following formulas $a,b,r$ denote those residues, $z=a^2$, and
\[
 A_c=1+zr,\quad g=(1+z)b-1+r,\quad
 \Delta=1+(1+z)b+z(1+z)b^2+2z(1+z)br+z^2r^2+2zr.
\]
The source formulas give
\[
\begin{gathered}
 n=\frac{(b+r)A_c^2-r\Delta}{2r\Delta},\qquad
 m=\frac{(1-r)A_c^2}{2g\Delta},\qquad
 \alpha=\frac{z^2(1+z)b(b+r)^2}{2(1-zb)\Delta},\\
 H_0=\overline{k^2CD/(C+D)}
       =\frac{z(1+z)(b+r)b}{2\Delta},\qquad
 A_0^{\rm weak}=\overline{A/k^2}
       =\frac{A_c^2(b+r)}{2zb\Delta}.
\end{gathered}                                                   \tag{RC3}
\]
Normalize the actual contact by $W=1$. With $v=kV$, $X=k\xi$,
$Z=k\zeta$ and the leading constraint $Y=aW$, the leading
polynomial of $F/k^2$ is
\begin{align*}
 f={}&4n(1+a^2)\xi^2
 +4m\{V(1+a^2)W-\xi-a\zeta\}^2
 +4H_0V^2(a\xi-\zeta)^2\\
 &+4A_0^{\rm weak}(1+a^2)a^2W^2
 +(4\alpha+2)(1+a^2)(VaW-\zeta)^2\\
 &+2(\xi+a\zeta)^2-2(1+a^2)\zeta^2.               \tag{RC4}
\end{align*}
It has the form of the lemma, with $q=(\xi,\zeta)^T$,
\[
 q_0=(-r/b,1/(ab))^T,
\]
\[
 q^TA_0q=4n(1+a^2)\xi^2+(4m+2)(\xi+a\zeta)^2
                      +4\alpha(1+a^2)\zeta^2,
 \qquad q^TA_2q=4H_0(a\xi-\zeta)^2.
\]
Here $A_0\succ0$ and $H_0>0$. Put
$s_0=aq_{0,1}-q_{0,2}=-(1+a^2r)/(ab)\ne0$.
The exact source identities (RC3) verify all three coefficient
relations in the lemma. Hence the apparently problematic relative
normal and scale difference combine into the single root
\[
             \frac{v(aX-Z)-s_0k^2W}{k}.             \tag{RC5}
\]
The ordinary centered normals $(X,Z)-vq_0W$ supply the remaining
positive directions. Requiring separate positive pivots for the
relative normal and the scale difference would obscure this
completion.

This proves an exact leading-model identity on its stated bounded
subdomain. It does not discard the additional large source scales:
those are handled algebraically by the centered-packet reduction in
Section~\ref{pc3daa35e9a9b:sec:ui-full-parent-reduction}. The complete local lifted
Gram is proved in Section~\ref{p76545700d90b:sec:local-ui-certificate}; the outer
projective join remains a separate obligation.


\section{\texorpdfstring{Five consecutive lower-bounded event packets}{Five consecutive lower-bounded event packets}}
\label{component:p80253fe9bc1b}

\label{p80253fe9bc1b:sec:five-retained-events}

The following argument permits every retained event coefficient to be
unbounded. It requires no canonical omission, actual-contact coordinates,
or source-shape entry. Five consecutive events in the eight-cycle suffice.

\begin{theorem}[Five consecutive lower bounds]
\label{p80253fe9bc1b:thm:five-retained-events}
Let the original canonical source be physical, with $d>0$ and all eight
event coefficients nonnegative. Suppose
\[
 d/m,\ d/c,\ d/A,\ d/D,\ d/\alpha
 \quad\hbox{are positive bounded coefficient elements}.
\]
Then the entire original parent is tensor SOS at every independent
projective measurement pair and every real or complex projective state.
There is no upper bound on any event coefficient divided by $d$.
By exact CHSH relabeling the same conclusion holds if any five
consecutive event coefficients in the cycle
$(\lambda,D,c,A,\alpha,m,\beta,n)$ have these lower bounds.
\end{theorem}

\begin{proof}
Divide by the positive pure coefficient $d$, so $d=1$.
Let $\mu=\min(1,m,c,A,D,\alpha)$. It is a positive bounded unit.
Select the parent $F_\mu$ with $m=c=A=D=\alpha=\mu$ and
$\lambda=n=\beta=0$. Keep every nonnegative excess and all three
unselected coefficients as their original event squares. This uses only
pure coefficient square roots. It makes no assertion that $F_\mu$ is
globally physical.

At noncommuting residue tests apply
the inherited noncommuting-residue theorem to the full physical source. It
performs its own positive coefficient normalization if the event
weights are unbounded.
It remains to certify $F_\mu$ near every commuting residue test. Its
coefficients are bounded, so a positive-definite real residue Gram is
sufficient by the bounded positive-residue Gram lemma.

Write $\mathcal H$ for the signed core with $d=1$ and
$D_u=1+u^2,D_v=1+v^2$. On the four open commuting axes its exact
homogeneous restrictions are
\[
\begin{aligned}
\mathcal H|_{u=0}&=2D_v(X^2+Y^2),\\
\mathcal H|_{v=0}&=2\{(uW-Y)^2+(X+uZ)^2\},\\
\mathcal H|_{u=\infty}&=2\{(W+vX)^2+(vY-Z)^2\},\\
\mathcal H|_{v=\infty}&=2D_u(X^2+Y^2).
\end{aligned}                                                   \tag{NW1}
\]
The remaining finite angle has nonzero residue. At $u=0$, the retained
$m,D$ packets force $W=Z=0$ on the core kernel. At $v=0$, the retained
$A,\alpha$ packets force $Y=Z=0$, and then $W=X=0$. At $u=\infty$,
the retained $c,A$ packets force $W=Y=0$, and then $X=Z=0$. At
$v=\infty$, the retained $c,D$ packets force $W=Z=0$ after the core
has forced $X=Y=0$. Thus the
selected residue state matrix is positive definite on every open axis,
for either sign of the other finite angle.

Normalize a nonzero testing state by a coordinate of largest absolute
value. Its coordinates are then bounded and one is a unit. At the four
double corners, write $F_{ab}$ for the homogeneous restriction of
$F_\mu$ to $(u,v)=(a,b)$. Directly,
\[
\begin{aligned}
F_{00}&=(2+4\mu)X^2+(2+12\mu)Y^2+4\mu Z^2,\\
F_{\infty0}&=(2+8\mu)(W^2+Z^2)+4\mu Y^2,\\
F_{0\infty}&=4\mu(W^2+Z^2)+2X^2+(2+12\mu)Y^2,\\
F_{\infty\infty}&=4\mu W^2+(2+4\mu)X^2+(2+12\mu)Y^2.
\end{aligned}                                                    \tag{NW2}
\]
Here infinity means the polynomial homogeneous chart, obtained by
dividing the finite parent by the appropriate pure-testing squared
angle. Every form is strict except its explicitly missing state
coordinate. Consequently the only nontrivial residue states are the
$W$ ray at $(0,0)$, the $X$ ray at $(\infty,0)$,
and the $Z$ ray at $(\infty,\infty)$.

At $(0,0)$ normalize $W=1$. The five local coordinates
$(u,v,X,Y,Z)$ are infinitesimal. The complete polynomial $F_\mu$
has no constant or linear term, and its degree-two part is
\[
 (2+8\mu)(u-Y)^2+2X^2
       +4\mu\{(v-X)^2+Y^2+Z^2\}.                    \tag{NW3}
\]
Its residue is positive definite: vanishing successively gives
$X=Y=Z=v=u=0$.

At the corner $(\infty,0)$ normalize $X=1$ and put $p=1/u$.
The five infinitesimal coordinates are $(p,v,W,Y,Z)$. The complete
homogeneous polynomial again has no terms below degree two, and its
quadratic part is
\[
 (2+4\mu)(W+v)^2+(2+4\mu)(p+Z)^2
                         +4\mu(W^2+Y^2+Z^2).        \tag{NW4}
\]
This is positive definite. At the corner $(\infty,\infty)$,
normalize $Z=1$ and put $p=1/u,q=1/v$. On $(p,q,W,X,Y)$ the exact
quadratic part is
\[
 (2+8\mu)(q-Y)^2+2X^2
                         +4\mu\{(p-X)^2+Y^2+W^2\}. \tag{NW5}
\]
It is positive definite as well. One can obtain (NW5) directly from
(NW3) by the exact symmetry
$(u,v;W,X,Y,Z)\mapsto(1/v,1/u;Z,X,Y,W)$, which permutes the five
equal retained weights.

For completeness, these quadratic calculations control the entire
polynomial. In each local chart, allocate every higher monomial to a
product of two of the five coordinate roots. The remaining monomial
has positive degree, hence is infinitesimal; its coefficient is bounded
because it is a real polynomial in the bounded unit $\mu$. Symmetrizing
gives an exact finite five-root Gram whose residue is (NW3), (NW4), or
(NW5). No inversion, root, or comparison mixing the two fields occurs.
The bounded positive-residue Gram lemma therefore certifies the full
selected polynomial, including all higher terms.

Restore the excess event squares, the positive pure-testing homogeneous
and projective state factors, and the positive coefficient $d$. The
zero state is immediate. Splitting a complex state into real and
imaginary parts proves the complex assertion.

Finally the exact full-parent Alice reflection acts on the displayed
cycle indices as $i\mapsto7-i$, and the swap/reflection
$(u,v;W,X,Y,Z)\mapsto(1/v,1/u;Z,X,Y,W)$ acts as $i\mapsto6-i$.
Their composition gives a one-step cycle rotation. Their full-parent
identities have only positive pure-testing homogeneous factors, so they
transport the five-packet conclusion to every cyclic translate.
\end{proof}

This theorem is a sufficient source domain, not an exhaustive zero-$d$
source cover. It requires five selected events to be positive and
noninfinitesimal relative to $d$. Its all-tests conclusion is independent
of actual contact angles and of all original event hierarchies. In
particular, it contains both C1 and C2 sources with all six retained
events bounded below relative to $d$.


\section{\texorpdfstring{An unbounded first actual angle forces a strong $C$ packet}{An unbounded first actual angle forces a strong C packet}}
\label{component:pb7189b9c1f6e}

\label{pb7189b9c1f6e:sec:c2-weak-hi-hu-empty}

\begin{theorem}[The directed weak-$C$ types HI and HU are empty]
Use the exact positive rational C2 source, with $d=1$ and $n\geq0$.
If the first actual half-angle $u_c$ is unbounded and the second
actual half-angle $v_c$ is bounded, then $1/C$ is bounded.
Consequently no exactly-one-weak source oriented with $C$ weak can
have directed actual type HI or HU.  Literal $n=0$ is included.
\end{theorem}

This is an original-source impossibility theorem.  It makes no
claim that a scalar contact identity by itself is a tensor-SOS
certificate.  All independent testing variables and state coordinates
remain unrestricted; the two source leaves are empty before any
testing restriction is imposed.

\begin{proof}
Write
\[
 q=u_c^{-2}>0,\qquad y=v_c^2>0,
 \qquad H=u_c^2h=v_cY_c\in(0,1),
\]
and retain the original $r\in(0,1)$.  Here $q$ is infinitesimal
and $y$ is bounded; $y$ may itself be infinitesimal.  Put
\[
 R=r/q=-u_cX_c,\quad B=u_cv_cW_c,\quad
 L=v_c(u_cW_c-Y_c),\quad T=B-yR.
\]
The exact positive gaps give
\[
 B=H+L=yR+T,\qquad 0<T<H+y,\qquad L>0.
\]
In particular $H,T,y,r$ are bounded.  The original source coordinates
become exactly
\[
 z=1/q,\quad a^2=yq,\quad b=yr+qT,\quad h=qH,
 \quad t=r/q+(T-H)/y,
\]
and the two gaps satisfy
\[
             g=H+(1+y)r-1+qT>0,
 \qquad j=(H+y-T)/y>0.                             \tag{WE1}
\]
The displayed expressions are pure coefficient identities; no
inverse of a mixed source/testing element is used.

We first show that $R$ is bounded.  Set $Q=q^2\Delta>0$.
Direct expansion of the unchanged original formula for $n$ gives
\[
                         n=-\frac{qP}{2rQ},        \tag{WE2}
\]
where the following exact decomposition of $P$ is sufficient:
\begin{align*}
P={}&r^2(1+y)(H+yr)-Tqr(y+2)-q^2(T^2+T)+P_+,\\
P_+={}&HTq^2+HTqr+Hqr(y+1)+T^2q^2r\\
     &+Tq^2r(y+1)+2Tqr^2(y+1)+qr^2y(y+1)>0.        \tag{WE3}
\end{align*}
Suppose for contradiction that $R$ is unbounded, and put
$\kappa=q/r=1/R$, which is infinitesimal.  Let $S=H+yr>0$.
The ratio $T/S$ is bounded.  Indeed, if $r$ is a unit, then
\[
                T<H+y\leq S/r.
\]
If $r$ is infinitesimal, (WE1) and boundedness of $y,T$ imply
$H_0=1$, so $S$ is a unit and $T/S$ is again bounded.
Since $T$ is bounded, $(T^2+T)/S$ is bounded as well.  Dividing
the lower bound in (WE3) by the positive $r^2S$ gives
\[
 \frac{P}{r^2S}\geq (1+y)
       -\kappa\frac{T(y+2)}S
       -\kappa^2\frac{T^2+T}S.
\]
The last two terms are infinitesimal, while $1+y\geq1$.
Thus $P>0$.  Equation (WE2) contradicts $n\geq0$, including
the equality case $n=0$.  Therefore $R$ is bounded.

Now $r=qR$ is infinitesimal.  Equation (WE1) becomes
\[
             g=H-1+q\{(1+y)R+T\}>0,
\]
so $H_0=1$.  Both $B=yR+T$ and $L=B-H$ are bounded.
The original denominator has the exact form
\begin{align*}
 \Delta={}&(1+q)\{B^2+2BR+(1+y)B\}
           +(1-yq)R^2+2(1+y)R+1+y.                \tag{WE4}
\end{align*}
It is a positive bounded unit: all quantities are bounded, and
$1-yq$ is a positive unit, with constant contribution $1+y\geq1$.
The unchanged numerator $N_{\rm src}$ satisfies
\[
 \frac{N_{\rm src}}q
  =(1+q)(B+R)(H+y)-(1+y)(1-qR)L.                  \tag{WE5}
\]
Since $H_0=1$ and $L=B-H$, its residue is
\[
       (N_{\rm src}/q)_0=(1+y_0)(1+R_0)>0.
\]
Thus $N_{\rm src}/q$ is a positive bounded unit.  Finally,
$A_c=1+R$ and $a^2t=qL$, so the original coefficient is exactly
\[
          C=\frac{(1+R)(N_{\rm src}/q)}{2L\Delta}.
                                                               \tag{WE6}
\]
The numerator has bounded positive inverse, and $L\Delta$ is
bounded.  Therefore $1/C$ is bounded, as asserted.  No lower bound
on $L$ was required; if $L$ is infinitesimal the conclusion is
stronger.
\end{proof}

The exact C2 reflection preserves directed type HI and exchanges
$C,\alpha$.  Applying the theorem in both orientations shows that
every physical HI source has both complementary packets strong.
For type HU the reflected type is UI; the reflected consequence is
that $\alpha$ is automatically strong at every UI actual contact.
These observations respect the directed weak-$C$ orientation; they
do not identify IH with HI.


\section{\texorpdfstring{Further actual-angle reductions from exact source entry}{Further actual-angle reductions from exact source entry}}
\label{component:pe56fb394b522}

\label{pe56fb394b522:sec:c2-angle-reductions}

All angles in this section are actual source contacts. They are strictly
positive finite field elements; ``infinite'' below means unbounded in
the coefficient field, not a literal projective endpoint. Write
$z=c^2=u_c^2$, $y=a^2z=v_c^2$, and
$A_c=1+zr$, $S=A_c+a^2zt$. The exact source identities used are
\[
 n\ge0\ \Longleftrightarrow\ (b+r)A_cS-r\Delta\ge0,
 \qquad
 N=(1+z)(b+r)h+a^2\{g+j+rt(1+y)\}>0.
\]

\begin{theorem}[The ordinary--ordinary source orbit is already closed]
If both actual C2 half-angles are positive bounded units, then the
source belongs either to the two-weak-complementary-event all-tests
exit or to the five-consecutive-strong-event all-tests exit. No upper
bound on an original event coefficient is required.
\end{theorem}
\begin{proof}
Here $a,z$ are units. The positive gaps give $0<h<1/z$ and
$0<t<1+zr$, so $h,t,b=h+a^2t$ are bounded. From
$g=b-1+r+zh>0$ we obtain $1<b+r+zh\le(1+z)b+r$;
thus $b+r$ is a positive unit. The positive expansion of $\Delta$
makes it a bounded unit too. The original $A$ formula has unit
numerator and bounded denominator $h$, so $A$ is strong.

The exact numerator for $n$, after substituting $t=(b-h)/a^2$,
satisfies the polynomial specialization
\[
 \left[(b+r)A_cS-r\Delta\right]_{r=1}
     =-z(a^2+h)(b+1)(z+1).
\]
If $1-r$ were infinitesimal, taking residues would contradict
$n\ge0$: the right side is strictly negative. Hence $1-r$ is a
unit. The original $m$ numerator is now a unit and $g\Delta$ is
bounded, so $m$ is strong. Apply the exact C2 reflection; it preserves
ordinary actual angles and exchanges $m,D$. Thus $D$ is strong.

If either $C$ or $\alpha$ is weak, its original rational formula
implies $N$ infinitesimal: all the multiplying denominator factors
are bounded and $z,a,b+r,\Delta$ have bounded inverses. The positive
identity for $N$ then forces $h,j$ infinitesimal. Consequently
$t=A_c-j$ and $1-zh$ are units. Both original complementary formulas
therefore give $C,\alpha$ infinitesimal. The two-weak theorem applies.
Otherwise $C,\alpha$ are both strong; with $m,A,D$ strong, the
five-consecutive-event theorem applies. These cases exhaust the source.
\end{proof}

\paragraph{Some directed one-weak angle types are empty.}
Orient an exactly-one-weak source so that $C$ is weak and $\alpha$
strong. In addition to the ordinary--ordinary exclusion above, its
actual angle types cannot be $(\mathrm{unit},\mathrm{Inf})$,
$(\mathrm{Inf},\mathrm{unit})$ or $(\mathrm{Inf},\mathrm{Inf})$.
Here $\mathrm{Inf}$ denotes positive infinitesimal.

For $(\mathrm{unit},\mathrm{Inf})$, $z$ is a unit and $a^2$ is
infinitesimal. Again $t,h,b,\Delta$ are bounded, $\Delta$ is a unit,
and $b+r$ is a unit. Weak $C$ implies $N/a^2$ infinitesimal.
The positive identity forces $j$ and $rt$ infinitesimal. Since
$t=A_c-j$ is a unit, $r$ is infinitesimal. Then $g>0$ makes
$b+zh$ a unit, and $b=h+a^2t$ makes $h$ a unit. The first positive
summand of $N$ is consequently a unit, contradicting $N/a^2$
infinitesimal.

For $(\mathrm{Inf},\mathrm{unit})$, bounded $b$ is a unit by the
inherited source-entry inequality. Explicitly,
\[
 xC(b-h)=\frac{A_c(xN)}{2\Delta},\qquad
 xN=(1+x)(b+r)(xh+k^2)-x(1+k^2)(1-r)(b-h).
\]
The last numerator has residue $k_0^2(b_0+r_0)>0$, independently
of any lower bound on $g$. Thus $xC(b-h)$ is a unit and $C$ is strong. For unbounded
$b$, (SR1) gives $C\asymp1/L$ with bounded $L$, again strong.
This exhausts $b$ and excludes weak $C$ in this angle orientation.

For $(\mathrm{Inf},\mathrm{Inf})$, both $z,y$ are infinitesimal
and $t<A_c$ is bounded. Put $H=zh$, $L=z(b-h)=yt$ and $B=zb=H+L$.
Here $0<H<1$, $L$ is infinitesimal, and $B$ is bounded. If $H$ is
a unit then $b$ is unbounded; direct division of the original
$\Delta$ by $b$ gives a positive unit of residue $1+B_0$.
The positive first summand of $N$ consequently gives
$C\ge c_0H/L$ for some positive real $c_0$, contradicting weak $C$.
Thus weak $C$ forces $H,B$ infinitesimal. The exact ratio is then
\[
 \frac\alpha C=
 \frac{z^2(b+r)(b-h)}{A_c(1-H)}
 =\frac{(B+zr)L}{A_c(1-H)},
\]
which is infinitesimal. Therefore $\alpha$ is weak as well, contrary
to the selected one-weak orientation.

After these exclusions, an exactly-one-weak source oriented with weak
$C$ has only the following candidate actual types:
\begin{gather*}
 (\mathrm{Inf},\mathrm{unbounded}),\quad
 (\mathrm{unit},\mathrm{unbounded}),\quad
 (\mathrm{unbounded},\mathrm{Inf}),\\
 (\mathrm{unbounded},\mathrm{unit}),\quad
 (\mathrm{unbounded},\mathrm{unbounded}).
\end{gather*}
The second type is the reflection of the single-angle domain and
is now fully covered by the homogeneous unbounded-$b$ theorem and the
other single-angle admissions. The four other listed directed types
remain conditional source leaves; no nonemptiness is inferred from
their listing.

\paragraph{Other strong-event angle types are not all empty.}
For $x>0$ infinitesimal, the exact positive chart
\[
 c=\sqrt x,\quad a=1,\quad b=1/x,\quad r=1/4,\quad t=1/2
\]
has $h=1/x-1/2$, $g=1/x+1/4-x/2$, $j=1/2+x/4$ and all
remaining chart gaps positive. Its original coefficients are
\begin{align*}
 n&=\frac{8-5x^2}{16(1+x)^2},&
 m&=\frac{3x^2(3x+4)}{16(1+x)(4+x-2x^2)},\\
 C=\alpha&=\frac{4-x}{8x},&
 A&=\frac{x(x+4)(3x+4)}{32(2-x)(1+x)^2},\\
 D&=\frac{(2-x)(3x+4)}{16x(1+x)}.
\end{align*}
Every coefficient is positive. The unchanged exact contact is
\[
 (u_c,v_c;W,X,Y,Z)
 =(\sqrt x,\sqrt x;1/x,-\sqrt x/4,1/\sqrt x-\sqrt x/2,1).
\]
Original source stationarity, positive transverse state inertia and
Laurent convexity prove physicality, including projective limits.
Here $C,\alpha,D$ are unbounded, $m\sim3x^2/16$, $A\sim x/4$,
and $n\to1/2$. Thus a genuine strong-complementary source can have
two infinitesimal actual angles and weak $m,A$. This is an obstruction
to deleting that angle orbit, not a counterexample to tensor SOS.


\section{\texorpdfstring{The joint rational C2 opening: a full-parent certificate}{The joint rational C2 opening: a full-parent certificate}}
\label{component:p8f912583feaa}

\label{p8f912583feaa:sec:c2-joint-centered}

\begin{theorem}[Both event openings may close]
Use the exact rational C2 source and original parent of
Section~\ref{component:pec59340d644f}.  Suppose $x>0$ is infinitesimal,
$k,b$ are positive bounded units, $0<r<1$, $0<h<b$, and
$g=b-1+r+xh>0$.  Suppose $r,g$ are infinitesimal, $n\geq0$, and
the source is physical.  Then the entire original parent is tensor
SOS at every independent real or complex projective testing state
and every independent projective pair of testing angles.  No lower
bound on either $h$ or $b-h$ is required.
\end{theorem}

Together with the three-opening and last-opening theorems this
removes the condition that at least one of $r,g$ be a unit in the
bounded-$k,b$ rational C2 source family.  It does not classify all
physical one-weak C2 sources, nor admit a nonunit $k$ or unbounded $b$.
All tensor-ring conventions and the strict bounded tensor Gram lemma
are those of the authoritative package.  In particular, a root vector
need not be linearly independent, and its entries may be unbounded.

\paragraph{Source bounds, including both complementary openings.}
Write $\rho=\sqrt{x}$, $\eta=h/b$, $\kappa=b/k$, $K=k_0>0$,
and $D_K=1+K^2$.  Since $r,g$ are infinitesimal, $b_0=1$.
The source formulas give
\[
 (nr)_0=\frac1{4D_K},\qquad (mg)_0=\frac14,
 \qquad \Delta_0=2D_K.
\]
Thus $n,m$ are positive unbounded pure coefficient quantities,
$nr,mg$ are bounded units, and $nr^2,mg^2$ are infinitesimal.
In particular $m-1>0$.  The existing opening cancellation identities,
with $T_C=xC(b-h)$ and $T_A=xAh$, give bounded positive units
\[
 (T_C)_0=(T_A)_0=\frac{K^2}{4D_K},\qquad
 (D/x)_0=\frac1{4K^2D_K},\quad
 (\alpha/x)_0=\frac{K^2}{4D_K}.
\]
They hold even if $h$ or $b-h$ is infinitesimal.  Specifically,
\[
 xC(1-\eta)=T_C/b,\quad xA\eta=T_A/b,\quad
 xC(1-\eta)^2=T_C(b-h)/b^2,\quad
 xA\eta^2=T_Ah/b^2.                                  \tag{JC1}
\]
Put $H_c=x(C+A)$.  This quantity is bounded below by a positive
real, although it can be unbounded.  Its reciprocal and reciprocal
square root are bounded pure coefficient elements.

\paragraph{Center the excess positive packets before taking normals.}
Retain exactly the original amplitudes
\[
 M=vW-X+uvY-uZ,\qquad J=uW+uvX-Y-vZ,
 \qquad D_u=1+u^2,\quad D_v=1+v^2.
\]
Let $F$ denote the full original parent, including its signed core
$H$.  Define $G$ by the following exact polynomial identity:
\[
 F=G+4nD_uD_v(X+ruZ)^2
             +4(m-1)(M-guZ)^2.                       \tag{JC2}
\]
An explicit expression, requiring no elimination or mixed inverse, is
\begin{align*}
G={}&4M^2+4CD_v(uW-Y)^2+4AD_uD_vY^2
       +4DJ^2+4\alpha D_u(vY-Z)^2+H\\
 &-8nrD_uD_vuZX-4nr^2D_uD_vu^2Z^2\\
 &+8(m-1)g\,uZM-4(m-1)g^2u^2Z^2.                  \tag{JC3}
\end{align*}
Every coefficient in (JC3) except $C,A$ is bounded.  At the exact
original contact
\[
 (u,v;W,X,Y,Z)=(\rho,k;\kappa\rho,-r\rho,xh/k,1)
\]
both subtracted squares and all of their first derivatives vanish.
Consequently $G$ has exactly the same zero and full stationary
contact as $F$.  This statement is essential: merely capping the
uncentered original packets would not preserve the contact.

On the finite $u=0$ residue axis use
\begin{gather*}
 X_g=-ruZ,\qquad Y_g=\eta uW,\qquad E_X=\sqrt n(X+ruZ),\\
 E_Y=\sqrt{C+A}(Y-\eta uW),\qquad
 E_M=\sqrt{m-1}(M-guZ).                              \tag{JC4}
\end{gather*}
Expand only $G$ in $E_X,E_Y$ and retain $4E_M^2$ as written.
The root $E_M$ is allowed to be redundant.  No substitution for
$W$ is made.  In particular, the positive $C$ packet introduces
no coefficient of the form $Cu^2/n$ or $Cu^2/m$.

\paragraph{The complete graph scalar and its global leading inequality.}
Write $f=G|_g=f_{ww}W^2+f_{wz}WZ+f_{zz}Z^2$.  The same exact monic
divisions as in the prior single-angle proof give
\[
 f=aP^2+b_sB^2+c_sV^2+eE^2+g_sBV,
 \qquad (P,B,V,E)=(uW/\rho,uZ,vW,\rho Z),              \tag{JC5}
\]
where
\[
 a=x\frac{f_{ww}-f_{ww}|_{u=0}}{u^2},\quad
 b_s=\frac{f_{zz}-f_{zz}|_{u=0}}{u^2},\quad
 c_s=\frac{f_{ww}|_{u=0}}{v^2},\quad
 e=\frac{f_{zz}|_{u=0}}x,\quad g_s=\frac{f_{wz}}{uv}.
\]
These are polynomial coquotients, not inverses of testing
infinitesimals.  Their coefficients are bounded by (JC1)--(JC3).
The exact constant $f_{zz}|_{u=0}=4Dv^2+4\alpha$ supplies the $x$
factor.  For finite $v$ their residues are
\[
 q(v)=\frac{K^2(1+v^2)}{D_K},\qquad
 b_s=4,\quad c_s=4,\quad g_s=-10,\qquad
 e(v)=\frac{v^2/K^2+K^2}{D_K}.                        \tag{JC6}
\]
Thus the leading scalar, using the exact identity $BV=vPE$, is
\[
                  qP^2-2vPE+eE^2+4(V-B)^2.           \tag{JC7}
\]
The key identity is now an identity for the graph of a full
bounded remainder, rather than for the overconstrained graph from
the earlier frontier note:
\[
                   q(v)e(v)-v^2
                    =\frac{(v^2-K^2)^2}{D_K^2}.      \tag{JC8}
\]
If $v$ is a finite nonzero residue with $v^2\ne K^2$, then
$10|v|<8|v|+2\sqrt{qe}$.  Splitting the exact $BV=vPE$ cross
term with the pure testing function
\[
             \theta=\frac{4|v|}{4|v|+\sqrt{q(v)e(v)}}
\]
gives strictly positive residue blocks on $(B,V)$ and $(P,E)$.
At $v_0=0$, allocate the entire cross term using $BV=vPE$;
(JC5) then has strictly positive diagonal residue
$(q(0),4,4,e(0))$ and zero cross residues.  There is no exceptional
$n_0=0$ corner in the joint regime.

Here is also a global equality and Hessian check needed at the
matched angles.  For $z=v^2>0$ and $s>0$, set
\[
 \mathcal D(s,z)=(q s+4z)(4s+e)-25zs.
\]
Exactly,
\begin{align*}
 \frac{\mathcal D(s,z)}s
 &=4qs+4ez/s+qe-9z\\
 &=4\bigl(\sqrt{qs}-\sqrt{ez/s}\bigr)^2
       +(\sqrt{qe}-\sqrt z)(\sqrt{qe}+9\sqrt z).    \tag{JC9}
\end{align*}
The second term is nonnegative by (JC8); its equality forces
$z=K^2$.  The first then forces $s=1$.  Moreover the second term is
\[
 \frac{(z-K^2)^2}{D_K^2}
       \frac{\sqrt{qe}+9\sqrt z}{\sqrt{qe}+\sqrt z},
\]
whose quadratic coefficient at contact is strictly positive.
The first square supplies a nonzero $s$ derivative there.
Hence $\mathcal D$ has a positive definite Hessian at $(1,K^2)$.
These radicals are used only in the residue-field inequality; no
radical of an element mixing coefficient and testing fields is used.

\paragraph{Restore every normal, with bounded coefficients.}
From (JC2)--(JC4) the entire parent is exactly
\[
 F=f+E_XL_X+E_YL_Y+
 H_{XX}E_X^2+2H_{XY}E_XE_Y+H_{YY}E_Y^2+4E_M^2.        \tag{JC10}
\]
The first normals are $L_X=G_X|_g/\sqrt n$ and
$L_Y=G_Y|_g/\sqrt{C+A}$.  The only unbounded weights $C,A$ do not
enter the $X$ row.  Thus
\[
 L_X=V\ell_w+B\ell_z
\]
with bounded coefficient functions carrying the infinitesimal
pure coefficient factor $1/\sqrt n$.  Both have identically zero
residue on the entire finite $u=0$ axis, including $v=0$.

For the second normal the exact form is
\[
 L_Y=\frac1{\sqrt{H_c}}\left\{
 P(B_2+xe_w)+B(\rho u e_z)+E e_*\right\},            \tag{JC11}
\]
where
\[
 B_2=8D_v\{(xA)\eta D_u-(xC)(1-\eta)\},\qquad
 e_*=8v(D-\alpha).
\]
All $e_w,e_z$ are bounded: all new terms in (JC3) have bounded
coefficients, and their $Y$ derivative has an additional $u^2Z$
factor.  Equation (JC1) makes $B_2$ bounded with identically zero
residue, even at either closing complementary opening.  Thus every
coefficient in (JC11) has zero residue.

The entire hard block is bounded and has residue
\[
             \operatorname{diag}(4D_{v,0},4D_{v,0},4)
                   \quad\hbox{on }(E_X,E_Y,E_M).      \tag{JC12}
\]
The first entry comes from the retained centered $n$ packet;
the bounded $G$ contribution is divided by $n$.  The second is
\[
 4D_v\frac{C+AD_u}{C+A}
   =4D_v\left(1+\frac A{C+A}u^2\right)
\]
plus bounded terms divided by $C+A$.  The mixed entry contains
$1/(\sqrt n\sqrt{C+A})$ times a bounded function.  The third
entry is the unexpanded centered $m$ packet.  This proves the
full-parent assertion at all remote finite-axis tests, by the
strict Gram in (JC5) and the zero-residue first normals.

\paragraph{The matched finite-angle residues.}
At $v_0=K$, normalize $Z=1$ and suppose $W$ is infinitesimal.
Use exactly the prior contact conductors
\[
 R=(uW-\kappa x)/\rho,\quad S=u-\kappa^{-1}W,
 \quad U=\rho(v-k),\quad T=u(v-k),
\]
with exact identities
\[
 u^2-x=uS+\kappa^{-1}\rho R,\quad
 W^2-\kappa^2x=-\kappa WS+\kappa\rho R,
 \quad uW-\kappa x=\rho R.                           \tag{JC13}
\]
Here $f=p^2Q_b/x+B_b$, with bounded $B_b$ and
\[
 Q_b=4D_v\{(xC)(1-\eta)^2+(xA)\eta^2D_u\}.
\]
Boundedness follows from (JC1); the remaining terms use only
the bounded coefficients of (JC3).  The full exact stationary
contact of $G$ proved after (JC3) supplies every vanishing required
by the prior monic-contact-division argument.  Successive monic
differences in $u^2-x,W^2-\kappa^2x,uW-\kappa x,v-k$ therefore
give a bounded Gram on $R,S,U,T$.

For precision, its strict residue follows from the same explicit
Schur-complement step in that argument, now with
\[
 \Phi=q(v)s^2w^2+4s^2+4v^2w^2-10vsw+e(v).
\]
Completing the $w$ square gives
\[
 \Phi=(qs^2+4v^2)
       \left(w-\frac{5vs}{qs^2+4v^2}\right)^2
       +\frac{\mathcal D(s^2,v^2)}{qs^2+4v^2}.
\]
Equation (JC9) proves a positive definite three-variable Hessian
at $(s,w,v)=(1,1/K,K)$.  Consequently the residue before the
conductor-zero correction has form
\[
 q(K)R^2+4S^2+A_1RU+B_1ST+E_1U^2,
 \qquad E_1>A_1^2/(4q(K))+B_1^2/16.
\]
Choose a real $\delta$ strictly between $B_1^2/16$ and
$E_1-A_1^2/(4q(K))$, and add $\delta$ times the exact zero
\[
 T^2-U^2-ST(v-k)-\kappa^{-1}RU(v-k)=0.
\]
The two residue blocks $(R,U)$ and $(S,T)$ are strictly positive.

An explicit equivalent residue check can be used in place of this
Schur choice.  The capped current has the exact conductor lift
\[
 (M-gu)|_g=-v\kappa S+\kappa T
                  +\eta uv\rho R+\eta u\kappa\rho U.
\]
Retaining its $T^2$ term, rather than first moving that term to $U^2$,
gives the scalar residue
\[
 K^2R^2+\frac{2(K^2-1)}{D_K}RU+\frac{U^2}{K^2}
                  +4(S-T/K)^2.
\]
Adding $2/D_K^2$ times the same conductor zero leaves the $(R,U)$
Schur margin $2/D_K^2>0$, and leaves
$4(S-T/K)^2+(2/D_K^2)T^2$ on $(S,T)$, also strictly positive.

The full-normal lifts are also the ones actually required for
(JC10), rather than scalar-only lifts.  Write
$L_X=uA_1'+WB_1'=-\kappa B_1'S+u(A_1'+\kappa B_1')$.
Both functions and all their bounded monic coquotients retain
$1/\sqrt n$.  Exact $G_X=0$ at contact supplies vanishing of
$A_1'+\kappa B_1'$.  Using (JC13) gives zero-residue coefficients
on $R,S,U,T$.  For (JC11), write its numerator as $RB_2+\rho E_2$.
Both $B_2,E_2$ have identically zero leading functions and $E_2$
vanishes at contact.  The same monic divisions give factors
$x,\rho u,\rho W$, or a zero-leading-function $(v-k)$ coquotient;
division by $\sqrt{H_c}$ preserves boundedness and zero residue.
Thus the complete seven-root Gram on
$(R,S,U,T,E_X,E_Y,E_M)$ is strict in residue.

If $W/Z$ is not infinitesimal, normalize $W=1$ and keep $Z$
bounded.  Then the scalar has a strict residue on $(P,1)$:
its positive constant contribution is $4K^2$.  All normalized
first normals have zero residue as above.  For $Z=0$, the original
parent is already the literal sum of its event packets and three
positive core squares.  The simultaneous sign copy gives $v_0=-K$.

\paragraph{Every other independent test.}
The outer-chart proof need not use the centered decomposition.
Use the original parent, cap $C,A$ at $\sigma/x$, where
$\sigma=\min(1,xC,xA)$ is a positive bounded unit, and cap $m$ at
$1$.  Retain the original nonnegative excess packets and the
entire original $n$ packet.  The four exact outer-core identities
in the prior single-angle theorem apply without change: their
strict residues require only $\sigma_0>0$, the selected $m=1$,
and $(D/x)_0>0$.  They never bound $n$.  This covers the finite
unit $u$ axis, both reciprocal axes, all their intersections, and
literal projective ends.  The inherited physical
noncommuting-residue theorem covers the remaining tests.

These charts and the entire finite $u=0$ axis exhaust all independent
projective measurement pairs.  All state charts were restored above;
complex states split into their real and imaginary parts as in the
prior proof.  Every inverse or radical used in a tensor coefficient
is pure coefficient or pure testing.  No comparison of a source
infinitesimal with a testing infinitesimal, and no mixed inverse
or radical, occurs.  Restoring the retained packets proves the theorem.


\section{\texorpdfstring{Strong C2 pitchfork with arbitrary complementary scales}{Strong C2 pitchfork with arbitrary complementary scales}}
\label{component:pf7943b34a7dc}

\label{pf7943b34a7dc:sec:strong-complement-pitchfork}

\begin{theorem}[No upper bound on the two complementary events]
\label{pf7943b34a7dc:thm:strong-complement-pitchfork}
Let $K_c,K_t\supset\mathbb R$ be independent real closed fields.
Let the original canonical C2 parent be physical over $K_c$, with
$\lambda=\beta=0$, $d$ a positive bounded unit, $C,\alpha>0$
having bounded reciprocals, and $n,m,A,D\geq0$ arbitrary.
Suppose it has an actual finite noncommuting stationary zero with
$u_c=k>0$ infinitesimal and $v_c=a>0$ a bounded unit.
Then the entire original parent is tensor SOS at every independent
projective measurement pair and every real or complex projective state.
There is no upper bound or relative-scale condition on any event weight.
\end{theorem}

The case of bounded $C,\alpha$ is the existing three-large-event
theorem. We prove both remaining cases, without modifying the parent
or identifying the source and testing fields. We use
\[
\begin{gathered}
 D_u=1+u^2,\quad D_v=1+v^2,\quad r=uW-Y,\quad p=vY-Z,\\
 J=uW+uvX-Y-vZ,\quad R_2=vW-X+up,\quad B=X+uZ,\\
 H=2\{J^2+D_up^2+D_vB^2-D_uD_vZ^2\},\\
 F=dH+4nD_uD_vX^2+4mR_2^2+4CD_vr^2
       +4AD_uD_vY^2+4DJ^2+4\alpha D_up^2.              
\end{gathered}
\tag{SC1}
\]
Every square in this formula retains its original coefficient.

\paragraph{Necessary reductions from the full parent.}
Cap only $C,\alpha$ at positive bounded units in the initial hard
state estimate. At the actual angles the bounded selected form is
bounded below by
$b(X^2+Y^2+Z^2)-B_0k^2W^2$, with positive real $b,B_0$.
Every retained packet is nonnegative. Hence the contact has $W_c\ne0$;
on setting $W_c=1$,
\[
 X_c,Y_c,Z_c=O(k),\qquad m=O(k^2).
\]
Each positive event packet is $O(k^2)$ as well. Thus, writing
$y=Y_c/k$, $z=Z_c/k$, $x_1=X_c/k$, $j=J_c/k$,
\[
 Ay^2,\ Dj^2,\ C(1-y)^2,\ \alpha(ay-z)^2=O(1),\qquad
 j=1-y-az+aX_c.                                      \tag{SC2}
\]
If $C$ is unbounded then $\bar y=1$, so $A$ is bounded. If $D$
were also unbounded, $j,z$ would be infinitesimal, and the residue of
\[
 H_c/(2k^2)=j^2+(1+k^2)(ay-z)^2
            +(1+a^2)(x_1+Z_c)^2-(1+k^2)(1+a^2)z^2
\]
would be $\bar a^2+(1+\bar a^2)\bar x_1^2>0$.
This contradicts $H_c\leq0$ at the zero of (SC1). Therefore
$C$ unbounded forces both $A,D$ bounded. If $\alpha$ is unbounded
and $A$ were unbounded, (SC2) would make $y,z$ infinitesimal and
$\bar j=1$; the same residue would instead be
$1+(1+\bar a^2)\bar x_1^2>0$. Thus $\alpha$ unbounded forces $A$
bounded. It remains to prove precisely
\[
 \text{(I) } C\text{ unbounded},\ A,D\text{ bounded};\qquad
 \text{(II) } C\text{ bounded},\ \alpha\text{ unbounded},\ A\text{ bounded}.
                                                               \tag{SC3}
\]

\paragraph{Uniform stationary scale entry.}
Set $U=k^2$, $V=a^2$, $N=d+2n$, $\theta=n/N$, and
\[
\begin{gathered}
 t_2=d+2(\alpha+m+A),\quad T=t_3=d+2(C+D),\quad
 t_4=d+2(\alpha+m),\\
 t_6=d+2(\alpha+C+A),\quad t_8=d+2(C+A+D).
\end{gathered}
\]
The exact original Laurent stationary equations are
\[
\begin{split}
 U^2 A t_4\{\theta T+C(1+2D/N)V\}
     &=mV\{\alpha t_8+Dt_6V\},\\
 \alpha A(1+2m/N)T+\theta A Tt_4U
     &=\{mDt_6/U+CDt_2\}V^2.                         
\end{split}\tag{SC4}
\]
Define only pure coefficient quantities
\[
 \delta=CD/T,\quad e=(D/T)(t_6/\alpha),\quad b=t_2/\alpha,
 \quad \mathcal B=\theta+(C/T)(1+2D/N)V,
 \quad \mathcal G=t_8/T+eV.
\]
Because one of $C,D$ is bounded in either case (SC3),
$\delta\leq\min(C,D)/2$ is bounded. The identities
\[
 e=(D/T)\{2+(d+2A)/\alpha\}+2\delta/\alpha,
 \quad \mathcal G=1+2A/T+eV,
 \quad \mathcal B=\theta+(C/T)V+(2\delta/N)V
\]
show boundedness of all these quantities. The quantities
$b,\mathcal G,t_4/\alpha$ are positive bounded units. So is
$\mathcal B$: since $C,N$ have bounded reciprocals and $d$ is
bounded,
\[
 \frac{T}{C(N+2D)}\leq\frac d{CN}+\frac2N+\frac1C=O(1),
\]
and consequently
\[
 \theta+(C/T)(1+2D/N)
 =\frac{n+C(N+2D)/T}{d+2n}
 \geq\frac{n+c_0}{d+2n}
\]
for a positive real $c_0$. The bounded-unit factor $V$ preserves
this lower bound. Dividing (SC4) by $\alpha T$ gives exactly
\[
 \frac m{U^2}=A\rho,\qquad
 A\mathcal H=\delta bV^2,\qquad
 \rho=\frac{(t_4/\alpha)\mathcal B}{V\mathcal G},\quad
 \mathcal H=1+2m/N+\theta(t_4/\alpha)U-U\rho eV^2.      \tag{SC5}
\]
Here $\rho$ is a positive bounded unit and $\mathcal H=1+O(U)$.
These identities were obtained before inverting $A$ or $D$.
They show $A\asymp\delta$. Since $C$ has a bounded reciprocal,
$\delta$ infinitesimal implies $D$ infinitesimal.
If $A$ were infinitesimal, therefore, so would $D,m$ be.
But the exact physical test
\[
 u=v=1,\qquad (W,X,Y,Z)=(1,0,1,1)
\]
has $X=r=p=0$, $R_2=1$, $J=-1$, $H=-2$, and hence
$F=-2d+4m+16A+4D$. Its residue would be negative. We conclude
\[
 A\asymp1,\qquad D^{-1}=O(1),\qquad
 m/k^4\asymp1,\qquad \gamma=4ma^2/k^4\asymp1.        \tag{SC6}
\]
In case (I), $D$ is a bounded unit; in case (II), it may be
arbitrarily large. In both cases $\delta$ is a positive bounded unit.

\paragraph{Full state equations and bounded graph data.}
For later use, the three entire state derivatives of (SC1) are
\[
\begin{split}
 F_X/4={}&D_v[d+2nD_u]X+(d+2D)uvJ+d uD_vZ-2mR_2,\\
 F_Y/4={}&-(d+2D)J+(d+2\alpha)vD_up-2CD_vr
                  +2AD_uD_vY+2muvR_2,\\
 F_Z/4={}&-(d+2D)vJ-(d+2\alpha)D_up+d uD_vB
                  -dD_uD_vZ-2muR_2.                
\end{split}\tag{SC7}
\]
In case (I), the $Z$ equation first yields $\alpha p_c=O(k)$,
because $D$ is bounded and every core coordinate is $O(k)$.
Then the $Y$ equation yields $Cr_c=O(k)$. In case (II), $C,A$
are bounded, so the $Y,Z$ equations are a system with bounded right
side $O(k)$ and unit coefficient matrix
\[
 \begin{pmatrix}-1&a\\-a&-1\end{pmatrix}
\]
on $((d+2D)J_c,(d+2\alpha)(1+k^2)p_c)$.
They yield $DJ_c,\alpha p_c=O(k)$. Finally the $X$ equation in
either case yields $NX_c=O(k^2)$. Thus put $x=X_c/k^2$ and
\[
\begin{array}{ll}
 \text{(I)}&r_0=r_c/k,\ p_0=p_c/k,
       \qquad Nx,Cr_0,\alpha p_0=O(1),\\
 \text{(II)}&j_0=J_c/k,\ p_0=p_c/k,
       \qquad Nx,Dj_0,\alpha p_0=O(1).
\end{array}                                                    \tag{SC8}
\]
All unweighted graph constants are bounded as well. The reciprocals
of $N,C,D,\alpha$ used below are bounded pure coefficient elements.

\paragraph{Two exact graphs retaining the large packets.}
For independent finite testing angles, case (I) uses
\[
 X_g=xu^2W,\quad r_g=r_0uW,\quad p_g=p_0uW,
 \quad Y_g=(1-r_0)uW,\quad Z_g=\{v(1-r_0)-p_0\}uW.    \tag{SC9}
\]
Its exact normals and inverse are
\[
\begin{gathered}
 E_X=\sqrt N(X-xu^2W),\quad E_r=\sqrt C(r-r_0uW),
 \quad E_p=\sqrt\alpha(p-p_0uW),\\
 X=xu^2W+E_X/\sqrt N,\quad
 Y=(1-r_0)uW-E_r/\sqrt C,\quad
 Z=vY-p_0uW-E_p/\sqrt\alpha.                         
\end{gathered}
\tag{SC10}
\]
Case (II) uses
\[
 X_g=xu^2W,\quad J_g=j_0uW,\quad p_g=p_0uW,
 \quad Y_g=\frac{u(1-j_0+vp_0+xu^2v)}{D_v}W,
 \quad Z_g=vY_g-p_0uW.                              \tag{SC11}
\]
Here the exact normals and inverse are
\[
\begin{gathered}
 E_X=\sqrt N(X-xu^2W),\quad E_J=\sqrt D(J-j_0uW),
 \quad E_p=\sqrt\alpha(p-p_0uW),\\
 X=xu^2W+E_X/\sqrt N,\quad J=j_0uW+E_J/\sqrt D,
 \quad p=p_0uW+E_p/\sqrt\alpha,\\
 Y=(uW+uvX+vp-J)/D_v,\qquad Z=vY-p.                 
\end{gathered}
\tag{SC12}
\]
Both graphs match the entire actual contact. All denominators in
these formulas are pure coefficient factors or the pure testing
unit $D_v$. No inverse of a mixed tensor element occurs.
The possibly unbounded packets become respectively
\[
\begin{array}{ll}
 \text{(I)}&4(n/N)D_uD_vE_X^2,quad4D_vE_r^2,quad4D_uE_p^2,\\
 \text{(II)}&4(n/N)D_uD_vE_X^2,quad4E_J^2,quad4D_uE_p^2
\end{array}
\]
in the quadratic normal remainder. Their graph and first-normal
coefficients are bounded by (SC8), for example
$nx,nx^2,Cr_0^2,\sqrt C r_0,\alpha p_0^2,\sqrt\alpha p_0$
in case (I), and $Dj_0^2,\sqrt D j_0$ in case (II).
Every remaining source coefficient is bounded.

\paragraph{Entire scalar and entire first-normal division.}
Set $W=1$, $K=k^2$, and
$R=u^2-(K/a)v$, $T_s=u(v-a)$.
Since $R_2|_g=v+(p_0-x)u^2$ for both graphs, the entire scalar has
the exact shape
\[
 \Phi=F|_g=4mv^2+u^2H_0(u^2,v),                    \tag{SC13}
\]
with bounded polynomial coefficients in case (I), and bounded
coefficients in $B_c[v,D_v^{-1}]$ in case (II).
The three scalar contact jets give
$H_0(K,a)=-\gamma K$, $H_{0,U}(K,a)=\gamma$, and
$H_{0,v}(K,a)=-2\gamma K/a$. In particular the explicit expression
\[
 \Psi=H_0-\gamma u^2+2\gamma(K/a)v
\]
satisfies $\Phi-\gamma R^2=u^2\Psi$, and its value and both first
derivatives vanish at $(u^2,v)=(K,a)$. Finite Taylor division of
$\Psi$ therefore gives
\[
 \Phi=\gamma R^2+u^2\{(u^2-K)^2H_{20}
       +(u^2-K)(v-a)H_{11}+(v-a)^2H_{02}\},           \tag{SC14}
\]
with bounded coquotients. This is an exact bounded Gram on
$(R,T_s)$, with separate residue pivot $\bar\gamma$ on $R$.
For the rational case, Taylor division is legitimate because
\[
 D_v^{-r}-D_a^{-r}
 =-(v-a)(v+a)
   \frac{\sum_{i=0}^{r-1}D_v^iD_a^{r-1-i}}{D_a^rD_v^r}.
\]
Iterating this identity gives the second-order division; only pure
factors have been inverted.

Here are the full first normals, rather than initial forms. In
case (I) they are
\[
 \mathcal N_X=F_X|_g,\quad
 \mathcal N_r=-(F_Y+vF_Z)|_g,\quad
 \mathcal N_p=-F_Z|_g,                              \tag{SC15}
\]
and in case (II) they are
\[
 \mathcal N_X=\{F_X+(uv/D_v)(F_Y+vF_Z)\}|_g,\quad
 \mathcal N_J=-(F_Y+vF_Z)|_g/D_v,\quad
 \mathcal N_p=(vF_Y-F_Z)|_g/D_v.                     \tag{SC16}
\]
The derivatives in (SC7) specify every term. In particular the
following exact packet expressions expose the cancellations:
\[
\begin{split}
 -(F_Y+vF_Z)/(4D_v)
   &=(d+2D)J+2Cr-2AD_uY-duvB+dvD_uZ,\\
 (vF_Y-F_Z)/(4D_v)
   &=(d+2\alpha)D_up-2Cvr+2AvD_uY+2muR_2-duB+dD_uZ.
\end{split}\tag{SC17}
\]
The actual coefficients of the normal roots are obtained by dividing
these rows by $\sqrt N,\sqrt C,\sqrt\alpha$ in case (I), and by
$\sqrt N,\sqrt D,\sqrt\alpha$ in case (II).
Using (SC8), their exact forms after those divisions are
\[
 \mathcal N_X/\sqrt N=-8mv/\sqrt N+u^2L_X(u^2,v),
 \qquad \mathcal N_i/\sqrt{w_i}=uL_i(u^2,v),          \tag{SC18}
\]
where all $L_i$ have bounded coefficients in the indicated ring.
The contact equations give
$L_X(K,a)=8ma/(K\sqrt N)$ and $L_i(K,a)=0$.
Writing $L_X=L_X(K,a)+(u^2-K)L_{X1}+(v-a)L_{X2}$, and similarly
for $L_i$, gives the full exact row identities
\[
\begin{split}
 \mathcal N_X/\sqrt N
  &=(8ma/(K\sqrt N)+u^2L_{X1})R
       +u\{(K/a)L_{X1}+L_{X2}\}T_s,\\
 \mathcal N_i/\sqrt{w_i}
  &=uL_{i1}R+\{(K/a)L_{i1}+L_{i2}\}T_s.             
\end{split}\tag{SC19}
\]
Thus the $X$ row has infinitesimal couplings to both soft roots,
and the other rows have infinitesimal couplings to $R$.
The exact state degree is two, so the bounded quadratic remainder
already accounts for every higher-normal term. Equations
(SC10)--(SC19) give an exact full-parent bounded five-root Gram.

\paragraph{A common exact real Schur certificate.}
The $X$ and $R$ residue pivots are $2(1+\bar a^2)$ and
$\bar\gamma$. The unscaled boundary hard form on $Y,Z$ is
\[
 \mathsf B(v)[Y,Z]
 =sD_vY^2+4D(Y+vZ)^2+4\alpha(vY-Z)^2,
 \qquad s=4(C+A)+2d.                               \tag{SC20}
\]
In case (I), $Y=-E_r/\sqrt C$,
$Z=-vE_r/\sqrt C-E_p/\sqrt\alpha$ at $u=0$.
The form has the explicit pivots $4D_vE_r^2+4E_p^2$ plus a
bounded nonnegative form. In case (II),
$Y=(vE_p/\sqrt\alpha-E_J/\sqrt D)/D_v$ and
$Z=(-E_p/\sqrt\alpha-vE_J/\sqrt D)/D_v$.
The explicit pivots are $4E_J^2+4E_p^2$, again with a bounded
nonnegative remainder. The hard residue is strict in both cases.

The determinant and the exact boundary Schur numerator are
\[
\begin{gathered}
 \det\mathsf B(v)=4(1+v^2)\{s(D\,v^2+\alpha)+4D\alpha(1+v^2)\},\\
 Q_{10}=4\alpha A[d+2(C+D)],\qquad
 Q_{12}=4CD[d+2(\alpha+A)],\\
 Q_0=8ACD+8AC\alpha-4ADd-2Ad^2-2Cd^2-4Cd\alpha
       -2Dd^2-4Dd\alpha-d^3-2d^2\alpha.             
\end{gathered}
\tag{SC21}
\]
Put $S=\alpha T$. Every coefficient of $\det\mathsf B/S$ and of
$Q_{10}/S,Q_0/S,Q_{12}/S$ is bounded. This follows directly from
(SC21), bounded $A$, bounded $C/T,D/T,\delta$, and bounded
$1/\alpha$. Moreover
\[
 Q_{10}/S=4A\asymp1,\qquad
 Q_{12}/S=4\delta[d+2(\alpha+A)]/\alpha\asymp1.       \tag{SC22}
\]
The normalized determinant has a positive lower bound for every
bounded unit real $v$. Indeed it equals $C/T$ times the determinant
of the scaled hard form in case (I), and equals
$(D/T)D_v^2$ times that determinant in case (II). The factors
$C/T$ and $D/T$, respectively, are positive units in their cases.

We now take residues of these bounded coefficients and use an
ordinary real variable $v$. The exact real hard Schur function is
\[
 \bar f(v)=
 \frac{8(1+v^2)}{\overline{\det\mathsf B/S}(v)}
 \left\{\overline{Q_{10}/S}+\overline{Q_0/S}v^2
                      +\overline{Q_{12}/S}v^4\right\}.          \tag{SC23}
\]
This formula is an ordinary real Schur calculation, never an
inverse of $\det\mathsf B(v)$ in $K_c\otimes K_t$.
The contact first-normal identities identify the graph at $\bar a$
with the unique minimizer of the scaled real hard form. The scalar
contact jets give $\bar f(\bar a)=\bar f'(\bar a)=0$.
Because $\bar a\ne0$, the quartic bracket in (SC23) is exactly
\[
 \overline{Q_{12}/S}(v^2-\bar a^2)^2.
\]
Its strict second derivative supplies the soft Schur pivot
\[
 \tfrac12\bar f''(\bar a)=
 \frac{32(1+\bar a^2)\bar a^2\overline{Q_{12}/S}}
      {\overline{\det\mathsf B/S}(\bar a)}>0.
\]
The entire five-root Gram therefore has strictly positive residue
when $u$ is infinitesimal or zero and $\bar v=\bar a$.
The strict-residue Gram lemma proves tensor SOS. The exact identity
\[
 F(u,-v;W,-X,Y,-Z)=F(u,v;W,X,Y,Z)
\]
supplies $\bar v=-\bar a$.

\paragraph{All other independent tests, including literal endpoints.}
If $u$ is infinitesimal and $v$ is a bounded unit with
$\bar v\ne\pm\bar a$, retain the full original $m,n$ packets.
In case (I) use the roots
\[
 X,\quad K_C=\sqrt C r,\quad K_\alpha=\sqrt\alpha p,\quad uW
\]
and the exact inverse
$Y=uW-K_C/\sqrt C$, $Z=vY-K_\alpha/\sqrt\alpha$.
In case (II) use
\[
 X,\quad K_D=\sqrt D J,\quad K_\alpha=\sqrt\alpha p,\quad uW
\]
and the inverse in (SC12), with unshifted $J,p$.
The entire remaining parent is a bounded quadratic Gram on these
four roots. Its separate $X$ residue pivot is
$2\bar d(1+\bar v^2)$; its hard pair is the strict form (SC20)
in the scaled coordinates just displayed, and its soft Schur
complement is $\bar f(\bar v)>0$ by (SC23).

For the two endpoint neighborhoods $u,v\sim0$ and
$u,q=1/v\sim0$, divide by $d$, choose
\[
 \mu=\min(1,C/d,A/d,D/d,\alpha/d)>0,
\]
a bounded unit by (SC6), and retain all excesses together with
the full $m,n$ packets. The selected bounded parent has signed
weight one, $C=A=D=\alpha=\mu$, and $m=n=0$.
In the finite chart it is exactly a Gram on $(X,Y,Z,uW)$.
At $(u,v)=(0,0)$ its hard block is positive and its soft Schur
complement is $8Q_{10,\mu}/\det\mathsf B_\mu(0)>0$.
For the reciprocal chart use the polynomial homogeneous parent
$q^2F_\mu(u,1/q)$, on the same four roots. At $(u,q)=(0,0)$
its hard block is
$\operatorname{diag}(2,4(C+A+\alpha)+2,4D)$ in the
$(X,Y,Z)$ variables, evaluated at the capped weights, and its
soft Schur complement is
$8Q_{12,\mu}/\det\mathsf B_{\mu,\infty}>0$.
These are bounded polynomial Grams, so they include the full
infinitesimal neighborhoods and literal $v=0$ or $q=0$.

For completeness, the other open commuting axes use the same cap;
here the other angle has a nonzero finite residue.
At $v=0$ the core is $2\{(uW-Y)^2+(X+uZ)^2\}$ and the
$A,\alpha$ packets force $Y=Z=0$; finite nonzero $u$ then forces
$W=X=0$. At $u=\infty$ the homogeneous core is
$2\{(W+vX)^2+(vY-Z)^2\}$; the $C,A$ packets force $W=Y=0$,
and finite nonzero $v$ forces $X=Z=0$. At $v=\infty$ the core
is $2(1+u^2)(X^2+Y^2)$; the $C,D$ packets force $W=Z=0$
for finite nonzero $u$. Thus these complete state matrices are strict.

Only the two corners with $u=\infty$ remain. Put $p_u=1/u$.
At $(p_u,v)=(0,0)$ the selected state form is
$(2+8\mu)W^2+4\mu Y^2+(2+4\mu)Z^2$.
Its sole zero ray is $X$. In the chart $X=1$, the five small
variables $(p_u,v,W,Y,Z)$ have exact quadratic part
\[
 (2+4\mu)(W+v)^2+2(Z+p_u)^2+4\mu(W^2+Y^2+Z^2)\succ0.
\]
At $(p_u,q)=(0,0)$ the selected state form is
$4\mu W^2+(2+4\mu)X^2+(2+8\mu)Y^2$.
In its sole zero-ray chart $Z=1$, the variables
$(p_u,q,W,X,Y)$ have exact quadratic part
\[
 2X^2+(2+4\mu)(Y-q)^2+4\mu\{W^2+Y^2+(X-p_u)^2\}\succ0.
\]
There are no lower-order terms. Allocating every higher monomial
to two coordinate roots leaves an infinitesimal factor, giving a
strict-residue Gram for the entire selected polynomial at each
corner. Retain the original excess packets throughout.
After pure testing projective normalization, states outside each
displayed zero ray have positive residue value and are covered by
the scalar strict-residue lemma.

Noncommuting-residue tests use the inherited arbitrary-physical-source
theorem. This exhausts all independent projective tests. Finally
homogenize the local graph identities in $W$, replacing $R,T_s$
by $RW,T_sW$; they then include $W=0$. Restore all pure positive
coefficient and testing factors, and split complex states into real
and imaginary parts. No source infinitesimal was compared with an
independent testing infinitesimal anywhere in the argument.

\paragraph{Exact scope.}
This completes the strong C2 finite pitchfork for $C/d,\alpha/d$
bounded below, with arbitrary upper scales of all six events.
The reflected actual-contact domain follows by the established C2
reflection. It does not cover other actual-angle residue types or
sources in which either complementary ratio has no positive real
lower bound; no exhaustive global C2 join is asserted here.


\section{\texorpdfstring{A homogeneous scalar lift at the unbounded-state contact}{A homogeneous scalar lift at the unbounded-state contact}}
\label{component:pbb29c73cc124}

\label{pbb29c73cc124:sec:unbounded-b-homogeneous-scalar}

This lemma is a scalar lifting statement.  It neither supplies the full
state normals nor proves positivity of its coefficient Gram by itself.
Its purpose is to avoid any bound on the contact state ratio.

Let $B,\sigma>0$ be pure coefficient infinitesimals, let $k>0$ be a
pure coefficient bounded unit, and put $x=\sigma B$, $\xi=u^2$.
Suppose
\[
 f=a(\xi,v)W^2+2u d(\xi,v)WZ+c(\xi,v)Z^2,
 \quad
 a=\frac\xi B Q+\sigma^2 a_0,\qquad
 c=B e+\xi b,                                      \tag{HL1}
\]
where $Q,a_0,d,e,b$ have bounded polynomial coefficients.  Bounded
pure testing unit denominators may be cleared beforehand.
Assume the value and all first derivatives of $f$ vanish at
\[
 u=\sqrt{x},\quad v=k,\quad
 W=\frac1k\sqrt{B/\sigma},\quad Z=1.                \tag{HL2}
\]
No bound on $B/\sigma$ is assumed.

Define the homogeneous roots
\[
 R=\frac{uW-(B/k)Z}{\sqrt B},\qquad
 S=uZ-\sigma kW,\qquad
 U=\sqrt B(v-k)Z,\qquad T=u(v-k)Z.                 \tag{HL3}
\]
They satisfy the exact conductors
\[
 (\xi-x)W=u\sqrt B R+(B/k)S,
 \qquad
 (\xi-x)Z=uS+\sigma k\sqrt B R.                    \tag{HL4}
\]
Then $f$ has an exact symmetric Gram on $(R,S,U,T)$ whose entries
are bounded polynomials.  In particular no inverse of $\sigma$,
$B$, a testing infinitesimal, or a mixed expression is required in
these Gram entries.  The roots themselves need not be bounded.

\paragraph{First set $v=k$.}
In this paragraph every scalar coefficient is evaluated at $v=k$.
State stationarity and $u$ stationarity give
\begin{gather*}
 a_c=-\sigma k d_c,\qquad c_c=-B d_c/k,\\
 B(a_\xi)_c+kd_c+2kx(d_\xi)_c
                      +\sigma k^2(c_\xi)_c=0,     \tag{HL5}
\end{gather*}
where $c$ on a subscript means evaluation at $\xi=x$.
These are identities in the pure coefficient field; deriving them
does not insert an inverse into a tensor Gram.

Set, by monic polynomial division in $\xi$,
\[
 s=\frac{c+Bd/k}{\xi-x},\qquad
 q=-kd-\sigma k^2s,\qquad
 h_2=\frac{Ba+k\xi d+\sigma k^2c+\sigma Bkd}
                    {(\xi-x)^2}.                 \tag{HL6}
\]
All three quotients have bounded polynomial coefficients.  The first
numerator vanishes at $x$ by (HL5), and the last numerator and its
$\xi$ derivative vanish there by the same three identities.
Moreover $Ba=\xi Q+B\sigma^2a_0$ is bounded.  Division by a monic
polynomial preserves bounded coefficients and is not evaluation of
an inverse of $\xi-x$.

Direct multiplication gives the exact identity
\[
 f(\xi,k)=qR^2+sS^2
                  +h_2\left(uR+\frac{\sqrt B}{k}S\right)^2.
                                                               \tag{HL7}
\]
For later residue computations, if $E_0=e+d/k$, then
\[
 s=b+B\frac{E_0(\xi)-E_0(x)}{\xi-x}
       +x\frac{b(\xi)-b(x)}{\xi-x}.                \tag{HL8}
\]
Thus $s_0=b_0$, $q_0=-k_0d_0$, and the $h_2$ correction has zero
residue in every Gram entry.  Here subscripts $0$ denote residues,
not evaluation at the contact polynomial variable.

\paragraph{The complete angular coquotient.}
Write $\tau=v-k$ and, for each of $Q,a_0,d,e,b$, form its bounded
monic difference in $v$, denoted respectively by
$Q_1,a_{01},d_1,e_1,b_1$:
\[
 Q(\xi,v)-Q(\xi,k)=\tau Q_1(\xi,v),
\]
and similarly for the other four functions.  Put
\[
 H_1=Q_1/k^2+2d_1/k+e_1,\qquad
 L_1=a_{01}/k^2+b_1,\qquad J_1=H_1+\sigma L_1.
\]
Angular stationarity at (HL2) gives exactly $J_1(x,k)=0$.
Consequently the successive monic differences
\[
 J_1=(\xi-x)J_\xi+\tau J_v,
 \quad
 J_\xi=\frac{J_1(\xi,v)-J_1(x,v)}{\xi-x},\qquad
 J_v=\frac{J_1(x,v)}{v-k}                            \tag{HL9}
\]
are bounded polynomials.  Let $D_1=L_1+BJ_\xi$.  The full exact
angular correction is
\begin{align*}
 f(\xi,v)-f(\xi,k)={}&
 \tau Q_1R^2+\frac{\tau a_{01}}{k^2}S^2\\
 &+\left\{2(Q_1/k+d_1)+\sigma kD_1\right\}RU\\
 &+\left\{D_1-2a_{01}/k^2\right\}ST+J_vU^2.       \tag{HL10}
\end{align*}
To check this without any division, use
$uW/\sqrt B=R+\sqrt B Z/k$ and
$\sigma W=(uZ-S)/k$.  The only term not already among the displayed
root products becomes
\[
 \tau Z^2\{BH_1+\xi L_1\}
 =D_1\tau(\xi-x)Z^2+BJ_v\tau^2Z^2
 =D_1(ST+\sigma kRU)+J_vU^2,
\]
by (HL4).  Equations (HL7) and (HL10) prove the asserted homogeneous
Gram lift for every state $W,Z$ simultaneously.  There is no state
normalization and no discarded endpoint $Z=0$.

The usual exact zero, now with the new source scale, is
\[
 T^2-\sigma U^2-ST(v-k)-\sigma kRU(v-k)=0.          \tag{HL11}
\]
This is the appropriate identity for redistributing the angular
diagonal.  Its $U^2$ coefficient is $\sigma$, which is infinitesimal;
replacing it by $1$ would erase the smaller radial scale.  Any
positivity argument must use the actual residue coefficients of
(HL7)--(HL10), rather than importing the bounded-state Schur test
without this scale correction.

\paragraph{Strict scalar corollary for the small-$B$ source.}
Suppose additionally that the coefficientwise leading functions on
the finite $u=0$ axis are
\[
 Q_0(v)=\frac{2K^2(1+v^2)}{1+K^2},\qquad
 d_0(v)=-2v,\qquad
 e_0(v)=\frac{2(v^2/K^2+K^2)}{1+K^2},
 \qquad b_0(K)>0,
\]
where $K=k_0$.  Then the $(R,U)$ residue block produced above is
\[
 2K^2R^2+\frac{4(K^2-1)}{1+K^2}RU+\frac2{K^2}U^2,
\]
with Schur margin $8/(1+K^2)^2>0$.  The $S^2$ coefficient has
residue $b_0(K)>0$, the $RS$ residue is zero, and the only other
possible residue entry is a bounded $ST$ coefficient, say $\ell$.
Add a positive real multiple $\delta$ of (HL11), with
$\delta>\ell^2/(4b_0(K))$.  Its correction to the $U^2$ entry has
zero residue because $\sigma$ is infinitesimal, whereas it supplies
the positive $T^2$ diagonal $\delta$.  Both residue blocks are then
strictly positive.  This proves the scalar strict bounded Gram for
all homogeneous states in the matched-angle chart.  Full normals
and the other testing charts remain separate requirements.


\section{\texorpdfstring{The unbounded-$b$ source with $xb$ infinitesimal}{The unbounded-b source with xb infinitesimal}}
\label{component:pab6222b899ef}

\label{pab6222b899ef:sec:c2-unbounded-b-small-B}

\begin{theorem}[The remaining unbounded-state rational C2 source]
Use the exact original rational C2 source.  Suppose $x>0$ is
infinitesimal, $k$ is a positive bounded unit, $b$ is positive
unbounded, $B=xb$ is infinitesimal, $0<r<1$, $0<h<b$, and all
original positive-chart inequalities hold.  Assume $n\geq0$ and
physicality of the source.  Then the full original parent has tensor
SOS for every independent real or complex projective state and
every independent projective pair of measurement angles.
\end{theorem}

The conclusion includes $n=0$, arbitrarily large $n$, arbitrary
complementary openings, and every value of the pure coefficient
ratio $B/\sigma$, where $\sigma=1/b$.  In particular, the actual
contact ratio $W/Z=k^{-1}\sqrt{B/\sigma}$ need not be bounded.
This proof uses the homogeneous scalar lemma in
Section~\ref{pbb29c73cc124:sec:unbounded-b-homogeneous-scalar}; it restores both
full state normals and exhausts the other testing charts.

\paragraph{Exact source scales.}
Put
\[
 \sigma=1/b,\quad x=\sigma B,\quad \eta=h/b,\quad
 y=B\eta,\quad z=B(1-\eta),\quad
 A_c=1+\sigma Br,\quad S_c=A_c+z.
\]
The positive chart gives $0<\eta<1$; both $y,z$ are infinitesimal.
Let $\widehat\Delta=\sigma\Delta$,
$\widehat N=\sigma^2 B N_{\rm src}$,
$\widehat g=\sigma g=1+\sigma(r-1+B\eta)$ and
$j=A_c-z/k^2$.  The cancelled exact formulas are
\begin{align*}
 nr&=\frac{(1+\sigma r)A_cS_c-r\widehat\Delta}
                       {2\widehat\Delta},\\
 m&=\frac{\sigma^2(1+k^2)(1-r)A_cS_c}
                       {2\widehat g\widehat\Delta},\\
 C&=\frac{A_c\widehat N}{2z\widehat\Delta},\qquad
 A=\frac{k^2A_c(1+\sigma r)S_c}{2y\widehat\Delta},\\
 D&=\frac{B(1+\sigma B)(1+\sigma r)(1-\sigma k^2r)S_c}
                       {2k^2j\widehat\Delta},\\
 \alpha&=\frac{B(1+\sigma r)\widehat N}
                       {2(1-y)\widehat\Delta}.
                                                               \tag{UB1}
\end{align*}
Here $\widehat\Delta,\widehat N,\widehat g,j,1-y$ are bounded
positive units.  Write $K=k_0>0$ and $D_K=1+K^2$.  Their relevant
residues imply
\begin{gather*}
 (nr)_0=\tfrac12(D_K^{-1}-r_0),\quad
 (m/\sigma^2)_0=(1-r_0)/2,\\
 (Cz)_0=(Ay)_0=K^2/(2D_K),\quad
 (D/B)_0=1/(2K^2D_K),\quad
 (\alpha/B)_0=K^2/(2D_K).                            \tag{UB2}
\end{gather*}
Since $n\geq0$, $r_0\leq D_K^{-1}<1$, so $1-r$ is a unit.
Set $N=4n+2\geq2$.  Exactly,
\[
 Nr=\frac{2(1+\sigma r)A_cS_c}{\widehat\Delta},      \tag{UB3}
\]
which is a bounded positive unit.  Thus $nr,nr^2$ are bounded even
when $n$ is unbounded, while $1/N$ and $1/\sqrt N$ are bounded.

All opening cancellations needed below are uniform:
\[
 BC(1-\eta)=Cz,\quad BA\eta=Ay,\quad
 BC(1-\eta)^2=Cz(1-\eta),\quad BA\eta^2=Ay\eta.
                                                               \tag{UB4}
\]
The quantity $H_B=B(C+A)$ is bounded below by a positive real;
its reciprocal and reciprocal square root are bounded.  No
assumption that $\eta$ or $1-\eta$ is a unit is made.

\paragraph{The whole finite $u=0$ axis.}
Use the original parent and the exact-contact graph
\[
 X_g=-ruZ,\qquad Y_g=\eta uW,
 \qquad E_X=\sqrt N(X+ruZ),\quad
 E_Y=\sqrt{C+A}(Y-\eta uW).                          \tag{UB5}
\]
The actual contact is
\[
 u=\sqrt{\sigma B},\quad v=k,\quad
 W=k^{-1}\sqrt{B/\sigma},\quad Z=1.
\]
Let $\xi=u^2$ and $f=F|_g$.  Exact polynomial coquotients give
\[
 f=a(\xi,v)W^2+2u d(\xi,v)WZ+c(\xi,v)Z^2,
 \quad a=\frac\xi B Q+\sigma^2a_0,\quad c=B e+\xi b.
                                                               \tag{UB6}
\]
Every coefficient of $Q,a_0,d,e,b$ is bounded by (UB3)--(UB4).
For example $a_0=4(m/\sigma^2)v^2$,
$e=4(D/B)v^2+4\alpha/B$, and
$Q=B(a-a|_{u=0})/u^2$.  These are monic polynomial divisions.
The full leading functions at $u=0$ are
\begin{align*}
 q(v)&=2K^2(1+v^2)/D_K,& d_0(v)&=-2v,\\
 e_0(v)&=2(v^2/K^2+K^2)/D_K,&
 (a_0)_0(v)&=2(1-r_0)v^2,\\
 b_0(v)&=2r_0(1+v^2)/D_K-4r_0+2.                  \tag{UB7}
\end{align*}

For remote tests use
\[
 P=uW/\sqrt B,\quad B_u=uZ,\quad
 V=\sigma vW,\quad E=\sqrt B Z.
\]
Then (UB6) has the exact bounded Gram
\[
 QP^2+b B_u^2+(a_0/v^2)V^2+eE^2+2dPE,             \tag{UB8}
\]
where $a_0/v^2$ is a polynomial coquotient.  The $(P,E)$ residue
block is strictly positive whenever $v_0^2\ne K^2$, since
\[
 q(v)e_0(v)-4v^2=4(v^2-K^2)^2/D_K^2.               \tag{UB9}
\]
This is a global identity for every finite residue $v$, not a
local-contact positivity inference.

\paragraph{Both complete state normals.}
The exact quadratic normal expansion is
\[
 F=f+E_X L_X+E_Y L_Y+H_{XX}E_X^2
                       +2H_{XY}E_XE_Y+H_{YY}E_Y^2.
                                                               \tag{UB10}
\]
Its entire hard block is bounded with residue
$\operatorname{diag}(D_{v,0},4D_{v,0})$.
Indeed the $X$ pivot is $ND_uD_v+O(1)$ before division by $N$,
with the remaining terms either infinitesimal in $u$ or of order
$m$; the $Y$ pivot from the strong packets is exactly
$4D_v(1+[A/(C+A)]u^2)$.  All mixed entries have zero residue and
contain no $C,A$ contribution.

The complete unnormalized $X$ derivative has the form
\[
 F_X|_g=uZ A_X+W(\xi B_1+\sigma^2 B_0),             \tag{UB11}
\]
with bounded polynomial coefficients, where
\begin{align*}
 B_0&=-8(m/\sigma^2)v(1+\eta\xi),\\
 B_1&=v(8D+4)(1-\eta),\\
 A_X&=-8nrD_uD_v+8m(1-r)
       -(8D+4)v^2(1+r\xi)+4D_v(1-r).
\end{align*}
Thus in the remote roots its $P,V$ coefficients have zero residue;
its only possible nonzero residue coupling is to $B_u$.
The exact residue Schur expression for the $(B_u,E_X)$ block is
\[
 (N r^2D_v-4r+2)
       -\frac{(4-2NrD_v)^2}{4ND_v}
                 =2-\frac4{ND_v}.                  \tag{UB12}
\]
It is strictly positive for every finite nonzero residue $v$,
because $N\geq2$.  At $v_0=0$ it is strict if $n$ is unbounded
or has positive residue.  The sole zero case $n_0=v_0=0$ is
treated below.

For the other complete first normal,
\[
 L_Y=\frac1{\sqrt{H_B}}
 \{P(B_2+B e_w)+B_u(\sqrt B\,u e_z)+E e_*\},       \tag{UB13}
\]
where $e_w,e_z$ are bounded,
\[
 B_2=8D_v\{(BA)\eta D_u-(BC)(1-\eta)\},\qquad
 e_*=8v(D-\alpha).
\]
All coefficients have zero residue, by (UB4), $B\to0$, and
$u\to0$.  This verifies both full normals and proves the full
parent at every remote finite-axis test covered by (UB12).

\paragraph{Matched residues without any state-ratio restriction.}
For $v_0=K$, apply the homogeneous scalar lemma to (UB6).  Its
contact and full stationarity hypotheses hold exactly because the
graph (UB5) passes through the original stationary contact.
Use the homogeneous roots
\[
 R=\frac{uW-(B/k)Z}{\sqrt B},\quad
 S=uZ-\sigma kW,\quad
 U=\sqrt B(v-k)Z,\quad T=u(v-k)Z.                  \tag{UB14}
\]
They satisfy
\[
 (\xi-x)W=u\sqrt B R+(B/k)S,\qquad
 (\xi-x)Z=uS+\sigma k\sqrt B R.                    \tag{UB15}
\]
The lemma constructs an exact bounded scalar Gram, valid for
every $W,Z$ simultaneously.  Its $(R,U)$ residue is
\[
 2K^2R^2+\frac{4(K^2-1)}{D_K}RU+\frac2{K^2}U^2,
\]
whose Schur margin is $8/D_K^2>0$.  Its radial $S^2$ residue is
$s_0=2(1-r_0)>0$; its only other possible residue coupling is
$\ell ST$, with a bounded real residue $\ell$.

To lift the full $X$ normal, set
\[
 L=A_X+(B/k)B_1+(\sigma/k)B_0.
\]
Equations (UB11) and (UB15) give exactly
\[
 F_X|_g=u\sqrt B B_1R-(\sigma/k)B_0S+uZL.
\]
At the source contact $L(x,k)=0$, because $F_X=0$ there.
Write its bounded monic differences as
$L=(\xi-x)L_\xi+(v-k)L_v$.  Then
\[
 F_X|_g=
 u\sqrt B(B_1+\sigma kL_\xi)R
 +\{\xi L_\xi-(\sigma/k)B_0\}S+L_vT.             \tag{UB16}
\]
After division by $\sqrt N$, all coefficients are bounded.
The $R,S$ coefficients have zero residue; the $T$ coefficient
may have a finite residue $\gamma$ (it is
$-8K/(D_K\sqrt{N_0})$ when $N$ is bounded).
This nonzero full-normal coupling is retained, not silently
discarded.

For (UB13), put $A_Y=B_2+B e_w$ and
$L_Y'=A_Y/k+\xi e_z+e_*$.  Its numerator is exactly
$RA_Y+\sqrt B ZL_Y'$.  Exact normal contact gives
$L_Y'(x,k)=0$.  Monic differences and (UB15) yield bounded
coefficients on $R,S,U$, all of zero residue: the $R,S$ corrections
carry factors $\sigma B$ or $u\sqrt B$, and the $U$ coefficient
is the angular coquotient of an identically zero leading function.
Division by $\sqrt{H_B}$ preserves these bounds.

Finally add $\delta$ times the exact zero
\[
       T^2-\sigma U^2-ST(v-k)-\sigma kRU(v-k)=0.    \tag{UB17}
\]
Choose a real $\delta>\ell^2/(4s_0)+\gamma^2/(4D_K)$.
This supplies a $T^2$ residue $\delta$ at zero residue cost to
the $U^2$ coefficient, since $\sigma$ is infinitesimal.  The block
on $(S,E_X,T)$ has strictly positive Schur complement; the $(R,U)$
block and $E_Y$ diagonal are already strict.  Thus the complete
six-root Gram on $(R,S,U,T,E_X,E_Y)$ is strictly positive in
residue.  The calculation is homogeneous in all four state
coordinates and needs no bound on $W/Z$, no comparison of $B$ with
$\sigma$, and no state endpoint exclusion.  The exact simultaneous
sign copy treats $v_0=-K$.

\paragraph{The double-zero test, including literal $n=0$.}
Suppose now $n_0=v_0=0$.  Use the pure coefficient quantities
\[
 H_c=4n+4m+2,\quad \zeta=(4m+2)/H_c,\quad
 \xi_c=4m/H_c,\quad \eta_c=C/(C+A),
\]
and graph $X_g=\xi_cvW-\zeta uZ$, $Y_g=\eta_cuW$.
Normalize its $X$ normal by $\sqrt{H_c}$ and its $Y$ normal by
$\sqrt{C+A}$.  Put
\[
 \epsilon=4\alpha+4n(4m+2)/H_c>0.
\]
Then $\epsilon$ is infinitesimal, $n/\epsilon$ and $B/\epsilon$
are bounded, and the harmonic product $BCA/(C+A)$ is a unit.
Use the five scalar roots
\[
 P=uW/\sqrt B,\quad V=\sigma vW,\quad E=\sqrt B Z,
 \quad B_\epsilon=\sqrt\epsilon\,uZ,\quad T_0=uvZ.
\]
The exact inherited identity, with its original coefficients,
is
\[
 f_{zz}=4\alpha+\epsilon u^2+4n\zeta^2u^4
                       +4Dv^2+u^2v^2 B_b,
 \qquad (B_b)_0=2.
\]
The $W^2$ coefficient at $u=0$ is exactly
\[
 \frac{4m(2n+1)}{2n+2m+1}v^2
              +(4n+2)\xi_c^2v^4.
\]
Therefore its coefficient on $V^2$ has positive residue
$4(m/\sigma^2)_0=2(1-r_0)$.  The scalar residue on the five roots
is strictly diagonal, with entries
\[
 \left(\frac{2K^2}{D_K},\ 2(1-r_0),\
        \frac{2K^2}{D_K},\ 1,\ 2\right).
\]
The cross term is allocated using $uvWZ=vPE$, so it has zero residue.

Both full normals are restored, including when $n=0$ exactly.
At $v=0$ the $X$ derivative is $-8n\zeta u^3Z$, whose coefficient
on $B_\epsilon$ is infinitesimal since $n/\epsilon$ is bounded.
Its remaining $Z$ coefficient is $uv^2$ times a bounded function,
so it has zero residue on $T_0$.  Its $W$ coefficient is a sum of
$u^2vW$ times a bounded function and $mvW$ times a bounded
function divisible by $v^2$ at $u=0$.  Allocate the former to
$P$ with coefficient $u\sqrt B\,v$, and the latter to $V$ with
coefficient $m/\sigma$; both have zero residue.  In particular one
must not divide the former by $\sigma$.
The $Y$ derivative is the inherited harmonic-opening identity:
$B_2=8D_v(BA)\eta_cu^2$, with bounded residual coquotients.
Its coefficients on $P,T_0,E$ all have zero residue.  The hard
residue is $\operatorname{diag}(1,4)$.  This proves the full-parent
corner without deleting a state coordinate.

\paragraph{All remaining projective tests.}
For the other commuting charts use the four original outer-core
identities with $\rho$ replaced by $\sqrt B$, and cap $C,A$ at
$s/B$, where $s=\min(1,BC,BA)$ is a positive bounded unit.
Retain both original nonnegative excess packets and the entire
original $n$ packet.  No positive lower bound on $m$ is needed:
the first two outer pivots are $4m_0+2=2$, while the other two
charts retain the whole $m$ packet and do not use it for strictness.
The reciprocal-$v$ chart uses $(D/B)_0>0$ exactly as before.
Thus the finite unit $u$ axis, both reciprocal axes, all double
corners, and literal projective endpoints have strict bounded
Grams.  The inherited physical noncommuting-residue theorem covers
all remaining testing pairs.

The finite-axis certificates above are homogeneous in the state,
including the previously problematic unbounded actual state
ratio.  Splitting a complex state into its real and imaginary
parts proves the complex assertion without an additional chart.
All inverses and radicals used in coefficients are pure coefficient
or pure testing operations; every displayed mixed quotient is
explicit monic polynomial division.  The roots may be unbounded.
The strict bounded tensor Gram lemma and restoration of the retained
packets complete the all-tests proof.


\section{\texorpdfstring{The exact unbounded-state-scale C2 complement}{The exact unbounded-state-scale C2 complement}}
\label{component:pb60046a37191}

\label{pb60046a37191:sec:c2-unbounded-b-reduction}

Use the original positive rational source with $d=1$, actual
$u_c=\sqrt x$ positive infinitesimal and $v_c=k$ a positive bounded
unit. The global entry theorem supplies its positive chart inequalities.
The preceding rational theorems, including the joint-opening theorem,
now cover every bounded $b$: the inherited $n\ge0$ entry inequality
makes bounded $b$ a unit. We classify the remaining $b$ unbounded
without changing any original coefficient.

Put
\[
 B=xb,\qquad H=xh,\qquad L=x(b-h)=B-H,\qquad
 a^2=k^2/x,\quad t=L/k^2.
\]
The original gaps give
\[
 0<H<1,\qquad0<L<k^2(1+xr),\qquad B=H+L.
\]
Thus $B,H,L$ are bounded and $x/B=1/b$ is infinitesimal.
Set $A_c=1+xr$ and $S=A_c+L$, both positive bounded units. Also
\[
 g/b=1+(r-1)/b+H/b\longrightarrow1,
 \qquad j=1+xr-L/k^2>0,\quad j\le1+x.
\]
The following exact rearrangements of the original short polynomials
control every event scale:
\begin{align*}
 x^2N={}&(1+x)(B+xr)(H+k^2)-x(1+k^2)(1-r)L,\\
 \Delta/b={}&(1+x)B+2x(1+x)r+(1+x)(1+k^2)\\
 &+(x^2-k^2x)r^2/b+2x(1+k^2)r/b+(1+k^2)/b.
\end{align*}
Consequently $xN/b=x^2N/B$ and $\Delta/b$ are positive bounded
units, with respective residues
\[
 H_0+k_0^2,\qquad B_0+1+k_0^2.
\]
Here $H,L\le B$ controls the divided corrections, and $x/B\to0$.

The exact physical condition $n\ge0$ is
\[
 r(\Delta/b)\le(1+r/b)A_cS.
\]
It implies
\[
 r_0\le\frac{1+L_0}{1+k_0^2+B_0},\qquad
 1-r_0\ge\frac{k_0^2+H_0}{1+k_0^2+B_0}>0.
\]
Thus $1-r$ is a positive bounded unit without any lower bound on $r$.
The original formulas now give the exact comparabilities
\[
 m\asymp b^{-2},\qquad C\asymp L^{-1},\qquad A\asymp H^{-1},
 \qquad D\asymp B/j,\qquad\alpha\asymp B/(1-H).
 \tag{SR1}
\]
Every omitted factor in $\asymp$ is a positive bounded unit with
bounded reciprocal. For instance,
\[
 CL=\frac{A_c(xN/b)}{2\Delta/b},\qquad
 AH=\frac{k^2A_c(1+r/b)S}{2\Delta/b}
\]
are positive units; the other identities follow using $g/b\to1$
and $(b-k^2r)/b\to1$.

There are exactly two unbounded-$b$ scale leaves:
\begin{enumerate}
\item If $B=xb$ is infinitesimal, so are $H,L$, while $j\to1$
and $1-H\to1$. Thus $C,A$ are unbounded and $m,D,\alpha$ are
weak, with $m\asymp b^{-2}$ and $D\asymp\alpha\asymp xb$.
This is the one-weak-complementary-event frontier, with weak
$\alpha$ and strong $C$ in this actual-angle orientation.
\item If $B$ is a positive bounded unit, then $C,A,D,\alpha$
are all strong and $m$ is weak. The new strong-complement pitchfork
applies to this entire leaf at all upper event scales. In particular
it includes $C$ unbounded when $L$ is infinitesimal and $\alpha$
unbounded when $1-H$ is infinitesimal.
\end{enumerate}
Unbounded $B$ is impossible by the original chart inequalities.
The homogeneous all-tests theorem of Section~\ref{pab6222b899ef:sec:c2-unbounded-b-small-B}
closes the remaining infinitesimal-$B$ leaf without assuming a bounded
actual state ratio. Thus both unbounded-$b$ leaves, as well as all
bounded-$b$ openings, are now covered. Literal $n=0$ is preserved.


\section{\texorpdfstring{Physical entry excludes a small state parameter}{Physical entry excludes a small state parameter}}
\label{component:p7294aae40ee2}

\label{p7294aae40ee2:sec:c2-bounded-state-entry}

This short reduction concerns the same exact rational source, before
assuming that $b$ is a unit.  Suppose $x>0$ is infinitesimal, $k$ is
a positive bounded unit, $b>0$ is bounded, $0<r<1$, $0<h<b$,
$g=b-1+r+xh>0$, and $n\geq0$.  No further source physicality is
needed for the following necessary bound.

All quantities in
\[
             (b+r)A_cS_c\geq r\Delta
\]
are bounded, so taking residues gives
\[
 r_0\{b_0+k_0^2(b_0+1)\}\leq b_0.
\]
The brace is strictly positive even if $b_0=0$.  On the other hand
$g>0$ gives $r_0\geq1-b_0$.  If $b_0\leq1$, combining the two
inequalities yields
\[
 (1-b_0)\{b_0+k_0^2(b_0+1)\}\leq b_0,
 \quad\text{or equivalently}\quad
              b_0^2\geq\frac{k_0^2}{1+k_0^2}.
\]
For $b_0\geq1$ this lower bound is automatic.  Hence in every case
\[
                b_0\geq\frac{k_0}{\sqrt{1+k_0^2}}>0.
\]
Thus bounded $b$ is necessarily a unit.  In particular there is no
additional $b$-infinitesimal source domain in this actual-contact
scaling.  Combining the two new single-angle certificates, the
uncaptured part of this bounded-$b$ entry domain has exactly the
joint necessary degeneration $r_0=g_0=0$, hence $b_0=1$; the other
remaining scale is unbounded $b$.  This is a necessary source
reduction, not a certificate for the joint degeneration.


\section{\texorpdfstring{The remaining opening when the opposite event is bounded}{The remaining opening when the opposite event is bounded}}
\label{component:pd7bd7f0d7fb6}

\label{pd7bd7f0d7fb6:sec:c2-last-opening}

\begin{theorem}[The unbounded-$m$ opening with bounded $n$]
\label{pd7bd7f0d7fb6:thm:c2-no-opening-bounds}
In Theorem~\ref{p81893ad08bed:thm:c2-three-openings}, replace the assumption that
$g=b-1+r+xh$ is a positive bounded unit by $g>0$ and the assumption that at least one of $r,g$ is a bounded
unit.  The all-tests tensor-SOS conclusion still holds.
\end{theorem}

The only new case is infinitesimal $g$; it is automatically bounded.
Our hypothesis now makes $r$ a unit.  Since $b_0+r_0=1$ and $b$ is
a positive unit, $r_0<1$.  Thus $n$ is bounded,
$0<b_0<1$, and $1-r$ is a unit.  The source formulas imply
\[
               mg=\frac{(1+k^2)(1-r)A_cS_c}{2\Delta},   \tag{CG1}
\]
which is a positive bounded unit.  Thus $m$ is unbounded.  The
two-opening bounds (CO1), positivity of $D/x,\alpha/x$, and the
consequence that $b-k^2r$ is a unit remain valid.

\paragraph{Use the original strong amplitude as the hard normal.}
Put $\eta=h/b$ and use the exact contact graph
\[
 Y_g=\eta uW,\qquad
 X_g=v(1+\eta u^2)W-(1+g)uZ.                            \tag{CG2}
\]
At contact this gives $X_g=\rho[b+xh-1-g]=-\rho r$.
The original $m$ amplitude satisfies exactly $M|_g=guZ$.
Normalize the full hard directions by
\[
 E_M=\sqrt m\{X-X_g-uv(Y-Y_g)\},\qquad
 E_Y=\sqrt{C+A}(Y-Y_g).                                 \tag{CG3}
\]
In these variables
\[
 M=guZ-E_M/\sqrt m,
 \quad
 X=X_g+E_M/\sqrt m+uvE_Y/\sqrt{C+A}.
\]
Consequently the unbounded packet has no $E_Y$ contribution at all.
This triangular change is essential: scaling $X$ alone would leave
the uncontrolled coefficient $\sqrt{m/(C+A)}uv$.
The entire scalar contribution of that packet is $4mg^2u^2Z^2$,
which is bounded by (CG1).  Its full $X$ derivative contributes
$-8mg\,uZ$, also bounded, before division by $\sqrt m$.

All other source weights entering the $X$ row are $n,D$ and the
unit signed weight, so the full $X$ derivative on (CG2) has bounded
coefficients on $V=vW,B=uZ$.  Hence the $E_M$ first normal has
zero-residue bounded coefficients on those roots.  The $E_Y$ first
normal is
\[
       \frac{(F_Y+uvF_X)|_g}{\sqrt{C+A}}.               \tag{CG4}
\]
The $m$ packet cancels identically in $F_Y+uvF_X$, before any graph
substitution.  The remaining large terms are exactly the $C,A$
terms controlled by (CO1); all other coefficients are bounded.
In particular the old (CO5) form on $P,B,E$ holds with bounded
coquotients and identically zero leading corrections.
The complete quadratic hard block is bounded and has residue
\[
                         \operatorname{diag}(4,4D_{v,0})  \tag{CG5}
\]
on the finite $u=0$ axis.  The $m$ packet gives precisely its first
diagonal entry, the $C,A$ packets give the second, and all mixed
entries from bounded packets contain $1/(\sqrt m\sqrt{C+A})$.

\paragraph{The remaining five Laurent terms already force strictness.}
The exact scalar divisions on $P,B,V,E$ are bounded for the same
reason: $C,A$ enter through (CO1), and $m$ only through $mg^2$.
Their leading weighted scalar is
\[
 q\,u^2W^2/x+b_su^2Z^2+c_sW^2-d_suWZ+exZ^2,
\]
where, for $z=v^2$ and $\nu=2n_0+1\geq1$,
\[
 \begin{gathered}
 q=Q(1+z),\quad b_s=2\nu(1+z)-2,\quad
 c_s=2z\nu(1+z),\quad d_s=4v\nu(1+z),\\
 e=E_2z+E_0,\qquad Q=4T_0/b_0,
 \quad E_2=4(D/x)_0>0,\quad E_0=4(\alpha/x)_0>0.
 \end{gathered}                                         \tag{CG6}
\]
Set $\mathcal D(s,z)=(qs+c_s)(b_ss+e)-d_s^2s/4$.
Its exact Laurent identity is
\begin{align*}
 \frac{\mathcal D(s,z)}{(1+z)sz}={}&
 2Q(\nu-1)s/z+2Q\nu s
 +2\nu E_2z/s+2\nu E_0/s\\
 &+QE_0/z+QE_2-4\nu.                                  \tag{CG7}
\end{align*}
Contact and stationarity give a zero at $(s,z)=(1,k_0^2)$.
After $z=k_0^2t$, write the five coefficients in order as
$a,b,c,d,e$.  Here $a\geq0$ and $b,c,d,e>0$.  Stationarity gives
$c=a+e$ and $b=d+e$.  The zero value fixes the constant, so exactly
\[
 \mathcal L=
 a(s/t+t/s-2)+d(s+1/s-2)+e(t/s+s+1/t-3).                 \tag{CG8}
\]
Every summand is nonnegative by AM--GM.  The last bracket alone
forces $s=t=1$ at equality: its three positive terms have product
one, and equality makes all three equal.  Its exponent vectors
$(-1,1),(1,0),(0,-1)$ span two dimensions; equivalently its Hessian
in $\log s,\log t$ at $(1,1)$ is
$\left(\begin{smallmatrix}2&-1\\-1&2\end{smallmatrix}\right)$.
Hence it already has a positive definite contact Hessian.
This proves global uniqueness and local nondegeneracy even when
$n_0=0$ and the first pair is absent.

For remote finite nonzero $v$ the same scalar cross-term split as
before gives a strict Gram.  The case $v_0=0,n_0>0$ is also strict
on $P,B,V,E$.  The full normals are restored by (CG3)--(CG5).
At matched finite-angle residues the whole scalar still has the
form $p^2Q_b/x+B_b$ with bounded $Q_b,B_b$.  The four monic
conductors $R,S,U,T$ of the prior theorem therefore give a bounded
Gram, and (CG8) proves its strict Schur inequality.

The full first normals have the required same conductor lifts.
For $E_M$, their coefficient is $(uA_1+WB_1)/\sqrt m$ with bounded
polynomial $A_1,B_1$.  Monic contact division retains the factor
$1/\sqrt m$ in every coquotient.  Exact contact makes the residual
coefficient $A_1+\kappa B_1$ vanish, giving the old four-root lift
with zero residue.  For $E_Y$, the cancellation in (CG4) leaves
precisely the old large-packet term $B_2$ plus bounded residual
terms; its exact contact and identically zero leading part give
the same $R B_2+\rho E_2'$ division, then the bounded factor
$1/\sqrt H$.  These are full normal identities, not just Hessian
checks.  Quadratic state degree and (CG5) finish the complete
matched-contact certificate.

\paragraph{The double-zero test when $n_0=0$.}
The graph (CG2) is replaced in this remote corner.  Put
\[
 H_c=4m+4n+2,\quad \xi=4m/H_c,\quad
 \zeta=(4m+2)/H_c,\quad \eta_c=C/(C+A).
\]
Use
\[
 Y_g=\eta_cuW,\quad
 X_g=\xi v(1+\eta_cu^2)W-\zeta uZ,
\]
and the triangular normals
\[
 E_X=\sqrt{H_c}\{X-X_g-uv(Y-Y_g)\},\qquad
 E_Y=\sqrt{C+A}(Y-Y_g).                                \tag{CG9}
\]
The entire graph amplitude is
\[
 M|_g=(1-\xi)v(1+\eta_cu^2)W+(\zeta-1)uZ.
\]
Both $m(1-\xi)$ and $m(\zeta-1)$ are bounded because $n$ is
bounded.  Thus the scalar and full first normals have bounded
coquotients after the same $x$ scaling as before.

Set $\epsilon=4\alpha+4n\zeta>0$.  It is infinitesimal, with
$n/\epsilon,x/\epsilon$ bounded.  The complete scalar $Z^2$
coefficient is exactly
\[
 f_{zz}=4\alpha+\epsilon u^2+4n\zeta^2u^4+4Dv^2
                            +u^2v^2 B_b,               \tag{CG10}
\]
where $B_b$ is bounded with residue $2$.  This identity follows
directly from the packets and the exact relation
$4n\zeta^2+(4m+2)(1-\zeta)^2=4n\zeta$.
The $W^2$ coefficients on $P=uW/\rho,V=vW$ have positive residues
$Q,2$.  The cross term is $uvWZ$ times a bounded function and
can be written as a zero-residue multiple of $PE$ using
$uvWZ=vPE$.  Thus the five scalar roots
\[
           P,\quad V,\quad E=\rho Z,\quad
           B_\epsilon=\sqrt\epsilon\,uZ,\quad T=uvZ
\]
have strict diagonal residue $(Q,2,E_0,1,2)$.

For the full $X$ derivative, at $v=0$ the $Z$ coefficient is exactly
$-8n\zeta u^3$.  After division by $\sqrt{H_c}$ its coefficient
on $B_\epsilon$ is infinitesimal because
$n/\sqrt\epsilon=(n/\epsilon)\sqrt\epsilon$ is infinitesimal.
The remaining $Z$ coefficient has the factor $uv^2$ and bounded
coquotient; its coefficient on $T$ has zero residue.  The $W$
coefficient has a factor $v$ and bounded coquotient, multiplied by
$1/\sqrt{H_c}$.  For $E_Y$ the exact $m$ cancellation in
$F_Y+uvF_X$ again applies; (CO6) gives
$B_2=8D_v(xA)\eta_c u^2$, and the remaining lift is
\[
 P(B_2+xe_w)+T(\rho u e_z)+E\,8v(D-\alpha),
\]
divided by $\sqrt H$, with bounded $e_w,e_z$.
Every coefficient has zero residue.  The hard residue of (CG9) is
$\operatorname{diag}(1,4)$, and all mixed hard entries vanish in
residue.  This proves the full-parent corner including literal $n=0$.

\paragraph{All other tests.}
For the remaining commuting axes use the earlier outer identities,
cap $C,A$ at $\sigma/x$ as before, and also cap $m$ at
$\min(m,1)$.  Retain the entire nonnegative excess $m$ packet.
The selected $m$ weight is a unit and all previous displayed
strict residues apply.  The noncommuting-residue theorem covers
the other tests.  Restore all positive factors and split complex
states as before.  This proves the theorem for arbitrary positive
$g$, including the new unbounded-$m$ regime.

The remaining assumptions in this source family are now only that
$k,b$ are units, $x$ is positive infinitesimal, all original contact
openings are positive, at least one of $r,g$ is a unit, $n\geq0$,
and the source is physical.  The joint degeneration $r,g\to0$,
where both $n,m$ can be unbounded, is not covered.  The
result does not yet treat unbounded $b$ or a degenerate residue of
the finite actual angle $k$.  Section~\ref{p7294aae40ee2:sec:c2-bounded-state-entry}
excludes infinitesimal $b$ under these entry hypotheses.


\section{\texorpdfstring{The one-weak C2 family with three arbitrarily small openings}{The one-weak C2 family with three arbitrarily small openings}}
\label{component:p81893ad08bed}

\label{p81893ad08bed:sec:c2-three-openings}

This theorem removes three unit assumptions from the source of
Theorem~\ref{pec59340d644f:thm:c2-general-single-angle}.  Both positive openings $h$
and $b-h$ may approach zero, and $r$ may approach zero, making $n$
arbitrarily large.  The conclusions concern the entire original parent
and every independent test.  Source physicality remains an explicit
hypothesis; no new source classification is claimed.

\begin{theorem}[Three-opening single-angle source]
\label{p81893ad08bed:thm:c2-three-openings}
Use exactly the rational C2 source and contact of
Section~\ref{pec59340d644f:sec:c2-general-single-angle}.  Suppose $x>0$ is
infinitesimal, $k,b$ are positive bounded units,
\[
             0<r<1,\qquad 0<h<b,
\]
and $g=b-1+r+xh$ is a positive bounded unit.  Assume $n\geq0$
and physicality of the original source.  There is no lower bound on
$r,h,b-h$.  Then the entire original parent has tensor SOS at all
independent projective measurement pairs and all real or complex
projective states.  Literal $n=0$ is included.
\end{theorem}

We write $\rho=\sqrt x$, $\eta=h/b$, and retain the exact source
quantities $A_c,S_c,\Delta$ from that section.  All arguments below
take place over the independent coefficient and testing fields.
Every coefficient inverse used below is pure coefficient or pure
testing, as specified; no mixed denominator is inverted.

\paragraph{Entry and the two complementary openings.}
First the former assumptions that $1-r$ and $b-k^2r$ are units follow
from $n\geq0$.  Indeed all quantities in
\[
                 (b+r)A_cS_c\geq r\Delta
\]
are bounded, and $\Delta_0=(1+k_0^2)(b_0+1)$.  Taking residues yields
\[
 r_0\leq\frac{b_0}{b_0+k_0^2(b_0+1)}<1,
 \qquad
 b_0-k_0^2r_0\geq
 \frac{b_0^2(1+k_0^2)}{b_0+k_0^2(b_0+1)}>0.
\]
Thus $m$ is a positive bounded unit, and $D/x,\alpha/x$ are positive
bounded units, regardless of whether $r$ has zero residue.

Set
\[
 T_C=xC(b-h),\qquad T_A=xAh,\qquad H=x(C+A).
\]
The original source formulas, with their apparent opening denominators
cancelled, show that $T_C,T_A$ are bounded positive units and have the
same residue
\[
 T_0=\frac{k_0^2(b_0+r_0)}{2\Delta_0}>0.
\]
Consequently $H$ is bounded below by a positive real, although it may
be unbounded.  The four identities
\begin{align*}
 xC(1-\eta)&=T_C/b,&xA\eta&=T_A/b,\\
 xC(1-\eta)^2&=T_C(b-h)/b^2,&
 xA\eta^2&=T_Ah/b^2                                      \tag{CO1}
\end{align*}
are the reason no new scalar singularity appears when either opening
collapses.  In particular their last two terms add to a bounded unit,
with residue $T_0/b_0$.  The effective scalar coefficient is
\[
                         Q=4T_0/b_0>0.                  \tag{CO2}
\]
It is independent of the limiting value of $h$.

\paragraph{A bounded $n$: restoration with the normalized hard root.}
If $r$ is a unit, the exact formula for $n$ makes $n$ bounded.  Use the
contact graph $X_g,Y_g=\eta uW$ of the prior theorem.  All of its scalar
coefficients remain bounded: the only potentially unbounded weights
are $C,A$, and on the graph they occur solely through the last two
products in (CO1), divided by $x$ and multiplied by $u^2W^2$.
Its leading scalar is exactly the scalar in that theorem with $Q$
replaced by (CO2).  The quantities $n_0,m_0,D_*,E_*$ have their former
values.  Hence the entire six-term Laurent identity, its balanced
AM--GM decomposition, global equality characterization, and positive
contact Hessian apply without change.

Replace the old hard root $N_Y=(Y-Y_g)/\rho$ by
\[
        \widehat N_Y=\sqrt H\,N_Y=\sqrt{C+A}(Y-Y_g).
                                                               \tag{CO3}
\]
Its exact hard coefficient from the $C,A$ packets is
\[
             4D_v\frac{C+AD_u}{C+A}
       =4D_v\left(1+\frac{A}{C+A}u^2\right),             \tag{CO4}
\]
which is bounded and has residue $4D_{v,0}$ on the finite $u=0$
residue axis.  The other packet contributions to this coefficient
and to the mixed hard entry are bounded: $1/(C+A)=x/H$ and
$1/\sqrt{C+A}=\rho/\sqrt H$ are bounded and infinitesimal.

For clarity the complete first-normal verification is also retained.
The previous unnormalized $Y$ derivative has exactly the form
\[
 L_Y=P(B_2+xe_w)+B(\rho u e_z)+E e_*,\qquad
 B_2=8D_v\{(xA)\eta D_u-(xC)(1-\eta)\},
 \quad e_*=8v(D-\alpha).                                \tag{CO5}
\]
Here $P=uW/\rho$, $B=uZ$, $E=\rho Z$.  The functions $e_w,e_z$
come only from the bounded $n,m,D,\alpha$ packets and the signed
core, so remain bounded.  Equation (CO1) makes $B_2$ bounded with
identically zero residue for every finite $v$ on $u=0$; $e_*$ has
the same property.  The coefficient of $\widehat N_Y$ is exactly
$L_Y/\sqrt H$, so its bounded zero-residue lifts are preserved.
The $X$ derivative contains neither $C$ nor $A$ and is unchanged.
Thus all four-root remote certificates restore both full normals.

For matched finite-angle residues the scalar has exactly the old
form $f=p^2Q_b/x+B_b$, with
\[
 Q_b=4D_v\{(xC)(1-\eta)^2+(xA)\eta^2D_u\}.
\]
Both $Q_b$ and $B_b$ have bounded coefficients by (CO1).  The monic
divisions in $u^2-x$, $W^2-\kappa^2x$, $uW-\kappa x$, and $v-k$,
where $\kappa=b/k$, therefore give precisely the same bounded Gram
on $R,S,U,T$.  Its residue is unchanged.  The old first-normal
contact divisions use only bounded $B_2,e_w,e_z$ and exact stationary
contact; multiplication by the pure coefficient $1/\sqrt H$
preserves their zero-residue lifts.  Hence the complete six-root
Gram has strictly positive residue.

If $n_0=0$ at the double-zero test, use the previous exact graph
with $\eta_c=C/(C+A)$.  Now
\[
 xC(1-\eta_c)=xA\eta_c=\frac{xCA}{C+A},
 \qquad
 \left(\frac{xCA}{C+A}\right)_0=T_0/b_0.                \tag{CO6}
\]
These quantities are bounded units even when $H$ is unbounded.
Thus the exact $\epsilon$ reserve and five scalar roots in that
proof remain valid.  Its $Y$ normal is again divided by $\sqrt H$,
and its hard $Y$ residue is now $4$ rather than $4(C_*+A_*)$.
This proves the literal $n=0$ case as well as its infinitesimal
neighborhood without deleting a state coordinate.

\paragraph{An infinitesimal $r$: the opposite event can be unbounded.}
Suppose now $r$ is infinitesimal.  Let $N=4n+2$.  The exact source
formula gives
\[
             Nr=\frac{2(b+r)A_cS_c}{\Delta}.             \tag{CO7}
\]
Thus $Nr$ is a bounded positive unit, $N$ is unbounded, and both
$nr$ and $nr^2$ are bounded.  Use the exact-contact graph
\[
 X_g=-ruZ,\qquad Y_g=\eta uW,
 \qquad
 \widehat N_X=\sqrt N(X+ruZ),\quad
 \widehat N_Y=\sqrt{C+A}(Y-\eta uW).                    \tag{CO8}
\]
At the original contact $u=\rho,v=k,W=\rho b/k,Z=1$ this graph
has exactly the required $X,Y$ values.  The potentially large $n$
packet contributes $4nr^2D_uD_vu^2Z^2$ to the graph scalar, and
$-8nrD_uD_vuZ$ to its unscaled $X$ derivative.  Equation (CO7)
therefore controls the entire scalar and every full first normal,
not only their values at contact.

The hard residue on $(\widehat N_X,\widehat N_Y)$ is
\[
                  \operatorname{diag}(D_{v,0},4D_{v,0}). \tag{CO9}
\]
Indeed $H_{XX}/N=D_uD_v+O(1/N)$, since the remaining source
weights that enter $H_{XX}$ are $m,D$; (CO4) controls $H_{YY}$.
The mixed entry is bounded times
$1/(\sqrt N\sqrt{C+A})$ and has zero residue.

The five exact scalar monic divisions on $P,B,V,E$ from the previous
proof still apply.  Their residues, with $z=v^2$, are now
\[
 q=Q(1+z),\quad b_s=2\mu,\quad c_s=2(\mu-1)z,
 \quad d_s=4\mu v,\quad e=E_2z+E_0,                    \tag{CO10}
\]
where $\mu=2m_0+1>1$, $E_2=4(D/x)_0>0$,
$E_0=4(\alpha/x)_0>0$, and $Q$ is given by (CO2).
Define
$\mathcal D(s,z)=(qs+c_s)(b_ss+e)-d_s^2s/4$.
Exactly,
\begin{align*}
 \frac{\mathcal D(s,z)}{sz}={}&
 2Q\mu s/z+2Q\mu s
 +2(\mu-1)E_2z/s+2(\mu-1)E_0/s\\
 &+QE_0/z+QE_2z+Q(E_0+E_2)-4\mu.                    \tag{CO11}
\end{align*}
All six nonconstant Laurent coefficients are positive.  Exact
stationary contact, applied after the bounded graph substitution
$u=\rho s_1,W=\rho w_1$, gives a stationary zero at
$(s,z)=(1,k_0^2)$.  Alternatively this follows by direct substitution
of
\[
 m_0=\frac1{2(b_0-1)(b_0+1)},\quad
 Q=\frac{2k_0^2}{\Delta_0},\quad
 E_2=\frac{2b_0^2}{k_0^2\Delta_0},\quad
 E_0=\frac{2k_0^2b_0^2}{\Delta_0}.
\]
Here $b_0>1$ follows from $g_0=b_0-1>0$.
Rescale $z=k_0^2t$ and use exactly the balanced reciprocal-pair and
three-cycle identity of the previous theorem.  Its $s+1/s-2$ and
$t+1/t-2$ coefficients are strictly positive.  Thus (CO11) has a
unique zero in the positive quadrant and a positive definite Hessian
there.  This proves remote scalar strictness and the matched-contact
Schur inequality, with no appeal to numerical evidence.

For the full $X$ normal, direct differentiation and parity give
\[
 \left.F_X\right|_g=V A_X+B B_X,
\]
where the complete coefficients are
\begin{align*}
 A_X={}&-8m(1+\eta u^2)+(8D+4)(1-\eta)u^2,\\
 B_X={}&-8nrD_uD_v+8m(1-r)
       -(8D+4)v^2(1+ru^2)+4D_v(1-r).
\end{align*}
These are bounded polynomials in the bounded testing coordinates.
Consequently the coefficient of $\widehat N_X$ has bounded
zero-residue coefficients on $V,B$ because $1/\sqrt N$ is
infinitesimal.  The full $Y$ normal is (CO5), with the bounded
graph (CO8); (CO1) gives all of its required bounds and zero residues.
Its complete residual coefficients are
\begin{align*}
 e_w={}&8mv^2(1+\eta u^2)-(8D+4)(1-\eta)
                     +(8\alpha+4)D_uv^2\eta,\\
 e_z={}&v\{-8m(1-r)+(8D+4)r-(8\alpha+4)\}.
\end{align*}
This restores the remote scalar certificates, including $v_0=0$,
where (CO10) has strictly positive coefficients on $P,B,V,E$ and
zero cross residues.  No special $n_0=0$ corner occurs in this case.

At a matched residue write, exactly,
\[
 \frac{F_X|_g}{\sqrt N}=uA_1+WB_1
   =-\kappa B_1S+u(A_1+\kappa B_1).
\]
Both $A_1,B_1$ have identically zero residue for every finite $v$;
their finite monic coquotients also retain the factor $1/\sqrt N$.
Exact normal contact makes $A_1+\kappa B_1$ vanish at the source
contact.  The same four conductors $R,S,U,T$ therefore give bounded
zero-residue coefficients, exactly as in the bounded-$n$ argument.
For the $Y$ normal the previous division
$L_Y=RB_2+\rho E_2'$ has bounded $B_2,E_2'$ with identically zero
leading functions; division by $\sqrt H$ preserves those properties.
All quadratic normal coefficients are already bounded by
(CO4),(CO7),(CO9).  Thus the matched-contact Gram for the entire
parent is strictly positive in residue.  The other state charts and
the simultaneous sign copy use precisely the previous proof.

\paragraph{Every other commuting test.}
Choose the pure coefficient
\[
                         \sigma=\min(1,xC,xA).
\]
It is a positive bounded unit, since $T_C,T_A$ are units while
$h,b-h$ are bounded.  Retain the nonnegative excess packets
\[
 4(C-\sigma/x)D_v(uW-Y)^2,
 \qquad4(A-\sigma/x)D_uD_vY^2,
\]
and use the four outer core identities in the previous theorem with
$C=A=\sigma/x$.  These identities retain the entire $n$ packet;
they never use boundedness of $n$.  Their displayed hard entries
$4C_*,4A_*$ are simply replaced by $4\sigma_0$, while $m_0>0$
and $(D/x)_0>0$ are unchanged.  Every residue Gram stays strictly
positive.  This covers the finite-unit $u$ axis, both reciprocal
axes and their double corners, including literal projective ends.

The finite $u=0$ axis, those outer axes, and the inherited theorem for
noncommuting-residue tests exhaust all independent measurements.
Restore the retained packets and positive homogeneous factors, and
split complex states into real and imaginary parts.  This proves
Theorem~\ref{p81893ad08bed:thm:c2-three-openings}.  In particular the newly admitted
source scales permit either $A$ or $C$ to be arbitrarily larger than
$x^{-1}$, independently of an arbitrarily large $n$.  They cannot
both have unbounded ratios to this scale because $h+(b-h)=b$ is a unit.  There is no
prescribed comparison between these coefficient scales and any
testing infinitesimal.


\section{\texorpdfstring{C2 pitchfork with three arbitrary event weights}{C2 pitchfork with three arbitrary event weights}}
\label{component:p6c62a53b1454}

\label{p6c62a53b1454:sec:c2-three-large}

\begin{theorem}[All scales of the three retained events]
\label{p6c62a53b1454:thm:c2-three-large}
Let $K_c,K_t$ be independent real closed extensions of $\mathbb R$.
Let a physical canonical C2 source over $K_c$ have $\lambda=\beta=0$,
with $d,C,\alpha$ positive bounded units and
$n,m,A,D\geq0$ arbitrary. Suppose an actual finite noncommuting
stationary zero has $u_c=k>0$ infinitesimal and $v_c=a>0$ a bounded
unit. Then the entire original parent is tensor SOS at every
independent projective measurement pair and every real or complex
projective state. No ratios among $n,A,D$ are restricted.
\end{theorem}

The bounded-$A$ result is a special case. The mechanism is to keep
all three possibly large event amplitudes as state coordinates.
The exact parent and its amplitudes are (TD1); in particular
$J=uW+uvX-Y-vZ$.

\paragraph{Source entry with three large weights.}
The initial coercivity estimate in the preceding proof uses only
$d,C,\alpha$, not $A$. Thus after $W_c=1$,
\[
 X_c,Y_c,Z_c=O(k),\qquad m=O(k^2).
\]
Now use $F_Z=0$ first. The $A$ and $n$ packets do not enter this
state equation. Every term other than $-8DvJ$ is $O(k)$, because
$v_c=a$ is a unit and the just displayed bounds hold. Therefore
$DJ_c=O(k)$. The $Y$ equation then gives $AY_c=O(k)$, and the $X$
equation gives $(d+2n)X_c=O(k^2)$. Write
\[
 N=d+2n,\quad T_A=d+2A,\quad L=d+2D,\qquad
 x=X_c/k^2,\quad y=Y_c/k,\quad j=J_c/k.
\]
Consequently
\[
 x,y,j,Nx,Ay,Dj\ \hbox{are bounded},\qquad N^{-1}\ \hbox{bounded}.
                                                               \tag{TA1}
\]
These estimates do not yet invert $A$ or $D$.

Use the exact stationary equations (TD3) and divide both by
$T_A L$. Now $t_2/T_A,t_6/T_A,t_3/L,t_4$ are positive bounded
units. In the first equation the coefficient multiplying $U^2$ is
$(A/T_A)$ times a positive bounded unit: its remaining bracket is
\[
 4t_4\left\{\theta\frac{t_3}{L}
            +CV\frac{1+2D/N}{L}\right\},
\]
whose lower bound follows from (TD4). The coefficient multiplying
$m$ is also a positive bounded unit:
\[
 4V\left\{\alpha\frac{t_8}{T_A L}
              +\frac{D}{L}\frac{t_6}{T_A}V\right\}.             \tag{TA2}
\]
Indeed, if $D/L$ is bounded below, the second summand is a unit.
Otherwise $D$ is infinitesimal and $L$ is a unit, while
$t_8/T_A=1+2(C+D)/T_A$ is a unit. All displayed factors have
bounded upper bounds in either case. Hence
\[
 \frac{m}{U^2}=\frac{A}{T_A}\cdot(\hbox{positive bounded unit}). \tag{TA3}
\]
The second stationary equation becomes
\[
 \frac{A}{T_A}\cdot(\hbox{positive unit})
       +O\left(\frac{A}{T_A}U\right)
 =\frac{D}{L}\cdot(\hbox{positive unit})
       +O\left(\frac{A}{T_A}U\right).                          \tag{TA4}
\]
Thus $A/T_A$ and $D/L$ are comparable. If they were infinitesimal,
$A,D,m$ would all be infinitesimal, contradicting exactly the
physicality test printed after (TD5). Therefore
\[
 A/T_A,\ D/L,\ m/k^4\quad\hbox{are positive bounded units}.       \tag{TA5}
\]
In particular $A,D$ have bounded reciprocals, but may be arbitrarily
large. Set $\gamma=4ma^2/k^4$, also a positive bounded unit.

\paragraph{The three-packet graph.}
On the local testing domain $\bar v=\bar a\ne0$, put
$b=1-j-y$ and use
\[
 X_g=xu^2W,\qquad Y_g=yuW,\qquad J_g=juW,\qquad
 Z_g=u\left(\frac b v+xu^2\right)W.                            \tag{TA6}
\]
This graph matches the whole actual contact: the definition of $J_c$
gives $Z_c/k=b/a+xk^2$. Its three hard normals are
\[
 E_X=\sqrt N(X-xu^2W),\quad
 E_Y=\sqrt A(Y-yuW),\quad E_J=\sqrt D(J-juW).                    \tag{TA7}
\]
The exact inverse is
\[
 X=xu^2W+E_X/\sqrt N,\quad Y=yuW+E_Y/\sqrt A,
\quad Z=\frac{uW+uvX-Y-juW-E_J/\sqrt D}{v}.                    \tag{TA8}
\]
Every denominator is either pure coefficient or the pure testing
unit $v$. All graph coefficients and all scaled normal coefficients
are bounded. For the three large packets this follows explicitly
from
\[
 n x^2,\ nx,\ n/N,\quad A y^2,\ \sqrt A\,y,\quad
 D j^2,\ \sqrt D\,j\quad\hbox{bounded},                       \tag{TA9}
\]
using (TA1),(TA5). Their quadratic normal contributions are exactly
$4(n/N)D_uD_v E_X^2$, $4D_uD_vE_Y^2$, and $4E_J^2$.
Every other original source coefficient is bounded, including $m$.

For clarity, the finite Taylor division used below works in the
bounded Laurent polynomial ring in $v$: for positive integer $r$,
\[
 v^{-r}-a^{-r}
 =-(v-a)\frac{\sum_{i=0}^{r-1}v^i a^{r-1-i}}{a^rv^r}.
\]
Thus a Laurent polynomial vanishing at $(U,v)=(K,a)$ belongs to
$(U-K,v-a)$ with bounded Laurent coquotients; repeating the same
identity gives the second-order division. There is no localization
at a difference of the two fields.

Set $K=k^2$, $R=u^2-(K/a)v$ and $T=u(v-a)$. The entire graph
scalar is again $4mv^2+u^2H_0(u^2,v)$, now with bounded Laurent
coefficients in $v$. Its three scalar contact jets give (TD10),
and hence an exact bounded Gram on $(R,T)$ with $R$ residue pivot
$\bar\gamma$. The full first normals before their pure source
scale divisions are
\[
\begin{split}
 (F_X+uF_Z)|_g&=-8mv+u^2L_X(u^2,v),\\
 (F_Y-F_Z/v)|_g&=uL_Y(u^2,v),\\
 (-F_Z/v)|_g&=uL_J(u^2,v).                                     \end{split} \tag{TA10}
\]
Divide the rows by $\sqrt N,\sqrt A,\sqrt D$, respectively.
The scaled coefficients are bounded by (TA9). At contact they
vanish. Laurent Taylor division gives exactly the root membership
and infinitesimal $R$ couplings proved in (TD11)--(TD12).
The $X$ row has infinitesimal coefficients on both $R,T$ because
$m/K=O(K)$; the other rows have infinitesimal coefficients on $R$.
The exact quadratic state remainder is bounded by (TA8)--(TA9),
so the entire original parent has a bounded five-root Gram on
$(E_X,E_Y,E_J,R,T)$.

\paragraph{Uniform strictness without a largest-event choice.}
The $X$ and $R$ residue pivots are $2(1+\bar a^2)$ and
$\bar\gamma$. The remaining hard state form is still (TD13).
At $u=0$, (TA8) gives
\[
 Y=E_Y/\sqrt A,\qquad
 Z=-\frac{E_Y}{v\sqrt A}-\frac{E_J}{v\sqrt D}.
\]
In these coordinates (TD13) is the sum of the two explicit pivots
\[
 4(1+v^2)E_Y^2+4E_J^2
\]
and the nonnegative bounded form
$(4C+2d)(1+v^2)Y^2+4\alpha(vY-Z)^2$.
Its residue is therefore positive definite for every combination of
bounded or unbounded $A,D$.

The boundary Schur function is the exact expression
\[
 f(v)=\frac{8(1+v^2)}{\det B(v)/(T_A L)}
 \left(\frac{Q_{10}}{T_A L}+
       \frac{Q_0}{T_A L}v^2+\frac{Q_{12}}{T_A L}v^4\right),     \tag{TA11}
\]
with $Q_{10},Q_{12}$ as in (TD14). For $v$ a bounded unit all its
coefficients and denominator are bounded, and
\[
 \frac{Q_{10}}{T_A L}=4\alpha\frac A{T_A}\frac{t_3}{L},\qquad
 \frac{Q_{12}}{T_A L}=4C\frac D L\frac{t_2-2m}{T_A}
\]
are positive units. The determinant has the positive term
$16AD(1+v^2)v^2$, so its quotient by $T_A L$ has a positive lower
bound. Its other terms are bounded after that division.
The exact middle coefficient is polynomial of degree at most one
in each of $A,D$, so $Q_0/(T_A L)$ is bounded too.

The contact normal equations and scalar jets give
$\bar f(\bar a)=\bar f'(\bar a)=0$ as before. Therefore the real
quartic bracket equals
$\overline{Q_{12}/(T_A L)}(v^2-\bar a^2)^2$.
Its strict second derivative gives the positive soft Schur pivot on
$T$. The whole five-root residue Gram is strict. This proves the
local all-state certificate; the Bob sign identity gives its negative
sign copy. Homogenization in $W$ includes $W=0$.

\paragraph{Outer tests when both $A$ and $D$ are large.}
For $u$ infinitesimal and $v$ a bounded unit with
$\bar v\ne\pm\bar a$, retain the entire $m,n$ packets and use the
four roots
\[
 X,\quad K_A=\sqrt A Y,\quad K_D=\sqrt D J,\quad uW.
\]
Substitute $Y=K_A/\sqrt A$ and
$Z=(uW+uvX-K_A/\sqrt A-K_D/\sqrt D)/v$.
The exact remainder is a bounded quadratic Gram. Its hard residue
has the positive $X$ pivot $2\bar d(1+\bar v^2)$ and the two
explicit pivots just displayed; its soft Schur complement is
$\bar f(\bar v)>0$ by (TA11).

At $u,v$ both infinitesimal, and at $u,1/v$ both infinitesimal,
one need not divide by $v$. Cap $C,A,D,\alpha$ at the positive
unit $\mu=\min(1,C/d,A/d,D/d,\alpha/d)$ after division by $d$,
and retain all excesses and the full $m,n$ packets. The bounded
selected parent has an exact Gram on $(X,Y,Z,uW)$ in the finite
chart and on the same four roots in the homogeneous reciprocal
chart. At $(u,v)=(0,0)$ its soft Schur complement is
$8Q_{10,\mu}/\det B_\mu(0)>0$; at $(u,1/v)=(0,0)$ it is
$8Q_{12,\mu}/\det B_{\mu,\infty}>0$. The hard blocks are positive
by (TD13) and its homogeneous limit. Hence both entire endpoint
neighborhoods have strict-residue Grams, including literal zeros
and infinities. These expressions use only the bounded selected
parent; no residue of an unbounded $A$ or $D$ is taken.

All other commuting tests use the same cap and the already printed
open-axis and corner forms in the preceding all-tests proof. Their
hypotheses require only the four lower bounds supplied by (TA5).
Noncommuting residue tests use the retained arbitrary-physical-source
theorem. Restore the positive pure factors and split complex states
into real and imaginary parts. This completes every independent
test for the stated source domain.

\paragraph{Coverage gained and remaining restriction.}
The theorem removes every relative-scale restriction among $n,A,D$
from the finite pitchfork. Its reflected form covers the corresponding
ordinary--infinite actual contacts. It still assumes $C/d$ and
$\alpha/d$ are bounded units and this actual-angle residue type.
Unbounded $C/d$ or $\alpha/d$, other actual-angle types, and the
exhaustive C2 join remain outside this theorem.


\section{\texorpdfstring{C2 pitchfork with two arbitrarily large retained events}{C2 pitchfork with two arbitrarily large retained events}}
\label{component:p0c7a15f800d9}

\label{p0c7a15f800d9:sec:c2-two-large}

\begin{theorem}[The $n,D$-unbounded C2 pitchfork]
\label{p0c7a15f800d9:thm:c2-two-large}
Let the physical canonical C2 source over a real closed coefficient
field have $\lambda=\beta=0$, with $d,C,\alpha$ positive bounded
units, $A\geq0$ bounded, and $n,m,D\geq0$ arbitrary. Suppose there is
an actual finite noncommuting stationary zero with $u_c=k>0$
infinitesimal and $v_c=a>0$ a bounded unit. Then the entire original
parent is tensor SOS at every independent projective measurement
pair and every real or complex projective state.
There is no bound on $n,D$, or on their ratio. In fact physicality and
contact force $A$ to be a positive unit, $D$ to be bounded below by
a positive real, and $m/k^4$ to be a positive bounded unit.
\end{theorem}

This extends the preceding pitchfork by removing the upper bound on
$D$. A bound on $m$ need not be assumed: it follows from physicality.
The original $D$ packet supplies a state normal, so an unbounded
weight is not put into a bounded Gram coefficient.

Write the unchanged parent as
\[
\begin{split}
F={}&dH+4nD_uD_vX^2+4mR_2^2+4CD_vr^2
       +4AD_uD_vY^2+4DJ^2+4\alpha D_up^2,\\
H={}&2\{J^2+D_up^2+D_vB^2-D_uD_vZ^2\},\\
R_2={}&vW-X+uvY-uZ,\quad r=uW-Y,\quad
J=uW+uvX-Y-vZ,\\
p={}&vY-Z,\quad B=X+uZ,\quad D_u=1+u^2,\quad D_v=1+v^2.
\end{split}                                                    \tag{TD1}
\]

\paragraph{Entry and bounded contact coordinates.}
At $u=0,v=a$, the state block on $(X,Y,Z)$ of
$dH+4CD_vr^2+4\alpha D_up^2$ has positive definite residue.
Its $W$ column is $O(k)$ and its $W^2$ entry is $O(k^2)$.
Completing this ordinary bounded hard block gives, for positive
real constants $b,B_0$, the coefficient-field inequality
\[
 dH+4CD_vr^2+4\alpha D_up^2
 \geq b(X^2+Y^2+Z^2)-B_0 k^2W^2
\]
at the actual angles. All omitted packets are nonnegative.
Thus an actual zero has $W\ne0$ and, on setting $W=1$,
$X,Y,Z=O(k)$. Since $R_2=a+O(k)$ is a unit, the same inequality
and its retained $m$ packet give $m=O(k^2)$.

The state equation $F_Y=0$ gives $DJ=O(k)$: every other term is
$O(k)$, using bounded $A,C,\alpha,d$, $m=O(k^2)$ and the bounds
just obtained. The equation $F_X=0$ then gives
$(d+2n)X=O(k^2)$, since the $J$ contribution is $O(kDJ)$ and the
$m$ contribution is $O(m)$, while $Z=O(k)$. Consequently, with
\[
 N=d+2n,\qquad x=X_c/k^2,\quad z=Z_c/k,\quad j=J_c/k,
\]
we have
\[
 Nx,\ z,\ j,\ Dj\quad\hbox{bounded},\qquad N^{-1}\ \hbox{bounded}.
                                                               \tag{TD2}
\]
The bound on $j$ follows directly from $X,Y,Z=O(k)$, even before a
lower bound for $D$ has been obtained. No inverse of $D$ has yet
been used.

For the scale entry, put $U=k^2,V=a^2$, $L=d+2D$,
$\theta=n/N$ and use
\[
\begin{gathered}
t_2=d+2(\alpha+m+A),\quad t_3=d+2(C+D),\quad
 t_4=d+2(\alpha+m),\\
t_6=d+2(\alpha+C+A),\quad t_8=d+2(C+A+D).
\end{gathered}
\]
The original Laurent-determinant stationary equations, divided only
by $N$, are
\[
\begin{split}
U^2\{4\theta A t_3t_4+4CA(1+2D/N)t_4V\}
 &=mV\{4\alpha t_8+4Dt_6V\},\\
4\alpha A(1+2m/N)t_3+4\theta A t_3t_4U
 &=\{4mDt_6/U+4CDt_2\}V^2.
\end{split}                       \tag{TD3}
\]
Here $t_2,t_4,t_6,t_3/L,t_8/L$ are positive bounded units.
After division by $L$, the braces in the first equation are
$A$ times a positive unit on the left and a positive unit on the
right. To check the lower bound uniformly when both $n,D$ diverge,
observe
\[
 \theta+\frac{1+2D/N}{L}
 =\frac{n+1}{N}+\frac{2n}{NL},                    \tag{TD4}
\]
which is a positive bounded unit; all its neighboring factors in
(TD3) have positive upper and lower real bounds. Thus
$m/U^2$ is $A$ times a positive unit. The second equation, again
divided by $L$, has the form
\[
 A\cdot(\hbox{positive unit})+O(AU)
    =(D/L)\cdot(\hbox{positive unit})+O(AU).       \tag{TD5}
\]
If $A$ were infinitesimal, then $D/L$ and hence $D$ would be
infinitesimal, as would $m$. This is incompatible with physicality:
on the test state $(W,X,Y,Z)=(1,0,u,zu)$ the $n$ packet vanishes
exactly, and at the residue source $m=A=D=0$ choose
\[
 z=\frac{v(2\bar\alpha+\bar d)}{2\bar\alpha}.
\]
The coefficient of $u^2$ in the full polynomial is then
$-\bar d v^2(2\bar\alpha+\bar d)/\bar\alpha<0$ for real $v\ne0$.
The bounded residue polynomial cannot be nonnegative for all real
small $u$. Therefore $A$ is a unit. Equations (TD3)--(TD5) now prove
that $D/L$ and $m/k^4$ are positive bounded units. In particular
$D>0$ and $D^{-1}$ is bounded. Set $\gamma=4ma^2/k^4$.

\paragraph{A graph that keeps both large packets bounded.}
For arbitrary independent testing angles use the exact graph
\[
 X_g=xu^2W,\qquad Z_g=zuW,\qquad J_g=juW,
\]
which fixes
\[
 Y_g=u(1-j-vz+xu^2v)W.                            \tag{TD6}
\]
It matches the entire actual contact, including its $Y$ coordinate.
The three exact state normals are
\[
 E_X=\sqrt N(X-xu^2W),\qquad
 E_J=\sqrt D(J-juW),\qquad E_Z=Z-zuW.             \tag{TD7}
\]
The inverse uses only pure coefficient inverses:
\[
 X=xu^2W+E_X/\sqrt N,\quad Z=zuW+E_Z,\quad
 Y=u(1-j)W+uvX-vZ-E_J/\sqrt D.                   \tag{TD8}
\]
The $n$ packet on the graph is
$4nx^2u^4D_uD_vW^2$, and the $D$ packet is
$4Dj^2u^2W^2$. Both have bounded coefficients by (TD2).
Their linear coefficients are bounded too, since $nx$ and
$\sqrt D\,j=(Dj)/\sqrt D$ are bounded. The exact quadratic
normal block is bounded: its only large packets become
$4(n/N)D_uD_vE_X^2$ and $4E_J^2$.

Set $W=1$, $K=k^2$, and
\[
 R=u^2-(K/a)v,\qquad T=u(v-a).
\]
The entire graph scalar has the exact shape
\[
 F|_g=4mv^2+u^2H_0(u^2,v),\qquad H_0\in B_c[u^2,v].             \tag{TD9}
\]
The three contact scalar jets give exactly the same monic division
as for the preceding pitchfork:
\[
 F|_g=\gamma R^2+u^2\{(u^2-K)^2H_{20}
                   +(u^2-K)(v-a)H_{11}+(v-a)^2H_{02}\},         \tag{TD10}
\]
with bounded polynomial coquotients. After
$u^2-K=R+(K/a)(v-a)$ this is a bounded Gram on $(R,T)$.
Its $R$ residue pivot is $\bar\gamma>0$.

The full first normals, before the indicated scale divisions, have
exact shapes
\[
 (F_X+uvF_Y)|_g=-8mv+u^2L_X(u^2,v),\quad
 (-F_Y)|_g=uL_J(u^2,v),\quad
 (F_Z-vF_Y)|_g=uL_Z(u^2,v).                       \tag{TD11}
\]
The coefficient of $E_X$ is the first expression divided by
$\sqrt N$; that of $E_J$ is the second divided by $\sqrt D$.
These scaled polynomials have bounded coefficients, as shown above.
Contact makes $L_X(K,a)=8ma/K$, and makes the scaled $J,Z$ functions
vanish at $(K,a)$. Monic division by $u^2-K,v-a$ therefore puts
every first normal in $(R,T)$ with bounded coquotients.
Every coefficient coupling $E_X$ to $R,T$ is infinitesimal; the
$R$ coefficients in the $E_J,E_Z$ rows are infinitesimal.
For example, the $X$ row is exactly
\[
\frac1{\sqrt N}\left[
 (8ma/K+u^2L_{X1})R
 +u\{(K/a)L_{X1}+L_{X2}\}T\right],              \tag{TD12}
\]
and the other two rows have form
$uL_{i1}R+\{(K/a)L_{i1}+L_{i2}\}T$ after their scale division.
State degree is two, so there are no omitted higher-normal terms.
Equations (TD7)--(TD12) give an exact bounded full-parent Gram on
$(E_X,E_J,E_Z,R,T)$.

\paragraph{The residue Gram is strict.}
The $X$ pivot has residue $2(1+\bar a^2)$ and the $R$ pivot has
residue $\gamma$. For the remaining hard pair the unscaled state
form is
\[
 B(v)[Y,Z]=\{4(C+A)+2d\}(1+v^2)Y^2
           +4D(Y+vZ)^2+4\alpha(vY-Z)^2.           \tag{TD13}
\]
In the scaled coordinates $Y=-vE_Z-E_J/\sqrt D$, $Z=E_Z$ its
coefficients are bounded and its residue is positive definite.
If $D$ is bounded, this follows from the invertible residue change
and the positive form (TD13). If $D$ is unbounded, its residue is
\[
 4E_J^2+\bigl(\{4(\bar C+\bar A)+2\bar d\}(1+v^2)v^2
                  +4\bar\alpha(1+v^2)^2\bigr)E_Z^2\succ0.
\]
The exact minimized boundary coefficient of $u^2$, before residue,
is
\[
 f(v)=\frac{8(1+v^2)}{\det B(v)/L}
       \left(\frac{Q_{10}}L+\frac{Q_0}L v^2+\frac{Q_{12}}L v^4\right),
                                                               \tag{TD14}
\]
where
\[
 Q_{10}=4\alpha A[d+2(C+D)],\qquad
 Q_{12}=4CD[d+2(\alpha+A)].
\]
The middle coefficient $Q_0$ is defined by the exact Schur numerator;
its explicit value is unnecessary. All quantities in (TD14) are
bounded, $\det B/L$ is a positive unit for bounded real $v$, and
$Q_{10}/L,Q_{12}/L$ are positive units. One may verify boundedness
of $Q_0/L$ directly from the determinant expansion, which is affine
in $D$. This is only a residue calculation of a real hard Schur
complement, not inversion of a mixed tensor element.

The full contact normal and scalar equations imply that the residue
of (TD6) at $v=\bar a$ minimizes the hard form, and that
$\bar f(\bar a)=\bar f'(\bar a)=0$. Thus the real quartic bracket
in (TD14) equals
\[
 \overline{Q_{12}/L}\,(v^2-\bar a^2)^2.
\]
Its second derivative at $\bar a$ is strictly positive. This is
exactly the Schur complement on $T$ of the remaining residue block
$(E_J,E_Z,T)$. Hence the whole five-root residue Gram is positive
definite. The bounded strict-residue Gram lemma proves the local
certificate at $\bar u=0,\bar v=\bar a$. The Bob sign copy covers
$\bar v=-\bar a$.

\paragraph{All other independent tests.}
For $u$ infinitesimal and $v$ bounded with
$\bar v\ne\pm\bar a$, retain the original $m,n$ packets. Set
$K_J=\sqrt D J$ and substitute
\[
 Y=uW+uvX-vZ-K_J/\sqrt D.
\]
The exact remainder is a bounded quadratic Gram on
$(X,Z,K_J,uW)$. Its hard residue is the block just considered with
$X$ pivot $2\bar d(1+\bar v^2)$; its soft Schur complement is
$\bar f(\bar v)>0$. This includes $v=0$ and all infinitesimal $v$.
For $u,q=1/v$ infinitesimal use the homogeneous parent $q^2F(u,1/q)$.
Here $J_h=quW+uX-qY-Z$. Retain $m,n$ and substitute
$Z=quW+uX-qY-K_J/\sqrt D$ with $K_J=\sqrt D J_h$.
The remainder is a bounded Gram on $(X,Y,K_J,uW)$. Its soft residue
is $\lim_{q\to0}q^2\bar f(1/q)>0$, since
$\overline{Q_{12}/L}>0$ and
$\overline{\det B_\infty/L}>0$.
These are polynomial homogeneous expressions, so literal projective
endpoints are included.

At every other commuting test the capping proof from the preceding
all-tests corollary applies verbatim: choose
$\mu=\min(1,C/d,A/d,D/d,\alpha/d)$, a positive unit, retain every
excess and the whole $m,n$ packets, and use its displayed open-axis
state forms and two corner quadratic forms. That proof used lower
bounds, not upper bounds, on these four source weights. Noncommuting
residue tests use the retained arbitrary-physical-source theorem.
Finally homogenize (TD6)--(TD12) in $W$, replacing $R,T$ by $RW,TW$;
restore pure testing projective factors and split complex states into
real and imaginary parts. This proves the theorem for all states,
including $W=0$, without comparing $k,n,D$ with independent testing
infinitesimals.

\paragraph{Exact scope.}
This is an all-tests theorem for the stated actual-contact domain.
It does not bound or remove the residual regimes where $C/d$ or
$\alpha/d$ is unbounded, $A/d$ is unbounded, or the actual contact
has a different angle-residue type. The C2 reflection
$(m,C,D,\alpha;u_c,v_c)\mapsto(D,\alpha,m,C;1/v_c,1/u_c)$ supplies
its corresponding reflected domain. No exhaustive C2 join follows
from these two domains alone.


\section{\texorpdfstring{A bounded-shape one-weak C2 family, at all tests}{A bounded-shape one-weak C2 family, at all tests}}
\label{component:pec59340d644f}

\label{pec59340d644f:sec:c2-general-single-angle}

This module gives a source-domain extension of the fixed rational
family.  Its scalar uniqueness proof is an exact six-term Laurent
identity and a three-cycle arithmetic--geometric mean inequality.
The full-normal and outer-chart arguments include literal $n=0$.

Put $x>0$ infinitesimal in the coefficient field and $\rho=\sqrt{x}$.
Choose positive bounded coefficient units $k,b,r,h$ such that
\[
 1-r,\quad b-h,\quad g=b-1+r+xh,\quad b-k^2r
\]
are positive bounded units.  Set
\[
 c=\rho,\quad a=k/\rho,\quad t=x(b-h)/k^2,
 \quad j=1+xr-t,\quad a_2=k^2/x,
 \quad A_c=1+xr,\quad S_c=A_c+a_2xt.
\]
Thus $h=b-a_2t$, $g=(1+x)b-(1-r+a_2xt)$, and $j=1+xr-t$.
Use the exact original rational source
\begin{align*}
 \Delta={}&(1+x)xb^2+2x(1+x)br+(1+x)(1+a_2x)b
       +x^2(1-a_2)r^2+2x(1+a_2x)r+1+a_2x,\\
 N={}&(1+x)(b+r)(b+a_2(1-t))-a_2(1+a_2x)(1-r)t,\\
 n={}&\frac{(b+r)A_cS_c-r\Delta}{2r\Delta},\\
 m={}&\frac{(1+a_2x)(1-r)A_cS_c}{2g\Delta},\\
 C={}&\frac{A_cN}{2a_2t\Delta},\\
 A={}&\frac{a_2A_c(b+r)S_c}{2h\Delta},\\
 D={}&\frac{(1+x)(b+r)(b-a_2xr)S_c}{2a_2j\Delta},\\
 \alpha={}&\frac{x^2(b+r)N}{2(1-xh)\Delta}.
\end{align*}
Assume $n\geq0$ and physicality of this canonical source.  The latter
is the same inherited source-physicality hypothesis used for the
general noncommuting-residue exit; no conclusion below is predicated
on a new source enumeration.  Literal $n=0$ is allowed.

The parent, with $d=1$ and $\lambda=\beta=0$, is
\begin{align*}
 F={}&4nD_uD_vX^2+4mM^2+4CD_v(uW-Y)^2+4AD_uD_vY^2
                +4DJ^2+4\alpha D_u(vY-Z)^2+H,\\
 M={}&vW-X+uvY-uZ,\qquad J=uW+uvX-Y-vZ,\\
 H/2={}&J^2+D_u(vY-Z)^2+D_v(X+uZ)^2-D_uD_vZ^2,
 \qquad D_u=1+u^2,\quad D_v=1+v^2.
\end{align*}
Its exact stationary contact is
\[
 (u,v;W,X,Y,Z)=(\rho,k;\rho b/k,-\rho r,xh/k,1).
\]
Writing $\kappa=b/k$ and using a subscript zero for residues gives
\begin{align*}
 \Delta_0&=(1+k_0^2)(b_0+1),\qquad
 n_0=\frac{b_0+r_0-r_0\Delta_0}{2r_0\Delta_0}\geq0,\\
 m_0&=\frac{1-r_0}{2(b_0-1+r_0)(b_0+1)}>0,\\
 C_*:=(xC)_0&=\frac{k_0^2(b_0+r_0)}{2(b_0-h_0)\Delta_0},\qquad
 A_*:=(xA)_0=\frac{k_0^2(b_0+r_0)}{2h_0\Delta_0},\\
 D_*:=(D/x)_0&=\frac{(b_0+r_0)(b_0-k_0^2r_0)}{2k_0^2\Delta_0},\qquad
 E_*:=(\alpha/x)_0=\frac{k_0^2(b_0+r_0)^2}{2\Delta_0}.
\end{align*}
All starred quantities are positive.  Thus $n$ is bounded, $m$ is a
positive unit, $C,A$ have scale $x^{-1}$, and $D,\alpha$ have scale $x$.

\begin{theorem}[General bounded-shape single-angle source]
\label{pec59340d644f:thm:c2-general-single-angle}
Every source in the displayed domain has tensor SOS for the full
original parent at all independent real or complex projective testing
states and all projective testing angles.  No comparison between a
coefficient infinitesimal and a testing infinitesimal is needed.
The conclusion includes $n=0$.  Applying the exact C2 reflection gives
the corresponding original orientation with weak $C$ and strong
$\alpha$; this is a sufficient source domain, not a global one-weak
C2 classification.
\end{theorem}

\paragraph{The finite $u=0$ axis and its full normals.}
Set $a_0=4m_0$, $b_0'=a_0+2$, and
\[
 H_0(v)=(4n_0+2)(1+v^2)+a_0,\quad \eta=h/b,
 \quad \gamma=-r-\frac{a_0b-b_0'}{H_0(k)}.
\]
Here $\gamma$ is infinitesimal and $\eta_0=C_* /(C_*+A_*)$.
Use the exact contact graph
\[
 X_g=\frac{a_0vW-b_0'uZ}{H_0(v)}+\gamma uZ,
 \qquad Y_g=\eta uW.
\]
The denominator $H_0(v)$ is a pure testing unit and $H_0(k)$ is a
pure coefficient unit.  Let $N_X=X-X_g$, $N_Y=(Y-Y_g)/\rho$.
Writing $f=F|_g=f_{ww}W^2+f_{wz}WZ+f_{zz}Z^2$, exact monic divisions give
\[
 f=aP^2+b_s B^2+c_s V^2+eE^2+g_sBV,
 \qquad(P,B,V,E)=(uW/\rho,uZ,vW,\rho Z),
\]
where
\[
 \begin{gathered}
 a=x(f_{ww}-f_{ww}|_{u=0})/u^2,\qquad
 b_s=(f_{zz}-f_{zz}|_{u=0})/u^2,\\
 c_s=f_{ww}|_{u=0}/v^2,\qquad e=f_{zz}|_{u=0}/x,\qquad
 g_s=f_{wz}/(uv).
 \end{gathered}
\]
These displayed quotients are polynomial divisions, not inverses of
testing infinitesimals.  Parity gives the divisions by $u^2,v^2,uv$;
$f_{zz}|_{u=0}=4Dv^2+4\alpha$ gives the $x$ factor.  Only $C,A$ can
produce coefficient poles, and they occur in $f_{ww}$ with a factor
$u^2$; its prefactor $x$ removes them.

Put $Q=4C_*A_* /(C_*+A_*)>0$.  The bounded coefficient residues are
\begin{align*}
 q(v)=a_0^{\rm sc}&=Q D_v,&
 b(v)&=b_0'\frac{H_0-a_0-2}{H_0},\\
 \bar c(v)&=a_0\frac{H_0-a_0}{H_0},&
 g(v)&=-2b_0'\frac{H_0-a_0}{H_0},&
 e(v)&=4D_*v^2+4E_*.
\end{align*}
Thus $c(v):=v^2\bar c(v)$ and $d(v):=-vg(v)$ give the weighted
leading scalar
\[
 q(v)u^2W^2/x+b(v)u^2Z^2+c(v)W^2-d(v)uWZ+e(v)xZ^2.
\]
The full expansion is
\[
 F=f+N_XL_X+N_YL_Y+H_{XX}N_X^2+2H_{XY}N_XN_Y+H_{YY}N_Y^2.
\]
Its hard residue is $\operatorname{diag}(H_0(v),4D_v(C_*+A_*))$.
Direct differentiation gives
\[
 L_X=V\ell_w+B\ell_z,\qquad
 L_Y=P(B_2+xe_w)+B(\rho u e_z)+E e_*,
\]
where
\[
 B_2=8D_v\{(xA)\eta D_u-(xC)(1-\eta)\},\qquad
 e_*=8v(D-\alpha).
\]
All $\ell_w,\ell_z,e_w,e_z$ are bounded; the residues of
$\ell_w,\ell_z,B_2,e_*$ vanish identically in finite $v$.
For the second identity, after subtracting $uWB_2/x$ from $F_Y|_g$,
the remaining $W$ coefficient is $ue_w$ and the remaining $Z$
coefficient is $u^2e_z+e_*$.  The constant signed-core derivatives
cancel.  This proves full normal restoration wherever the scalar
Gram on these four roots is strict.

\paragraph{Global scalar strictness by an exact cycle identity.}
Let $\nu=2n_0+1\geq1$, $\mu=2m_0+1>1$, $z=v^2$,
$\mathcal H=\nu(1+z)+\mu-1$, $e_2=4D_*$ and $e_0=4E_*$.  Then
\[
 b=\frac{2\mu(\nu(1+z)-1)}{\mathcal H},\quad
 c=\frac{2(\mu-1)z\nu(1+z)}{\mathcal H},\quad
 d=\frac{4\mu v\nu(1+z)}{\mathcal H},\quad
 e=e_2z+e_0.
\]
For $s,z>0$ define
$\mathcal D(s,z)=(q s+c)(b s+e)-d^2s/4$.
The following exact identity is the key:
\begin{align*}
 \frac{\mathcal H\mathcal D(s,z)}{(1+z)sz}
 ={}&2Q\mu(\nu-1)s/z+2Q\mu\nu s
 +2(\mu-1)\nu e_2z/s+2(\mu-1)\nu e_0/s\\
 &+Qe_0(\nu+\mu-1)/z+Qe_2\nu z\\
 &+Q\{e_2(\nu+\mu-1)+e_0\nu\}-4\mu\nu.
\end{align*}
Every nonconstant coefficient is positive, except the first may be
zero when $n_0=0$.  Exact contact and stationarity, followed by taking
residues, give a stationary zero at $(s,z)=(1,k_0^2)$.  Here is an
algebraic proof that it is the unique zero on the whole positive
quadrant.  Rescale $z=k_0^2t$ and write the Laurent expression as
\[
 \mathcal L=a_\ell s/t+b_\ell s+c_\ell t/s+d_\ell/s
                      +e_\ell/t+f_\ell t-\Sigma,
 \qquad \Sigma=a_\ell+b_\ell+c_\ell+d_\ell+e_\ell+f_\ell.
\]
The constant term has this form by its zero value at $(1,1)$.
Stationarity gives
$a_\ell-c_\ell=d_\ell-b_\ell=f_\ell-e_\ell=:\tau$.
Put $p_\ell=\min(a_\ell,c_\ell)\geq0$,
$q_\ell=\min(b_\ell,d_\ell)>0$,
$r_\ell=\min(e_\ell,f_\ell)>0$.  Exactly,
\begin{align*}
 \mathcal L={}&p_\ell(s/t+t/s-2)+q_\ell(s+1/s-2)
                              +r_\ell(t+1/t-2)\\
 &+\begin{cases}
 \tau(s/t+1/s+t-3),&\tau\geq0,\\
 (-\tau)(t/s+s+1/t-3),&\tau\leq0.
 \end{cases}
\end{align*}
Each bracket is nonnegative by the two- or three-term
arithmetic--geometric mean inequality.  The two strictly positive
pair terms force $s=t=1$ at equality.  They also give a positive
definite Hessian there: the other brackets are nonnegative and
vanish at that point.  Thus both remote uniqueness and local
nondegeneracy are proved by the same finite algebraic identity, over
any real closed residue field.  No inference from a merely local
Hessian is being used.

Consequently, for finite nonzero $v$ with $v^2\ne k_0^2$,
\[
 |d(v)|<2\sqrt{b(v)c(v)}+2\sqrt{q(v)e(v)}.
\]
Indeed the minimum of $\mathcal D(s,z)/s$ over $s>0$ is
$(\sqrt{bc}+\sqrt{qe})^2-d^2/4$ and is attained, so positivity is
strict.  Split the exact cross term using $BV=vPE$ and the pure
testing function
\[
 \theta=\frac{|v|\sqrt{b(v)\bar c(v)}}
 {|v|\sqrt{b(v)\bar c(v)}+\sqrt{q(v)e(v)}}.
\]
The resulting residue blocks on $(B,V)$ and $(P,E)$ are strictly
positive.  The denominator is a pure testing unit.  At $v_0=0$ the
same conclusion holds when $n_0>0$, with zero cross residues.
The case $n_0=0,v_0=0$ is treated separately below.

\paragraph{The matched finite-angle residue.}
For $v_0=k_0$, normalize $Z=1$ when $Z\ne0$.  If $W$ is
infinitesimal, put
\[
 R=(uW-\kappa x)/\rho,\quad S=u-\kappa^{-1}W,
 \quad U=\rho(v-k),\quad T=u(v-k).
\]
The exact conductors are
\[
 u^2-x=uS+\kappa^{-1}\rho R,\qquad
 W^2-\kappa^2x=-\kappa WS+\kappa\rho R,
 \qquad uW-\kappa x=\rho R.
\]
Write $\xi=u^2,\omega=W^2,p=uW$ and
\[
 f=p^2Q_b/x+B_b,\qquad
 Q_b=4D_v\{(xC)(1-\eta)^2+(xA)\eta^2D_u\}.
\]
Here $B_b$ is bounded and polynomial in these three invariants with
only the separate unit denominators already specified.  Put
$G=\kappa^2xQ_b+B_b$ and take successive monic differences
\[
 G=G_0(v)+(\xi-x)G_\xi+(\omega-\kappa^2x)G_\omega
                    +(p-\kappa x)G_p.
\]
Exactly,
\[
 f=R^2Q_b+\kappa^2S^2G_\omega+\rho R\mathcal A+uS\mathcal B+G_0,
\]
where $\mathcal A=2\kappa Q_b+\kappa^{-1}G_\xi+
\kappa G_\omega+G_p$ and $\mathcal B=G_\xi-\kappa^2G_\omega$.
The contact Jacobian of $(R,S)$ with respect to $(u,W)$ has
determinant $-2$.  Exact stationarity therefore gives
$\mathcal A_c=\mathcal B_c=0$.  Cancelling $\rho$ here proves
vanishing and does not introduce a Gram coefficient inverse.
Monic contact division of these two functions and the conductors
above give a bounded Gram on $(R,S,U,T)$.
Also $G_0=xg_0(v)$ with bounded $g_0$, and exact contact gives a
double factor $(v-k)^2$, producing its bounded $U^2$ coefficient.

The full normals have zero-residue lifts on these roots.  For $L_X$,
write $L_X=uA_1+WB_1=-\kappa B_1S+u(A_1+\kappa B_1)$.
Both coefficient functions have identically zero leading part for
every $v$; exact normal contact makes $A_1+\kappa B_1$ vanish at
contact.  Monic division gives coefficients with factors $u\rho$,
$u^2$, or $uW$, and a $(v-k)$ quotient of an identically zero
leading function.  All have zero residue.  For $L_Y$, write
$L_Y=RB_2+\rho E_2$.  Its bounded even $E_2$ vanishes exactly at
contact, and both $B_2,E_2$ have identically zero leading part.
Monic division gives factors $x,\rho u,\rho W$, or the same
zero-leading-function quotient.  Thus the entire normal restoration
is bounded and has zero residue.

The soft residue has the form
\[
 q(k_0)R^2+\kappa_0^2c(k_0)S^2+A_1'RU+B_1'ST+E_1'U^2.
\]
It satisfies
\[
 E_1'>\frac{(A_1')^2}{4q(k_0)}+
          \frac{(B_1')^2}{4\kappa_0^2c(k_0)}.
\]
To prove this strict inequality, the leading scalar at
$u=\rho s,W=\rho w$ is
$\Phi=q s^2w^2+b s^2+c w^2-dsw+e$.
Complete its $w$ square:
\[
 \Phi=(qs^2+c)\left(w-\frac{ds}{2(qs^2+c)}\right)^2
           +\frac{\mathcal D(s^2,v^2)}{qs^2+c}.
\]
The global Laurent calculation gives a positive definite Hessian of
the second term in $(s,v)$ at $(1,k_0)$; both coordinate derivatives
are nonzero there.  The first square supplies the missing $w$
direction.  Hence the full three-variable Hessian is positive
definite.  On first variations $T=U$; the displayed inequality is
exactly its Schur complement after the invertible $(R,S)$ change.
Choose a real $\delta$ strictly between the second summand and
$E_1'-(A_1')^2/(4q(k_0))$.  Add $\delta$ times the exact zero
\[
 T^2-U^2-ST(v-k)-\kappa^{-1}RU(v-k)=0.
\]
The two blocks $(R,U)$ and $(S,T)$ are now strictly positive, and so
is the full six-root residue Gram.

If the ratio $W/Z$ is not infinitesimal, instead normalize $W=1$;
then $Z$ is bounded.  The earlier scalar is $aP^2$ plus a term of
positive residue $c(k_0)$; its normal lifts have zero residue on
$(P,1)$.  This is strict.  The case $Z=0$ is a literal sum of the
event squares and the three positive core squares.  The simultaneous
sign change $(u,v,X,Y)\mapsto(-u,-v,-X,-Y)$ covers $v_0=-k_0$.

\paragraph{The double-zero corner when $n_0=0$.}
Use instead the pure coefficient quantities
\[
 H_c=4n+4m+2,\quad \xi_c=4m/H_c,\quad
 \zeta=(4m+2)/H_c,\quad \eta_c=C/(C+A),
\]
and graph $X_g=\xi_cvW-\zeta uZ$, $Y_g=\eta_cuW$.
Put $\epsilon=4\alpha+4n(4m+2)/H_c>0$.
Then $\epsilon$ is infinitesimal, while $n/\epsilon$ and
$x/\epsilon$ are bounded.  Use the hard roots as before and scalar
roots
\[
 P=uW/\rho,\quad V=vW,\quad E=\rho Z,\quad
 B_\epsilon=\sqrt\epsilon\,uZ,\quad T=uvZ.
\]
Direct packet expansion gives exactly
\[
 f_{zz}=4\alpha+\epsilon u^2+4n\zeta^2u^4+4Dv^2
                         +u^2v^2 B_b,
\]
where $B_b$ is bounded with residue $2$.  Its first four terms
give coefficients $4\alpha/x$, $1+(4n\zeta^2/\epsilon)u^2$,
and $(4D/x)v^2$ on $E^2,B_\epsilon^2,E^2$ respectively.  The
$W^2$ contribution has bounded coefficients on $P^2,V^2$ with
positive residues $Q$ and $4m_0/(2m_0+1)$.  The $WZ$ contribution
is $uvWZ$ times a bounded function; use $uvWZ=vPE$ to give it zero
cross residue.  Thus the five scalar roots have strict diagonal
residue $(Q,4m_0/(2m_0+1),4E_*,1,2)$.

For completeness, $L_X$ has a zero-residue bounded coefficient on
$V$.  At $v=0$ its $Z$ coefficient is exactly $-8n\zeta u^3$;
on $B_\epsilon$ this is $-8n\zeta u^2/\sqrt\epsilon$, which is
infinitesimal because $n/\epsilon$ is bounded.  The remaining
$Z$ coefficient has a factor $uv^2$, giving a zero-residue multiple
of $T$.  For $L_Y$, the large-packet coefficient is now exactly
$B_2=8D_v(xA)\eta_cu^2$; the remaining terms lift as
\[
 L_Y=P(B_2+xe_w)+T(\rho u e_z)+E\,8v(D-\alpha),
\]
with bounded $e_w,e_z$.  All coefficients have zero residue.  The
hard residue is $\operatorname{diag}(4m_0+2,4(C_*+A_*))$.
This includes literal $n=0$ without deleting a coordinate.

\paragraph{The remaining commuting charts.}
The following exact core identities retain all unused event packets.
Every retained coefficient is nonnegative; in particular the $n$
packet can be retained even when $n=0$.  They prove strict bounded
Grams without dividing by a testing infinitesimal.

If $u$ is a finite unit and $v\to0$, take
\[
 R=(uW-Y)/\rho,\quad L=Y/\rho,\quad B=X+uZ,\quad V=vZ.
\]
Then $M=-B+(\rho v/u)(R+D_uL)$, and exactly
\begin{align*}
 H/2={}&(1+D_uv^2)B^2+u^2D_uV^2-2uvD_uBV\\
 &+2\rho R(uvB-D_uV)+xR^2-2D_u\rho LV+D_uxv^2L^2.
\end{align*}
Retain $n,D,\alpha$.  The residue on $(B,V,R,L)$ is diagonal
with entries $(4m_0+2,2u_0^2D_{u,0},4C_*,4A_*D_{u,0})$.

If $u\to\infty$ and $v$ is finite, put $p=1/u$ and use $p^2F$.
Take $R=(W-pY)/\rho$, $L=Y/\rho$, $B=Z+pX$, $V=vX$.
The homogenized $m$ amplitude is
$-B+\rho pvR+\rho vD_pL$, and the exact core is
\begin{align*}
 \widehat H/2={}&B^2+D_pV^2+2\rho D_pRV-2\rho pvRB+xR^2\\
 &+2D_p\rho pLV-2D_p\rho vLB+D_pxv^2L^2.
\end{align*}
Retain $n,D,\alpha$.  The residue on $(B,V,R,L)$ is diagonal
with entries $(4m_0+2,2,4C_*D_{v,0},4A_*D_{v,0})$.

If $v\to\infty$ and $u$ is bounded, put $q=1/v$ and use $q^2F$.
Take $R=(uW-Y)/\rho$, $L=Y/\rho$, $T=uqZ$, $N_1=\rho Z$.
Exactly,
\begin{align*}
 \widehat H/2={}&(D_u+q^2)X^2+D_uxL^2+T^2+2qXT
       +q^2xR^2+2uq\rho XR-2qRN_1-2D_uqLN_1,\\
 4D\widehat J^2={}&4(D/x)(-N_1+\rho uX+xqR)^2.
\end{align*}
Retain $n,m,\alpha$.  The residue on $(X,R,L,T,N_1)$ is diagonal
with entries
\[ (2D_{u,0},4C_*,4A_*D_{u,0},2,4D_*). \]

Finally if both angles are infinite, use $p=1/u,q=1/v$ and
$p^2q^2F$.  Set $R=(W-pY)/\rho$, $L=Y/\rho$, $T=qZ$.  Then
\[
 \widehat H/2=(1+p^2D_q)X^2+xq^2R^2+xD_pL^2+T^2
       +2\rho qXR+2pqXT-2\rho pRT-2\rho D_pLT.
\]
Retain $n,m,D,\alpha$.  The residue on $(X,R,L,T)$ is
diagonal with entries $(2,4C_*,4A_*,2)$.

\paragraph{Exhaustion and independence.}
These charts cover every commuting angle residue: the whole finite
$u=0$ axis, the finite-unit $v=0$ axis, either infinite axis, and
their intersections.  When both angle residues are noncommuting,
apply the inherited physical noncommuting-residue theorem at arbitrary
coefficient scales.  Complex states require no additional chart:
if $Z\ne0$, multiply by a phase and normalize $Z=1$; the real parent
is the sum of its real-state value and its imaginary-state value,
the latter having $Z=0$ and hence being a literal square sum.
If $Z=0$ this is already true of both parts.  Homogeneity covers all
projective state charts.  Roots may be unbounded throughout.
All inverses or radicals used in coefficients are pure coefficient
or pure testing operations on units, except $\sqrt\epsilon$ in
roots and its controlled coefficient ratio $n/\sqrt\epsilon$.
No mixed inverse and no comparison of independent infinitesimals
occurs.  Applying the strict bounded-tensor Gram lemma completes
the proof.

